% concatenated from EJS2495.tex
\documentclass[ejsv2]{imsart}

%% Packages
\RequirePackage{amsthm,amsmath}%,amssymb}
%\RequirePackage[numbers]{natbib}
\RequirePackage[authoryear]{natbib}
\RequirePackage[colorlinks,linkcolor=blue,citecolor=blue,urlcolor=blue]{hyperref}
\RequirePackage{graphicx}
\RequirePackage{bm,algorithm,float,algpseudocode,accents}
\RequirePackage[utf8,applemac]{inputenc}
%\usepackage[affil-it]{authblk}
%\usepackage[title]{appendix}
%\usepackage{titling}
\newcommand{\ubar}[1]{\underaccent{\bar}{#1}}

%\arxiv{2010.00000}
\startlocaldefs
%%%%%%%%%%%%%%%%%%%%%%%%%%%%%%%%%%%%%%%%%%%%%%
%%                                          %%
%% Uncomment next line to change            %%
%% the type of equation numbering           %%
%%                                          %%
%%%%%%%%%%%%%%%%%%%%%%%%%%%%%%%%%%%%%%%%%%%%%%
%\numberwithin{equation}{section}
%%%%%%%%%%%%%%%%%%%%%%%%%%%%%%%%%%%%%%%%%%%%%%
%%                                          %%
%% For Axiom, Claim, Corollary, Hypothesis, %%
%% Lemma, Theorem, Proposition              %%
%% use \theoremstyle{plain}                 %%
%%                                          %%
%%%%%%%%%%%%%%%%%%%%%%%%%%%%%%%%%%%%%%%%%%%%%%
\theoremstyle{plain}
\newtheorem{axiom}{Axiom}
\newtheorem{claim}[axiom]{Claim}
\newtheorem{theorem}{Theorem}[section]
\newtheorem{lemma}[theorem]{Lemma}
\newtheorem{proposition}{Proposition}[section]
\newtheorem{corollary}{Corollary}[section]
%%%%%%%%%%%%%%%%%%%%%%%%%%%%%%%%%%%%%%%%%%%%%%
%%                                          %%
%% For Assumption, Definition, Example,     %%
%% Notation, Property, Remark, Fact         %%
%% use \theoremstyle{definition}            %%
%%                                          %%
%%%%%%%%%%%%%%%%%%%%%%%%%%%%%%%%%%%%%%%%%%%%%%
\theoremstyle{definition}
\newtheorem{definition}[theorem]{Definition}
\newtheorem{example}{Example}[section]
\newtheorem*{fact}{Fact}
%%%%%%%%%%%%%%%%%%%%%%%%%%%%%%%%%%%%%%%%%%%%%%
%%                                          %%
%% For Case use \theoremstyle{remark}       %%
%%                                          %%
%%%%%%%%%%%%%%%%%%%%%%%%%%%%%%%%%%%%%%%%%%%%%%
\theoremstyle{remark}
\newtheorem{case}{Case}
\newtheorem{remark}{Remark}[section]
%%%%%%%%%%%%%%%%%%%%%%%%%%%%%%%%%%%%%%%%%%%%%%
%% Please put your definitions here:        %%
%%%%%%%%%%%%%%%%%%%%%%%%%%%%%%%%%%%%%%%%%%%%%%
\endlocaldefs

\begin{document}
\begin{frontmatter}
\title{Negative Moment Bounds for Sample Autocovariance Matrices of Stationary Processes Driven by Conditional Heteroscedastic Errors and Their Applications}
%\title{A sample article title with some additional note\thanksref{t1}}
\runtitle{Negative Moment Bounds}
%\thankstext{T1}{A sample additional note to the title.}

\begin{aug}
%%%%%%%%%%%%%%%%%%%%%%%%%%%%%%%%%%%%%%%%%%%%%%%
%% Only one address is permitted per author. %%
%% Only division, organization and e-mail is %%
%% included in the address.                  %%
%% Additional information can be included in %%
%% the Acknowledgments section if necessary. %%
%% ORCID can be inserted by command:         %%
%% \orcid{0000-0000-0000-0000}               %%
%%%%%%%%%%%%%%%%%%%%%%%%%%%%%%%%%%%%%%%%%%%%%%%
\author[A]{\fnms{Hsueh-Han}~\snm{Huang}\ead[label=e1]{link4917@stat.sinica.edu.tw}},
\author[B]{\fnms{Ching-Kang}~\snm{Ing}\ead[label=e2]{cking@stat.nthu.edu.tw}}
\and
\author[C]{\fnms{Shu-Hui}~\snm{Yu}\ead[label=e3]{
shuhui@nuk.edu.tw}}
%%%%%%%%%%%%%%%%%%%%%%%%%%%%%%%%%%%%%%%%%%%%%%
%% Addresses                                %%
%%%%%%%%%%%%%%%%%%%%%%%%%%%%%%%%%%%%%%%%%%%%%%
\address[A]{Institute of Statistical Science,
Academia Sinica\printead[presep={,\ }]{e1}}

\address[B]{Institute of Statistics and Data Science,
National Tsing Hua University\printead[presep={,\ }]{e2}}

\address[C]{Institute of Statistics,
National University of Kaohsiung\printead[presep={,\ }]{e3}}
\runauthor{Huang, Ing and Yu}
\end{aug}

\begin{abstract}
This paper addresses an important gap in time series analysis by establishing rigorous results on prediction and model selection for weakly stationary processes with conditionally heteroscedastic errors, using mean squared prediction error (MSPE) as the evaluation metric. A key contribution is the derivation of a negative moment bound for the minimum eigenvalue of the sample autocovariance matrix. This result enables an explicit asymptotic decomposition of the MSPE for least squares predictors into interpretable components reflecting model complexity, misspecification, and time-varying volatility. Leveraging this decomposition, we develop a model selection criterion that consistently identifies the MSPE-optimal subset autoregressive model, even in the presence of conditional heteroscedasticity and model misspecification. Simulation studies support the theoretical developments and demonstrate the practical effectiveness of the proposed method.
\end{abstract}

\begin{keyword}[class=MSC]
\kwd[Primary ]{62M10}
\kwd[; secondary ]{62M20}
\end{keyword}

\begin{keyword}
\kwd{Autocovariance matrix}
\kwd{conditional heteroscedasticity}
\kwd{least squares predictor}
\kwd{mean squared prediction
error}
\kwd{negative moment bound}
\kwd{uniform integrability}
\end{keyword}

\end{frontmatter}
%%%%%%%%%%%%%%%%%%%%%%%%%%%%%%%%%%%%%%%%%%%%%%
%% Please use \tableofcontents for articles %%
%% with 50 pages and more                   %%
%%%%%%%%%%%%%%%%%%%%%%%%%%%%%%%%%%%%%%%%%%%%%%
%\tableofcontents

\section{Introduction}\label{sec:1}
\numberwithin{equation}{section}
We consider a weakly stationary process defined as
\begin{equation}\label{AR mean part}
x_t=\sum_{i=0}^{\infty}\alpha_i\varepsilon_{t-i},
\end{equation}
where
$\{\varepsilon_t\}$ is
a martingale difference sequence
with respect to an increasing sequence of $\sigma$-fields
$\{\mathcal{F}_{t}\}$,
and satisfies
$0<E(\varepsilon_t^2)=\sigma_{\varepsilon}^2<\infty$
for all $t$,
with constant $\sigma_{\varepsilon}^2$.
The coefficients
$\{\alpha_i\}$ are real and satisfy
%$\alpha_i$ssatisfy
%{\color{blue} 
%[To be removed: If $E(\varepsilon_{t}^2)$
%is a constant, then $x_t$ is weakly stationary.
%This is ensured by the fact that $\{z_t\}$ an independent %sequence
%with constant variance in the case of ARCH models,
%where $\varepsilon_{t}=z_t\sqrt{f(z_{t-1}, z_{t-2}, %\ldots)}$.
%Note that $\{\varepsilon_{t}\}$ is automatically weakly %stationary
%because it is a martingale difference sequence
%regardless of whether $z_t$ is i.i.d. or not.
%}
\begin{equation}\label{alpha rate}
\alpha_0=1,\,\,|\alpha_i|=O(i^{-\iota}), \,\,\textrm{for some}\,\,
\iota >1,
\end{equation}
%$\sum_{i=0}^{\infty}|\alpha_i|< \infty$,
and
\begin{equation}\label{stationary char poly}
\sum_{i=0}^{\infty}\alpha_iz^i\neq0\ \textrm{for complex}\ |z|\leq1.
\end{equation}
Statistical inference for model \eqref{AR mean part},
including prediction and model selection, has been extensively studied in the literature.
However, most existing work focuses on the case where the conditional variance of
$\varepsilon_{t}$ given
$\mathcal{F}_{t-1}$, denoted
$\sigma_t^{2}=E(\varepsilon_{t}^{2}|\mathcal{F}_{t-1})$,
remains constant almost surely (a.s.) over time.
In contrast,
when $\varepsilon_{t}$
exhibits conditional heteroscedasticity, that is, $\sigma_t^{2}$
is a non-degenerate, 
$\mathcal{F}_{t-1}$-measurable random variable, rigorous results on estimation, prediction, and model selection, in terms of mean squared (prediction) error, are far more limited.
In particular, it remains unclear whether theoretical guarantees underlying conventional methods remain valid under conditional heteroscedasticity, primarily because moment properties of estimators and predictors in such settings are not well understood. 
 A concrete manifestation of this gap arises when approximating
\eqref{AR mean part}
with a finite-order autoregressive (AR) model: it is still unknown whether the least squares (LS) estimators of the AR coefficients possess finite moments, thereby hindering a rigorous analysis of the associated LS predictor's
mean squared prediction error (MSPE). This represents a significant open issue in the theory of time series analysis.







%is a weakly stationary process obeying
%For sequences of real numbers $\{a_m\}$ and $\{b_m\}$,
%$a_m=O(b_m)$ denotes the existence of a constant 
%$C>0$ and an integer $m_0\geq 1$ such that
%$|a_m| \leq C|b_m|$ for all $m\geq m_0$.

To address this gap, we study the behavior of the minimum eigenvalue of the sample autocovariance matrix, a quantity that plays a central role in analyzing prediction and model selection procedures based on LS estimation.
Given observations $x_1, \ldots, x_n$, define the sample autocovariance matrix of order $k$ as
\begin{eqnarray*}
\hat{\mathbf{R}}_n(k)=\frac{1}{n-k}\sum_{j=k}^{n-1}\mathbf{x}_j(k)\mathbf{x}^\top_j(k),
\end{eqnarray*}
where $k$ is a positive integer and $\mathbf{x}_j(k)=(x_j,\ldots,x_{j-k+1})^\top$. The exclusion of $x_n$ in $\hat{\mathbf{R}}_n(k)$ is made for notational convenience, allowing 
a direct link to the LS estimator
(see \eqref{ing1}). 
%However,
% whether $x_n$
% is included or excluded does not impact the estimator's asymptotic behavior. 
Let $\lambda_{\min}(\mathbf{M})$ denote the minimum eigenvalue of matrix $\mathbf{M}$.
In this article, we aim to establish, for any finite
integer $k\geq 1$
and real number $q>0$, the negative moment bound,
\begin{eqnarray}
\label{t.fisher.inverse moment}
E[\lambda_{\min}^{-q}(\hat{\mathbf{R}}_n(k))]=O(1), \,\, \mbox{as}\,\, n \to \infty,
\end{eqnarray}
under a conditional heteroscedasticity assumption, (CH), detailed in Section
\ref{sec:2}.



Negative moment bounds such as \eqref{t.fisher.inverse moment} have been used to
establish the uniform integrability of the LS estimator and play a pivotal role in time series prediction and model selection. See, for example, 
\cite{Fuller1981}, \cite{Kunitomo1985a},
%\cite{Bhansali1981},
\cite{Lewis1988}, \cite{Shaman1988},
\cite{Ing2003},  \cite{Scho2005},
\cite{Chan2011},  \cite{Ryan2013},
\cite{W2016}, and \cite{Chi2021}.
%in the special case where $\varepsilon_t$s are i.i.d.
%{\color{blue}try}
%These bounds have two important applications. The first one is to derive the mean squared error (MSE) and moment bounds of the least squares (LS) estimator,
%$$\hat{\bm{\beta}}_{n}(k)=\hat{\mathbf{R}}^{-1}_n(k)\frac{1}{n-k}\sum_{j=k}^{n-1}\mathbf{x}_j(k)x_{j+1};$$
%see Remark \ref{remark betahat MSE and bias} for more discussion. The second one is to develop an asymptotic
%expression for the mean squared prediction error (MSPE),
%$E\{x_{n+1}-\hat{x}_{n+1}(k)\}^2$, of the LS predictor
%$$\hat{x}_{n+1}(k)=\hat{\bm{\beta}}_{n}^\top(k)\mathbf{x}_n(k).$$
When $\{\varepsilon_t\}$ is a sequence of independent random variables,
the validity of
\eqref{t.fisher.inverse moment} has been well explored by \cite{Bhansali1991}, 
\cite{Papangelou1994},
\cite{Findley2002},
\cite{Ing2003}, and
\cite{Chan2011}, among others.
However, 
the assumption of independence excludes
many economic or financial time series that exhibits
conditional heteroscedasticity; see, e.g., \cite{Bollerslev1994} %or Chapter 3 of
and \cite{Tsay2010}.
To tackle this limitation, Assumption (CH) permits
$\{\varepsilon_t\}$ to follow a broad family of conditional heteroscedasticity models,
such as the generalized autoregressive
conditional heteroscedasticity (GARCH) model (\citealp{Bollerslev1986}), the GJR-GARCH model
(\citealp{Glosten1993}), the asymmetric power GARCH model (\citealp{Ding1993}), and the stochastic volatility (SV) model
(cf. \citealp{Taylor1982}, \citeyear{Taylor1986}; \citealp{Tsay2010}).
Consequently, establishing \eqref{t.fisher.inverse moment} under (CH)
greatly broadens the applicability of LS-based methods to more realistic and widely encountered time series environments.
%As clarified in Section \ref{sec:2}, (CH) includes
%Since we do not impose a specific form on $\varepsilon_t$,
%our framework is more flexible than
%as special cases.

%In the case that $\{\varepsilon_t\}$ is independent ($\sigma_t$ is constant), the problem of finding an
%asymptotic expression for
%$E\{x_{n+1}-\hat{x}_{n+1}(k)\}^2$
%has been well addressed; see \cite{Fuller1981} and \cite{Ing2003a} when the underlying model is a stationary $m$-th ($m<\infty$) order autoregressive
%(AR($m$)) model, and \cite{Ing2003} and \cite{Kunitomo1985a} when the underlying model is a stationary AR($\infty$) model.
%However, assuming independent (or conditional homoscedastic) errors may preclude
%many economic or financial time series that exhibits
%conditionally heteroscedastic behavior; see, e.g., \cite{Bollerslev1994} %or Chapter 3 of
%and \cite{Tsay2010}. Therefore, to enhance the
%applicability of $\hat{x}_{n+1}(k)$, we relax the frequently used independent
%assumption on the error sequence $\{\varepsilon_t\}$ by only considering a general
%framework that accommodates conditional heteroscedasticity as in (CH) given in Section \ref{sec:2}. Since we do not postulate a specific form for $\sigma_t$,
%our framework is more flexible than the popular generalized autoregressive
%conditional heteroscedasticity (GARCH) (\cite{Bollerslev1986}), GJR-GARCH
%(\cite{Glosten1993}), asymmetric power GARCH (\cite{Ding1993}), and stochastic volatility
%(SV) (cf. \cite{Clark1973}, \cite{Hull1987}, \cite{Melino1990}, \cite{Taylor1982} and \cite{Tsay2010}) model.


In settings where
$\{\varepsilon_t\}$ is an independent sequence, a commonly used approach to proving \eqref{t.fisher.inverse moment}
is to assume that
the marginal
distributions
of $\{\varepsilon_t\}$ satisfy some smoothness conditions; see, e.g.,
%\cite{Chi2021}, 
\cite{Papangelou1994},
\cite{Findley2002}, and
\cite{Ing2003a}.
This approach, however, is no longer valid for conditional heteroscedastic $\{\varepsilon_t\}$,
which can be
serially dependent in a highly convoluted fashion.
Furthermore, the marginal distribution of a conditional heteroscedastic process is often mathematically intricate
(\citealp{Francq19}), which makes
it difficult to verify related assumptions in practice.
To address this challenge, we
introduce smoothness conditions,
\eqref{zt density}--\eqref{zt density 2},
on the (marginal) distribution of 
the standardized innovation,
$z_t=\varepsilon_t/\sigma_t$,
rather than on
$\varepsilon_t$ itself. 
%where $\{z_t\}$ is assumed to be an independent sequence.
 A key result of our analysis is that \eqref{t.fisher.inverse moment} holds when $\{\varepsilon_t\}$ follows Assumption (CH) and $\{z_t\}$  fulfills \eqref{zt density}--\eqref{zt density 2},
which are readily met in many common applications.








In addition to their fundamental theoretical significance, our negative moment bounds play a crucial role in analyzing the asymptotic behavior of the MSPE,
$E\{x_{n+h}-\hat{x}_{n+h}(k)\}^2$,
of the $h$-step LS predictor, $\hat{x}_{n+h}(k)$, for $x_{n+h}$, when \eqref{AR mean part} is approximated by an
AR model of order $k \geq 1$,
with $h \geq 1$
and $\hat{x}_{n+h}(k)$ defined in \eqref{ing1}.
In particular, we show in Section \ref{sec:3.1}
that %derive an asymptotic expression for
the {\it second-order} MSPE,
\begin{eqnarray}
\label{ing2}
n[E\{x_{n+h}-\hat{x}_{n+h}(k)\}^2-E(\varepsilon_{n,h,k}^2)],
\end{eqnarray}
%has a finite limit, which
can be asymptotically decomposed
as the sum of three terms related to
{\it model complexity,
model misspecification, and conditional heteroscedasticity.}
Here, $\varepsilon_{n, h, k}$, defined in \eqref{AR(k) misspecified multi},
denotes the model error associated with the working AR($k$) model,
and $E(\varepsilon_{n,h,k}^2)$
is referred to as the {\it population}
MSPE.
While
 \cite{Kunitomo1985a} and
\cite{Ing2003a} 
have analyzed the asymptotic behavior of
$\eqref{ing2}$ under independent errors,
 extending this analysis to the conditionally heteroscedastic case has remained unaddressed due to the technical challenges involved in establishing
\eqref{t.fisher.inverse moment}.





%To the best of our knowledge, suc rigorous analysis the MSPE
%of the LS predictor in misspecified AR models with
%conditional heteroscedastic errors.

 %$E\{x_{n+h}-\hat{x}_{n+h}(k)\}^2$ (defined in Section \ref{sec:3.1}),
%under a fairly general setting on innovations that encompasses a rich class of conditional heteroscedastic models.


%It should also be noted that the asymptotic expression of $E\{x_{n+h}-\hat{x}_{n+h}(k)\}^2$ is established in a possibly misspecified setting, and hence further broadening the range of its applications. When the underlying
%model is specified as a finite-order AR model with the errors sequence obeys a GARCH or SV process, we also
%provide explicit forms for the asymptotic expressions of $E\{x_{n+h}-\hat{x}_{n+h}(k)\}^2$.
This work also makes a contribution to model selection. Specifically, let $J$ be a finite subset 
of $\mathbb{N}=\{1, 2,\ldots\}$ 
representing a subset AR model employed to predict $x_{n+h}$ based on lagged variables 
$\{x_{n+1-i}, i \in J\}$. Let
$\hat{x}_{n+h}(J)$ 
stand for the corresponding LS predictors, as described in \eqref{hat xn+h J}, and
denote the population MSPE, $E(\varepsilon_{t,h,J}^2)$, of $J$ by $f_h(J)$,
where $\varepsilon_{n, h, J}$ (see \eqref{xt+h J}) is the model error corresponding to model $J$. Let $\mathcal{J}$ be a finite family of candidate subset AR models. Define
\begin{eqnarray}
\label{M1}
{\cal J}_{1}(h)=\{\tilde{J}\in \mathcal{J}: f_h(\tilde{J})=\min\limits_{J\in \mathcal{J}}f_h(J)\}
\end{eqnarray}
as the set of candidate models having the smallest population MSPE. Also, define
\begin{eqnarray}
\label{ing4}
g_{h}(J)=\lim_{n \to \infty}
n\{E\{x_{n+h}-\hat{x}_{n+h}(J)\}^2-f_h(J)\},
\end{eqnarray}
provided the limit exists, and
\begin{eqnarray}
\label{M2}
{\cal J}_{2}(h) = \{\tilde{J}\in {\cal J}_{1}(h): g_h(\tilde{J})=\min\limits_{J\in {\cal J}_{1}(h)}g_h(J)\}.
\end{eqnarray}
Our objective is to select a model 
$\hat{J}$ 
through a data-driven approach so that
\begin{eqnarray}
\label{ing3}
\lim_{n \to \infty}
P(\hat{J} \in {\cal J}_{2}(h))=1.
\end{eqnarray}
Essentially, our aim is to asymptotically identify the model that yields the lowest population MSPE
when ${\cal J}_{1}(h)$ contains only one model.
If ${\cal J}_{1}(h)$
includes multiple models, we then prefer the one with the smallest second-order MSPE among
those in ${\cal J}_{1}(h)$.
In the former case, this objective can be readily achieved by comparing residual mean-squared errors (or, equivalently, $R$-squared values). %when the target model is correctly specified. 
In the latter case, the goal remains tractable using BIC-type criteria, provided that
$\{\varepsilon_t\}$
is conditionally homogeneous {\it and}
 at least one model in
${\cal J}$ is correctly specified. 
%(or, equivalently, includes the smallest correct model for $h$-step prediction). 
However, 
when ${\cal J}_{1}(h)$ contains multiple models
and either 
$\{\varepsilon_t\}$ 
exhibits conditional heteroscedasticity or all models in
${\cal J}$ are misspecified,
identifying the optimal model becomes substantially more challenging.
Further discussion is provided in Section
\ref{sec:3.2}.



Recently, \cite{Hsu2019} partially addressed this challenge--under the assumption of independent errors--by introducing a misspecification-resistant information criterion (MRIC).
For candidate model $J$, its MRIC for $h$-step
prediction is defined by:
\begin{equation}\label{MRIC0}
\textrm{MRIC}_{h}(J)=\hat{\sigma}^2_{h}(J)+\frac{C_n}{n}\hat{g}_h(J),
\end{equation}
where 
$\hat{\sigma}^2_{h}(J)$ 
is a $\sqrt{n}$-consistent estimator
of $f_{h}(J)$, 
$\hat{g}_h(J)$ is a consistent estimator
of $g_h(J)$, and $C_n$ approaches $\infty$
at a suitable rate. 
%%{ to be continued....}
Unlike conventional information criteria
that penalize model dimension,
MRIC directly penalizes the second-order MSPE through the term
$C_n\hat{g}_h(J)/n$.
%Note that the construction of $\hat{g}_h(J_l)$ in
%\cite{Hsu2019} is based on a concise asymptotic expression for $g_h(J_l)$ that holds regardless of whether $J_l$ is correctly specified.
Assuming independent errors, 
\cite{Hsu2019} showed
that the model minimizing MRIC$_{h} (\cdot)$ over all candidates achieves
\eqref{ing3}, even when all $J$ are misspecified and ${\cal J}_{1}(h)$ contains multiple models.
However, whether this result extends to conditionally heteroscedastic errors remains unclear, as the existence of the limit in 
\eqref{ing4} 
has not been established in that context.
In this study, we confirm the existence of the limit 
using \eqref{t.fisher.inverse moment}, and demonstrate that the MRIC proposed by \cite{Hsu2019}
achieves \eqref{ing3}
regardless of model misspecification, the number of models in
${\cal J}_{1}(h)$, or
the presence of conditional heteroscedasticity in $\{\varepsilon_t\}$.



%the expression for $g_h(J)$ from \cite{Hsu2019} extends to cases of conditionally heteroscedastic errors.
%This discovery
%, along with our consistent estimate of $g_h(J)$ %established under conditional heteroscedasticity, 
%enables us to obtain a comprehensive model selection result, wherein 

%The major contributions of this article are that we develop  and apply it to\\
%{(This paragraph is used to be connected to the section of correctly specified examples)}\\
%Our results reveal that the conditional heteroscedasticity in $\{\varepsilon_t\}$ has a substantial effect on the one-step and multi-step MSPEs of the LS predictors, which, in turn, makes their asymptotic expressions quite different from the corresponding ones under the conventional independent assumption on $\{\varepsilon_t\}$. For instance, when the AR($k$) model is a correctly specified model, the limit of the one-step MSPE, $n[E\{x_{n+1}-\hat{x}_{n+1}(k)\}^2-\sigma^2]$, with $\sigma^2=E(\varepsilon_1^2)$, is known to be $k\sigma^2$ for the case of independent errors; see, e.g., \cite{Fuller1981} and \cite{Ing2003a}. However, if $\{\varepsilon_t\}$ exhibits conditional heteroscedasticity, the limit of $n[E\{x_{n+1}-\hat{x}_{n+1}(k)\}^2-\sigma^2]$ can be decomposed into a sum of two terms, $k\sigma^2$ and $\mathrm{tr}\{\mathbf{R}^{-1}(k)\mathbf{L}^\star_{0,1}(k)\}$ (defined in Section \ref{sec:3.1}), where the latter term depends on the degree of conditional heteroscedasticity and is generally away from 0; see Section \ref{sec:3.2} for more discussion.
%details.%can be much larger than 0. %far away from.
%This indicates the intriguing result that, depending on the model structure of $\varepsilon_t$, the limit of $n[E\{x_{n+1}-\hat{x}_{n+1}(k)\}^2-\sigma^2]$ under \eqref{AR mean part}--(CH)(i) may even exceed the ones corresponding to nonstationary AR processes with independent innovations; see Section 3.1 for more details.

The remainder of the paper is organized as follows.
The negative moment bound
\eqref{t.fisher.inverse moment} is established
in Theorems \ref{tfisher} and \ref{t SV negative bound}
of
Section \ref{sec:2}. %Armed with the result obtained in Section \ref{sec:2},
Section \ref{sec:3.1} derives an asymptotic expression for the second-order MSPE in \eqref{ing2}, which is extended in Section \ref{sec:3.2} to subset AR models under possible misspecification. The asymptotic validity of the MRIC is also confirmed in Section \ref{sec:3.2}. Sections
\ref{app:A} and \ref{appA00}
provide further discussion of the assumptions used in Sections \ref{sec:3.1} and \ref{sec:3.2}. Section \ref{sec:4} presents numerical simulations that support the theoretical findings of Section \ref{sec:3}. The proof of the main result, Theorem  \ref{tfisher}, is given in Section \ref{sec:5} and is divided into four steps. Additional technical details are provided in Appendix  \ref{appA} and the Supplementary Material. The latter also contains further numerical analyses regarding MSPE and model selection. 
We conclude this section with some notation used throughout the paper. For real
numbers $x$ and $y$, 
$x\wedge y= \min\{x,y\}$ and $x\vee y=\max\{x,y\}$. 
%and $\lfloor z\rfloor$ denotes the largest integer $\leq z$.
For square matrix $\mathbf{A}$,
$\|\mathbf{A}\|$ and $\mathrm{tr}(\mathbf{A})$
stand for its spectral norm and trace, respectively.
For vector $\bm{l}$,
$\|\bm{l}\|$ denotes its Euclidean norm. 
%For square matrix $A$, $\mathrm{tr}(A)$ stands for the trace operator.
%For a set $G$, $G^{o}$ denotes its interior.
For a sequence of random variables, $X_{t-1}, X_{t-2}, \ldots$, 
$\sigma(X_{t-1}, X_{t-2}, \ldots)$ represents the $\sigma$-algebra they generate.  The symbol '$\equiv$' denotes a definition. 



\section{Negative Moment Bounds for minimum eigenvalues}\label{sec:2}
%We begin this section by %providing negative moment bounds for $\hat{\mathbf{R}}_n(k)$, which is one
%of the key ingredients in deriving an asymptotic expression for the MSPE of $\hat{x}_{n+1}(k)$.
Let $\{{\cal F}_t\}$ be an increasing sequence of sub-$\sigma$-fields
on the probability space $(\Omega, {\cal F}, P)$,
and
$\{\varepsilon_t, {\cal F}_t\}$ be a martingale difference sequence.
We introduce an assumption of conditional heteroscedasticity. % on $\{\varepsilon_t, {\cal F}_t\}$.
\begin{description}
\item[\textbf{Assumption (CH)}.]
\noindent \normalfont There exist a $\mathcal{F}_t$-measurable random variable $z_t$,
with $E(z_t)=0$ and $E(z_t^2)=1$, and a
non-negative $\mathcal{F}_{t-1}$-measurable random variable $\sigma_t$ such that
\begin{eqnarray}
\label{varepsilon basic form}
\varepsilon_t=\sigma_tz_t.
\end{eqnarray}
Moreover, $z_t$ is independent of $\mathcal{F}_{t-1}$ and
\begin{itemize}
\item[(i)]
$\sigma_t$ is a $\sigma(\varepsilon_{t-1},\varepsilon_{t-2},\ldots)$-measurable random variable and for some 
$0<c_0<\infty$,
\begin{eqnarray}
\label{ing101}
\sigma_t \geq c_0 \,\,\mbox{a.s.},
\end{eqnarray}
\item[or]
\item[(ii)]
$\{z_t\}$
and $\{\sigma_t\}$ are
independent sequences, and for any $\theta_0>0$,
\begin{eqnarray}
\label{ing102}
\sup\limits_{-\infty<t<\infty}E\sigma_t^{-\theta_0}<\infty.
\end{eqnarray}


\end{itemize}
\end{description}

%A few comments of (CH) are in order.
%We demonstrate that
Under Assumption CH(i) and by
Theorem 1.4.5 of \cite{Chow1997}, there exists a measurable function
$\tilde{\sigma}_{t}$ on $\mathbb{R}^\infty$ such that
$\sigma_t=\tilde{\sigma}_{t}(\varepsilon_{t-1},\varepsilon_{t-2},\ldots)$ a.s. 
%Moreover, 
%$E(\tilde{\sigma}^2_{t}(\varepsilon_{t-1},\varepsilon_{t-2},\ldots))=\sigma_\varepsilon^2$, where $\sigma_\varepsilon^2$ is defined below \eqref{AR mean part}. 
Assumption (CH) includes
many conditional heteroscedastic models
as special cases.
%two major types of conditional heteroscedastic models. %most popular conditional heteroscedastic models.
For example, (CH)(i) is fulfilled by the asymmetric power GARCH model (\cite{Ding1993}),
which is \eqref{varepsilon basic form} with $z_t$ being i.i.d. random variables and
\begin{equation}\label{apGARCH}
 \sigma_t^\mu=  \varphi_{0}+\sum_{i=1}^{p^\prime}\varphi_{i}(|\varepsilon_{t-i}|-\lambda_i\varepsilon_{t-i})^\mu+\sum_{j=1}^{q^\prime}\psi_{j}\sigma_{t-j}^\mu,
\end{equation}
where $p^\prime,q^\prime\in\mathbb{N}$, $\varphi_{0}>0$, $\mu>0$, $\varphi_{i}\geq0$ and
$|\lambda_i|<1$ for $1\leq i \leq p^\prime$,
$\psi_{j} \geq 0$ for $1\leq j \leq q^\prime$, and
\begin{equation}\label{apGARCH stationary}
\sum_{i=1}^{p^\prime}\varphi_{i}E(|z_1|-\lambda_i z_1)^\mu+\sum_{j=1}^{q^\prime}\psi_{j}<1.
\end{equation}
%Note that \eqref{apGARCH stationary}
%ensures the stationarity of $\{\varepsilon_t\}$; see
Let $\mathcal{F}_{t}=\sigma(z_{t},z_{t-1},\ldots)$.
Then, %$z_t\in\mathcal{F}_{t}$,
$z_t$ is independent of
$\mathcal{F}_{t-1}$
and by
\eqref{apGARCH stationary}, $\sigma_t \in\mathcal{F}_{t-1}$ (see  \cite{Ling2002}).
In addition,
$\sum_{j=1}^{q^\prime}\psi_{j}<1$ ensures
\begin{equation}\label{sigma is a function of lag variables}
\sigma_t^\mu=(1-\sum_{j=1}^{q^\prime}\psi_{j}B^j)^{-1}\{\varphi_{0}+\sum_{i=1}^{p^\prime}\varphi_{i}(|\varepsilon_{t-i}|-\lambda_i\varepsilon_{t-i})^\mu\},
\end{equation}
where $B$ denotes the back-shift operator,
and hence
$\sigma_t\in\sigma(\varepsilon_{t-1},\varepsilon_{t-2},\ldots)$.
Finally, \eqref{apGARCH} implies
$\sigma_t\geq \varphi_{0}^{1/\mu}$ a.s.
Thus, \eqref{ing101} holds true.
It is also worth mentioning that the
GARCH($p^\prime,q^\prime$) process is a special case of \eqref{apGARCH}
with
$\mu=2$ and $\lambda_i=0$ for all $i$.


%On the other hand, (CH)(ii) is valid for SV models. To see this, consider
Another example is an SV model considered in Chapter 3 of \cite{Tsay2010}, 
which is \eqref{varepsilon basic form} with $z_t$ being i.i.d. random variables and
$\log \sigma_t^2$ satisfying
\begin{equation}\label{SV model}
      (1-\tilde{a}_1B-\cdots-\tilde{a}_{\tilde{p}}B^{\tilde{p}})\log(\sigma_t^2)=\tilde{a}_0+v_t,
\end{equation}
where $\{v_t\}$ is a sequence of i.i.d. mean zero normal random variables independent of $\{z_t\}$,
$\tilde{p}\in\mathbb{N}$, $\tilde{a}_0$ is a
constant, and $1-\sum_{i=1}^{\tilde{p}}\tilde{a}_iz^i\neq0$, for all $|z|\leq1$. 
%see Chapter 3 of \cite{Tsay2010} for more details. %{(May 8 at 2:50 pm)}
Letting $\mathcal{F}_{t}=\sigma(z_t,z_{t-1},\ldots,v_{t+1},v_t,\ldots)$,
we now argue that this SV model is a special case of (CH)(ii).
First, it is clear that $z_t\in\mathcal{F}_{t}$, $\sigma_t\in\mathcal{F}_{t-1}$,
and $\{z_t\}$ is independent of $\{\sigma_t\}$. Moreover, %$v_t$ are i.i.d. normal and $1-\sum_{i=1}^{\tilde{p}}\tilde{a}_iz^i\neq0$, for all $|z|\leq1$,
\eqref{ing102} is ensured by the fact that
$\sigma_t^2$ are identically distributed log-normal random variables.

It is important to highlight that because 
(CH) does not prescribe any particular parametric or non-parametric model for $\sigma_t$, the condition is considerably more flexible than the previously mentioned special cases. 
%, the second part of (CH)(ii). Note that the second part of (CH)(ii) also holds when $v_t$ are bounded random variables.
%Before providing one of the main results of this section, we further assume without loss of generality that $z_t$ have a common variance of 1.
%$z_t$ have common mean 0 and common variance 1.
%where $c_0$ is some positive constant. % and $z_t$ have common variance 1.
%Throughout this paper, we assume without loss of generality that $z_t$ have a common variance of 1 and $k$ is a finite positive integer.
The main result of this paper is presented in Theorem \ref{tfisher}, in which
(CH)(i) is assumed and
%In the rest of the paper, we assume that for each $$,
$z_t, t \in \mathbb{Z}=\{\cdots, -1,0,1,\cdots\}$, are allowed to have density functions,
$\phi_t(\cdot)$, with respect to the Lebesgue measure.
The proof of Theorem \ref{tfisher} can be found in Section \ref{sec:5}. %of the Supplementary Material. 
To state the theorem, 
define $\zeta_{t,x}(c)=(1/c)\phi_t(x/c)$, where 
$x \neq 0$ and
$c>0$.

 % provides one of the main results of this section.
%To provide one of the main results of this section, we further consider a distributional assumption.
%For all $t$,
\begin{theorem}\label{tfisher}
Assume \eqref{AR mean part}--\eqref{stationary char poly}
and {\rm (CH)(i)}. 
Suppose for any 
%$t \in \mathbb{Z}$and 
$0< x< y<\infty$, %that $\phi_t(\cdot)$ is bounded above and symmetric about zero and obeys
\begin{align}\label{zt density}
\begin{split}
\phi_t(y)\leq \phi_t(x)\,\,\mbox{and}\,\,\phi_t(-y)\leq \phi_t(-x).
\end{split}
\end{align}
In addition, 
for any $0<\delta \leq 1$, there exist a
finite positive constant $M_\delta$
and a pair ($\bar{\theta} ,\bar{C} $), independent of $\delta$, satisfying $0<\bar{\theta} \leq 1$ and $0<\bar{C}<\infty$,
such that
\begin{eqnarray}
\label{ing104}
\sup_{-\infty<t<\infty}\{\phi_t(\delta)+\phi_t(-\delta)\}\leq M_\delta,\,\,\sup_{-\infty<t<\infty}
\int_{-\delta}^{\delta}\phi_t(x)dx\leq\bar{C}\delta^{\bar{\theta}}.
\end{eqnarray}
Moreover, for any $t\in \mathbb{Z}$ and $x \neq 0$, there exists
a positive number $m_t(x)$ obeying $c_1|x| < m_t(x)< c_2|x|$, 
%$c_1|x| < m^{+}_t(x), m_t^{-}(x)< c_2|x|$,
where $0<c_1<c_2<\infty$ are constants that do not depend on $t$ or $x$, such that
\begin{align}
\begin{split}
\label{zt density 0}
\zeta_{t,x}(m_t(x))=\sup_{c>0} \zeta_{t,x}(c),
\end{split}
\end{align}
and 
\begin{equation}
\label{zt density 2}
\begin{split}
\zeta_{t,x}(c) \,\,\mbox{is non-increasing for} \,\,
c>m_t(x).
\end{split}
\end{equation}
Then, \eqref{t.fisher.inverse moment} follows.
\end{theorem}

%Since we assume that $k$ is fixed as $n$ increases, \eqref{t.fisher.inverse moment} directly follows from \eqref{t.fisher.inverse moment k dependence}. A discussion of the case where $k = k_n$ grows to infinity with $n$ is provided in Remark \ref{rm2.2} at the end of this section. 
%\begin{remark}\label{rm.sigma t2 variance}
%In fact, \eqref{t.fisher.inverse moment k dependence} holds without the need to assume that $\{x_t\}$ is weakly stationary or that $E(\varepsilon_t^2)$ is a finite constant. Instead, it suffices to assume that $\sup_{-\infty<t<\infty}E\sigma_t^2<\infty$, under which \eqref{t.fisher.inverse moment k dependence} remains valid.
%\end{remark}
Condition
\eqref{zt density}--\eqref{zt density 2},
in particular 
\eqref{zt density 0}
and 
\eqref{zt density 2},
enable us
to show that for any unit vector $(l_1,\ldots, l_k)^{\top}$ in $\mathbb{R}^{k}$
and any integer $t$,
the conditional distribution of $\sum_{j=1}^{k}l_j\varepsilon_{t+1-j}$
given $\{\varepsilon_{s}, s\leq t-k\}$
is sufficiently smooth near the origin. This is one of the most critical properties needed for proving 
\eqref{t.fisher.inverse moment}.
To establish this property, we construct 
a conditional joint probability density function (pdf)
of $(\varepsilon_{t}, \ldots, \varepsilon_{t-k+1})$
given $\sigma(\varepsilon_{s}, s\leq t-k)$,
based on the marginal densities
$\phi_t(\cdot), \ldots, \phi_{t-k+1}(\cdot)$.
However, the resulting conditional joint pdf consists of 
$k$ multivariate components with 
highly entangled arguments, making direct analysis analytically intractable. To overcome this, we invoke the smoothness assumptions in
\eqref{zt density}--\eqref{zt density 2}
 to construct a suitable univariate envelope function that bounds each component in the joint pdf. This strategy effectively decouples the multidimensional dependence, allowing us to derive tractable bounds on the conditional joint pdf and establish the desired regularity.
 For further details, refer to Steps 2 and 3 in the proof of 
 Theorem \ref{tfisher}.
%\eqref{pt.fisher.7.5} in Section \ref{sec:5}. %of the Supplementary Material.

\vspace{0.3cm}
Some additional remarks regarding 
Theorem \ref{tfisher} are in order.

\begin{remark}\label{rm2.0}
Assumptions \eqref{zt density}--\eqref{zt density 2} are fulfilled by many symmetric density functions encountered in common practice.
For example,
they are satisfied with $M_\delta=\bar{C}=\sqrt{2/\pi}$, $\bar{\theta}=1$, and
$m_t(x)=|x|$,
when $\phi_t(\cdot)$ is the standard normal density function,
and with
$M_\delta=\bar{C}=2\sigma_\nu\Gamma\{(\nu+1)/2\}/\{\sqrt{\nu\pi}\Gamma(\nu/2)\}$,
$\bar{\theta}=1$, and
$m_t(x)=\sigma_\nu |x|$,
when $\phi_t(\omega)=\sigma_\nu J(\sigma_\nu \omega)$
is the ``normalized" density function for the $t$-distribution
with $\nu$ degrees of freedom,
where $-\infty<\omega<\infty$, $\nu>2$,
$\sigma_\nu=\sqrt{\nu/(\nu-2)}$,
$J(\cdot)$ is the
density function for the $t$-distribution
with $\nu$ degrees of freedom, and
$\Gamma(\cdot)$ is the gamma function.
Assumptions \eqref{zt density}--\eqref{zt density 2} even hold when $\{\phi_t(\cdot)\}$
are unbounded. To see this, assume that for each $t \in \mathbb{Z}$,
$$\phi_t(\omega)=\frac{1}{2\Gamma(\xi)}|\omega|^{\xi-1}e^{-|\omega|}$$
is a symmetric Gamma density function, where $-\infty< \omega <\infty$ and $0<\xi<1$. Then,
\eqref{zt density}--\eqref{zt density 2} are satisfied with $M_\delta=\delta^{\xi-1}/\Gamma(\xi)$, $\bar{C}=1/\{\Gamma(\xi)\xi\}$, $\bar{\theta}=\xi$,
and $m_t(x)=|x|/\xi$.
\end{remark}

\begin{remark}\label{rm2.1}
The centrally monotonic property of 
$\phi_t(\cdot)$, as described in
\eqref{zt density}, can be readily relaxed to accommodate more general $\phi_t(\cdot)$, such as the mixture normal density, 
at the cost of replacing \eqref{ing104} with a slightly stronger assumption. For more details, refer to Appendix \ref{appA}.
\end{remark}



If
$\{z_t\}$ and $\{\sigma_t\}$ are two independent sequences,
then the distributional assumption,
\eqref{zt density}--\eqref{zt density 2}, on $\{z_t\}$,
and the lower bound condition, \eqref{ing101},
on $\{\sigma_t\}$
can be substantially relaxed. In particular,
when (CH)(i) is replaced by (CH)(ii),
we show in Theorem \ref{t SV negative bound} below that
\eqref{t.fisher.inverse moment} still follows, provided that
\eqref{zt density}--\eqref{zt density 2}
 are replaced by the weaker condition:
for some positive numbers $\rho$, $\eta$, and $\bar{M}$,
and all $|x-y|\leq\eta$,
%In the case of independent errors ($\sigma_t$ is constant), the following smoothness condition
%on the distribution functions of innovations $z_t$, $F_t(\cdot)$, is usually required to obtain .
\begin{equation}\label{smoothness condition zt}
\sup_{-\infty<t<\infty}
|F_t(x)-F_t(y)|\leq \bar{M}|x-y|^\rho,
\end{equation}
where $F_t(\cdot)$ denotes the distribution function of $z_t$.
The proof of Theorem \ref{t SV negative bound} is provided in Appendix \ref{appA}.
We note that \eqref{smoothness condition zt},
including the second part of \eqref{ing104}
as a special case, has been used by
\cite{Ing2003} to derive
\eqref{t.fisher.inverse moment}
in situations where $\{\varepsilon_t\}$
is a sequence of independent random variables.
Furthermore, \eqref{ing102} in (CH)(ii) is weaker than
\eqref{ing101} in (CH)(i).


%can be easily fulfilled if $z_t$ are i.i.d with bounded density (see Lemma 4 of ).
%Note that compared to the first part of (CH)(i), the first part of (CH)(ii) is much stronger.
%As a tradeoff, the distributional assumptions \eqref{zt density}--\eqref{zt density 2}
%are weakened to \eqref{smoothness condition zt} in Theorem \ref{t SV negative bound}.
%In addition, the negative moment assumption stated in (CH)(ii), \eqref{ing102},
%enables us to control the local behavior of $\sigma_t$ near 0, and can be viewed as a weaker version of \eqref{ing101}.
%It is used to provide an upper bound for the conditional probability
%\begin{equation}\label{conditional probability phi}
%P(|\bm{l}_j^\top\bm{\phi}_{ik}|\leq3u^{-\frac{1}{2q}}|\varepsilon_{s},s\leq ik-k),
%\end{equation}
%which is the most crucial step in proving \eqref{t.fisher.inverse moment}, where $\bm{l}_j^\top\bm{\phi}_{ik}$ is a linear combination of $x_{l}$, %$l=ik-k+1,\ldots,ik$, defined in the proof of Theorem \ref{tfisher} in Appendix \ref{appA}. See \cite{Findley2002}, \cite{Ing2003a}, and \cite{Ing2003} for the use of condition \eqref{smoothness condition zt} in the case of i.i.d. errors. However, such an argument does not hold here because condition \eqref{smoothness condition zt} is now insufficient to derive an upper bound for \eqref{conditional probability phi} due to the conditional heteroscedasticity in $\{\varepsilon_t\}$.
%their arguments do not hold here because condition \eqref{smoothness condition zt} is now insufficient to derive an upper bound for \eqref{conditional probability phi} due to the conditional heteroscedasticity in $\{\varepsilon_t\}$.
%To conquer this dilemma, we consider using \eqref{t.fisher.moment}--\eqref{zt density 2}, where the last three distributional assumptions together are slightly stronger than \eqref{smoothness condition zt} and the last moment assumption allows us to control the tail distribution of $\sigma_t$, instead to deal with %tackle
%\eqref{conditional probability phi} in the framework (CH)(i).
%The next theorem shows that \eqref{t.fisher.inverse moment} is also true for $\varepsilon_t$ satisfying (CH)(ii).

\begin{theorem}\label{t SV negative bound}
Assume \eqref{AR mean part}--\eqref{stationary char poly}, {\rm (CH)(ii)}, and
\eqref{smoothness condition zt}. Then, \eqref{t.fisher.inverse moment} holds.
\end{theorem}

%A comparison of Theorem \ref{tfisher} with Theorem \ref{t SV negative bound} is given below.
%A few comments are in order regarding Theorem \ref{t SV negative bound}. To start with, note that compared to the first part of (CH)(i), the first part of (CH)(ii) is much stronger. %As a tradeoff, the distributional assumption \eqref{zt density} and the moment assumption \eqref{t.fisher.moment} used in Theorem \ref{tfisher} are weakened to \eqref{smoothness condition zt} and \eqref{t.SV negative bound.1}, respectively, in Theorem \ref{t SV negative bound}.


%Moreover, in Theorem \ref{t SV negative bound}, since the third part of (CH)(i) is not assumed, the second part of (CH)(ii) is used as an alternative. Note that this negative moment assumption enables us to control the local behavior of $\sigma_t$ near 0.

%Although it is feasible to generalize the results obtained in





%Theorem \ref{t SV negative bound}
%can likewise be extended to accommodate growing
%$k=k_n$.
%Specifically, under the assumptions of Theorem \ref{t SV negative bound}, it can be shown that
%for any $q, \theta>0$ and $k=k_n=o(n^{1/2})$,
%\begin{equation}\label{SV inverse moment k}
%E[\lambda_{\min}^{-q}(\hat{\mathbf{R}}_n(k))]\leq C_{3}k^{(2+\theta)q}, 
%\end{equation}
%where $C_{3}$ is a positive constant independent of $k$ and $n$. This generalizes Lemma 1 of \cite{Ing2003} to settings
%with SV errors. Notably, due to the independence between $\{z_t\}$ and $\{\sigma_t\}$, the 
%bound in \eqref{SV inverse moment k} 
 %is considerably sharper than that in
 %\eqref{GARCH inverse moment k}. 
 %Applying the same techniques used in
 %the proof of Theorem 2 of \cite{Ing2003}, we conclude: (b1) If $k^{2+\delta}=O(n)$ for some $\delta>0$,
 %and
 %$\sup_{-\infty<t<\infty}E|\sigma_t|^{q_1}<\infty$ and $\sup_{-\infty<t<\infty}E|z_t|^{q_1}<\infty$ for all $q_1>0$,
 %then
%\eqref{t.fisher.inverse moment} holds under the condition , and (b2)
%when $k^{6+\delta}=O(n)$ for some $\delta>0$,
%\eqref{t.fisher.inverse moment} holds for $0<q<q_1$, provided that $\sup_{-\infty<t<\infty}E|\sigma_t|^{2(q_1\vee2)}<\infty$ and $\sup_{-\infty<t<\infty}E|z_t|^{2(q_1\vee2)}<\infty$ for some $q_1>0$. 
%Finally, as shown in the next section, in the analysis of MSPEs, \eqref{t.fisher.inverse moment} often needs to hold for a large $q$.
%Consequently, the moment conditions in (a) on 
%$\sigma_t$ appear restrictive for GARCH-type errors in practical applications (cf. \cite{Ling2002} and \cite{Ling2002joe}). %Therefore, we do not pursue the scenario of increasing $k$ in this work.




















\section{Prediction and Model Selection}\label{sec:3}
In Section \ref{sec:3.1}, we utilize Theorems
\ref{tfisher} and \ref{t SV negative bound} to derive an asymptotic expression for the multistep MSPE of the LS predictor, as detailed in Theorem \ref{tMSPE}. We concentrate on the "direct" multistep prediction using the LS method due to its robust performance in scenarios where the working AR model is misspecified. This approach is supported by discussions in \cite{Chevillon2007}, \cite{Ing2003a}, and \cite{OJ05}. Additionally, 
%we are drawn to the LS method because assumption (CH) is too general to be captured by a specific conditionally heteroscedastic model. Therefore, 
a more complex estimation method based on a specific conditionally heteroscedastic model could face greater risks of model misspecification, beyond the misspecification of the mean function.
%it difficult to specify using existing parametric or non-parametric models
%for considering the LS method is , which is difficult to
%model using any specific parametric or non-parametric approach.
%For any fixed $h$ and $k$, an asymptotic expression for
%the second-order MSPE, \eqref{ing2},
%is given in Theorem \ref{tMSPE}.
With the help of Theorem \ref{tMSPE},
Theorem \ref{t MRIC} of
Section \ref{sec:3.2} addresses the problem of subset selection
and establishes the validity of MRIC in possibly
misspecified AR models with conditionally heteroscedastic errors.
The proofs of Theorem \ref{tMSPE} and
\ref{t MRIC} are offered in Sections \ref{appB} and \ref{appC} of the Supplementary Material, respectively. %In Sections \ref{app:A} and \ref{appA00}, we discuss further about the details of how conditions imposed in Theorem \ref{tMSPE} and \ref{t MRIC} can be fulfilled by several commonly used weakly dependence models, including those whose serial dependencies are characterized using functional dependence measures (\cite{Wu2005Main}), strong mixing conditions, and 

\subsection{Asymptotic Expressions for the MSPEs}\label{sec:3.1}
Define
$\bm{\beta}_h(k)=\arg\min_{\mathbf{c}\in\mathbb{R}^k}E\{x_{t+h}-\mathbf{c}^\top\mathbf{x}_{t}(k)\}^2=\mathbf{R}^{-1}(k)E(\mathbf{x}_t(k)x_{t+h})$, where $\mathbf{R}(k)=E(\mathbf{x}_t(k)\mathbf{x}_t(k)^\top)$.
Then, $\bm{\beta}_h^\top(k)\mathbf{x}_{t}(k)$
is the best linear predictor of $x_{t+h}$
based on $\mathbf{x}_{t}(k)$,
with $\varepsilon_{t,h,k}=x_{t+h}-\bm{\beta}_h^\top(k)\mathbf{x}_{t}(k)$ representing the associated 
prediction error. Note that 
$E(\mathbf{x}_{t}(k)\varepsilon_{t,h,k})$
is the $k$-dimensional zero vector.
In what follows, the model
\begin{equation}\label{AR(k) misspecified multi}
x_{t+h}=\bm{\beta}_h^\top(k)\mathbf{x}_{t}(k)+\varepsilon_{t,h,k},
\end{equation}
is referred to as the $h$-step-ahead predictive model. 
Given \eqref{AR(k) misspecified multi}, the $h$-step LS predictor of $x_{n+h}$ 
is 
\begin{eqnarray}
\label{ing1}
\hat{x}_{n+h}(k)=\hat{\bm{\beta}}_{n,h}^\top(k)\mathbf{x}_n(k),
\end{eqnarray}
where
$$\hat{\bm{\beta}}_{n,h}(k)=(n-h-k+1)^{-1}\hat{\mathbf{R}}^{-1}_{n,h}(k)\sum_{j=k}^{n-h}\mathbf{x}_j(k)x_{j+h},$$
with
$$\hat{\mathbf{R}}_{n,h}(k)=(n-h-k+1)^{-1}\sum_{j=k}^{n-h}\mathbf{x}_j(k)\mathbf{x}^\top_j(k).$$
Since the value to be predicted, $x_{n+h}$, and the observed data, $x_1, \ldots, x_n$,
belong to the same realization, this type of prediction is referred to as the
``same-realization" prediction and is somewhat different from the ``independent-realization" prediction
in which the value to be predicted is independent of the observed data; see, e.g.,
\cite{Akaike69},
\cite{Bhansali1981},
\cite{Findley2002}, and
\cite{Scho2005}.

The next theorem provides
an asymptotic expression for $E\{x_{n+h}-\hat{x}_{n+h}(k)\}^2$. To streamline the presentation, we hereafter refer to the assumptions for GARCH-type errors (CH(i) and \eqref{zt density}--\eqref{zt density 2}) as Condition (G), and those for SV-type errors (CH(ii) and \eqref{smoothness condition zt}) as Condition (S). Conditions (G) and (S) are both used for establishing the negative moment bound in \eqref{t.fisher.inverse moment}, as discussed in Section \ref{sec:2}.

\begin{theorem}\label{tMSPE}
{\rm (a)} Assume \eqref{AR mean part}--\eqref{stationary char poly}, with $\iota$ in \eqref{alpha rate} satisfying $\iota>3/2$, and Condition (G). 
Let $\{\varepsilon_t\}$ be fourth-order weakly stationary, and suppose there exists a small positive number $\delta$ such that
\begin{equation}\label{tMSPE.moment}
\sup\limits_{-\infty<t<\infty}E|\varepsilon_t|^{6+\delta}<\infty.
\end{equation}
Furthermore, 
there exist 
a positive constant $C$
and a small positive number $\delta^{*}$
such that
for any sequences $\{\ubar{m}_{n}\}$ and $\{\bar{m}_{n}\}$ with $1\leq \ubar{m}_{n}\leq \bar{m}_{n}\leq n$,
\begin{equation}\label{tMSPE.moment2}
\sup\limits_{-\infty<s<\infty}E|\frac{1}{\sqrt{\bar{m}_{n}-\ubar{m}_{n}+1}}\sum_{t=\ubar{m}_{n}}^{\bar{m}_{n}}(\varepsilon_{t+s}^2-E(\varepsilon_{t+s}^2))|^{3+\delta^\star}\leq C,
\end{equation}
and
\begin{equation}\label{tMSPE.conditional moment}
\sup\limits_{-\infty<t<\infty} E|E(\varepsilon_t^2|\mathcal{F}_{t-j})-E(\varepsilon_t^2)|^{3/2}=o(j^{-3/2}),\ \textrm{as}\ j\to\infty.
\end{equation}
Then,
\begin{equation}\label{tMSPE.result}
\begin{split}
&\lim\limits_{n\to\infty}n[E\{x_{n+h}-\hat{x}_{n+h}(k)\}^2-E(\varepsilon_{n,h,k}^2)]\\
&= %g_{h}(k) \equiv
\mathrm{tr}\{\mathbf{R}^{-1}(k)\mathbf{L}_{0,h}(k)\}+2\sum_{s=1}^{h-1}\mathrm{tr}\{\mathbf{R}^{-1}(k)\mathbf{L}_{s,h}(k)\},
\end{split}
\end{equation}
where $\mathbf{L}_{s,h}(k)=E\{\mathbf{x}_k(k)\mathbf{x}_{k+s}^\top(k)\varepsilon_{k,h,k}\varepsilon_{k+s,h,k}\}$ and $\sum_{i=a}^{b}\cdot =0$ if $a>b$.
\\
\\
{\rm (b)} Equation \eqref{tMSPE.result} remains valid when
Condition (G) in (a) is replaced by Condition (S).
\end{theorem}

%\begin{remark}
%While the moment condition \eqref{tMSPE.moment} appears restrictive for GARCH-type errors, the asymptotic results established in Theorem \ref{tMSPE} represent a fisrt step toward understanding time series prediction under the joint challenges of conditional heteroscedasticity and model misspecification. The decomposition in \eqref{ptMSPE.1} in Section \ref{appB} of the Supplementary Material indicates that $n[E\{x_{n+h}-\hat{x}_{n+h}(k)\}^2-E(\varepsilon_{n,h,k}^2)]$ contains a term that is intrinsically of sixth-moment order. Consequently, it is unlikely that the required moment condition for $\varepsilon_t$ can be further relaxed in the presence of conditional heteroscedasticity and model misspecification. 


%While the moment condition \eqref{tMSPE.moment} is restrictive for GARCH-type errors, the asymptotic result established in Theorem \ref{tMSPE} is the first step toward the  context of the time series prediction under conditional heteroscedasticity and model misspecification. The decomposition in \eqref{ptMSPE.1} in Section \ref{appB} of the Supplementary Material indicates that $n[E\{x_{n+h}-\hat{x}_{n+h}(k)\}^2-E(\varepsilon_{n,h,k}^2)]$ consists of a term that is intrinsically of sixth moment order, and hence the required moment of $\varepsilon_t$ is hardly to be weakened under conditional heteroscedasticity and model misspecification. The numerical study in Example \ref{ex.S2.1} also support that the existence of $6+\delta$ moment of $\varepsilon_t$ is the minimal requirement for \eqref{tMSPE.result} to hold. In contrast, if the focus is restricted to establishing the moment properties of the LS estimator $\hat{\bm{\beta}}_{n,h}$, for instance, $E\lVert\sqrt{n}(\hat{\bm{\beta}}_{n,h}-\bm{\beta}_h)\rVert$, the required moment of $\varepsilon_t$ can be relaxed to the fourth moment. See, e.g., \cite{Bhansali1991} for an analysis of the moment properties of $\hat{\bm{\beta}}_{n,h}$ in the case of independent errors. Furthermore, if the mean absolute prediction error (MAPE) is considered instead of MSPE, the moment requirements for $\varepsilon_t$ would be further reduced. We leave the investigation of the MAPE for the LS predictor to future work. Last, when $\varepsilon_t$ follows a SV-type model, the $6+\delta$ moment condition is naturally satisfied, as $\varepsilon_t$ in this setting possesses finite moments of all orders.
%\end{remark}

%To reveal the effect of conditional heteroscedasticity on the MSPE, we can make use of the following decomposition: for any $0\leq s\leq h-1$,
%\begin{equation}\label{MSPE decomposition}
%\begin{split}
%    \mathrm{tr}\{\mathbf{R}^{-1}(k)\mathbf{L}_{s,h}(k)\}=   & \mathrm{tr}[\mathbf{R}^{-1}(k)E\{\mathbf{x}_k(k)\mathbf{x}_{k+s}^\top(k)\}]E(\tilde{\varepsilon}_{k,h}\tilde{\varepsilon}_{k+s,h}) \\
%     & +\mathrm{tr}\{\mathbf{R}^{-1}(k)\mathbf{L}^\ast_{s,h}(k)\}+\mathrm{tr}\{\mathbf{R}^{-1}(k)\tilde{\mathbf{L}}_{s,h}(k)\},
%\end{split}
%\end{equation}
%where $\tilde{\varepsilon}_{t,h}=\sum_{j=0}^{h-1}\alpha_j\varepsilon_{t+h-j}$, $\mathbf{L}^\ast_{s,h}(k)=E\{\mathbf{x}_k(k)\mathbf{x}_{k+s}^\top(k)(\tilde{\varepsilon}_{k,h}\tilde{\varepsilon}_{k+s,h}-E(\tilde{\varepsilon}_{k,h}\tilde{\varepsilon}_{k+s,h}))\}
%=\mathrm{cov}\{\mathbf{x}_k(k)\mathbf{x}_{k+s}^\top(k),\tilde{\varepsilon}_{k,h}\tilde{\varepsilon}_{k+s,h}\}$, and $\tilde{\mathbf{L}}_{s,h}(k)=E\{\mathbf{x}_k(k)\mathbf{x}_{k+s}^\top(k)(\varepsilon_{k,h,k}\varepsilon_{k+s,h,k}-\tilde{\varepsilon}_{k,h}\tilde{\varepsilon}_{k+s,h})\}$.
The second-order MSPE of $\hat{x}_{n+h}$ in \eqref{tMSPE.result} can be further decomposed into three terms. That is,
\begin{equation*}
\begin{split}
     & \mathrm{tr}\{\mathbf{R}^{-1}(k)\mathbf{L}_{0,h}(k)\}+2\sum_{s=1}^{h-1}\mathrm{tr}\{\mathbf{R}^{-1}(k)\mathbf{L}_{s,h}(k)\}=\bm{A}_{h, k}+\bm{B}_{h, k}+\bm{C}_{h, k},
\end{split}
\end{equation*}
where
\begin{equation*}
\begin{split}
    %=  \\
   \bm{A}_{h, k}\equiv &
   \left\{
   \mathrm{tr}
   \left(\mathbf{R}^{-1}(k)
E\big\{\mathbf{x}_k(k)\mathbf{x}_{k}^\top(k)
\big\}
\right)E\left(\tilde{\varepsilon}_{k,h}^2\right)\right.+\\ &2\sum_{s=1}^{h-1}\mathrm{tr}[\mathbf{R}^{-1}(k)E\{\mathbf{x}_k(k)\mathbf{x}_{k+s}^\top(k)\}]
E(\tilde{\varepsilon}_{k,h}\tilde{\varepsilon}_{k+s,h})\big\}, \\
     \bm{B}_{h, k}\equiv& \Big(\mathrm{tr}\{\mathbf{R}^{-1}(k)\mathbf{L}^\ast_{0,h}(k)\}+2\sum_{s=1}^{h-1}\mathrm{tr}\{\mathbf{R}^{-1}(k)\mathbf{L}^\ast_{s,h}(k)\}\Big),
\end{split}
\end{equation*}
\begin{equation*}
\begin{split}
     \bm{C}_{h, k}\equiv&\Big(\mathrm{tr}\{\mathbf{R}^{-1}(k)\tilde{\mathbf{L}}_{0,h}(k)\}+2\sum_{s=1}^{h-1}\mathrm{tr}\{\mathbf{R}^{-1}(k)\tilde{\mathbf{L}}_{s,h}(k)\}\Big),
\end{split}
\end{equation*}
with
\begin{align*}
\begin{split}
&\mathbf{L}^\ast_{s,h}(k)=E\{\mathbf{x}_k(k)\mathbf{x}_{k+s}^\top(k)(\tilde{\varepsilon}_{k,h}
\tilde{\varepsilon}_{k+s,h}-E(\tilde{\varepsilon}_{k,h}\tilde{\varepsilon}_{k+s,h}))\},\\
%&=\mathrm{cov}\{\mathbf{x}_k(k)\mathbf{x}_{k+s}^\top(k),\tilde{\varepsilon}_{k,h}\tilde{\varepsilon}_{k+s,h}\},
& \tilde{\mathbf{L}}_{s,h}(k)=
E\{\mathbf{x}_k(k)\mathbf{x}_{k+s}^\top(k)(\varepsilon_{k,h,k}\varepsilon_{k+s,h,k}-\tilde{\varepsilon}_{k,h}\tilde{\varepsilon}_{k+s,h})\},
\end{split}
\end{align*}
and
$\tilde{\varepsilon}_{t,h}=\sum_{j=0}^{h-1}\alpha_j\varepsilon_{t+h-j}$. The first term, $\bm{A}_{h, k}$, is related to the model complexity.
In particular, when $h=1$, $\bm{A}_{h, k}=k E(\varepsilon_{k+1}^2)$
is proportional to the number of parameters of the working AR($k$) model.
The second term,
$\bm{B}_{h, k}$,
%$\mathrm{tr}\{\mathbf{R}^{-1}(k)\mathbf{L}_{s,h}^\ast(k)\}$,
%arises from the correlation between $\mathbf{x}_k(k)\mathbf{x}_{k+s}^\top(k)$ and $\tilde{\varepsilon}_{k,h}\tilde{\varepsilon}_{k+s,h}$, which
is attributed to the conditional heteroscedasticity introduced by $\sigma_t$ in $\varepsilon_t$,
and vanishes when $\sigma_t=c$ is a positive constant (implying that $\{\varepsilon_t\}$ is a sequence of independent random variables).
The presence of the third term, $\bm{C}_{h, k}$, is owing to
$\varepsilon_{t,h,k} \neq \tilde{\varepsilon}_{t, h}$,
which occurs when the AR($k$) model is misspecified.
%Hence \eqref{tMSPE.result} can be rewritten as
%\begin{align}
%\label{ing121}
%\begin{split}
%\lim\limits_{n\to\infty}n[E\{x_{n+h}-%\hat{x}_{n+h}(k)\}^2-%E(\varepsilon_{n,h,k}^2)]= %g_{h}(k),
%\end{split}
%\end{align}

Note that \eqref{tMSPE.result}
has been reported in Theorem 2 of \cite{Ing2003a}
when $\sigma_t= c$ and
\eqref{AR(k) misspecified multi} is correctly specified
(i.e., $\varepsilon_{t,h,k} = \tilde{\varepsilon}_{t, h}$),
and Theorem 2.1 of \cite{Hsu2019} when
$\sigma_t=c$ but \eqref{AR(k) misspecified multi} is misspecified.
Therefore, Theorem \ref{tMSPE} can be viewed as an extension of these results
to misspecified AR models with conditional heteroscedastic errors.
This extension, however,
is far from being trivial
owing to the difficulty in developing
\eqref{t.fisher.inverse moment} under (CH), as shown in Section \ref{sec:2}.
Furthermore, the complexity of the formula on the right-hand side of \eqref{tMSPE.result} increases 
substantially in the presence of conditionally heteroscedastic errors compared to independent errors, even when 
\eqref{AR(k) misspecified multi} is correctly specified.
To illustrate this, assume that
$x_t=\beta_1 x_{t-1}+\varepsilon_{t}$, where
$|\beta_1|<1$, $\varepsilon_{t}=\sigma_t z_t$,
$z_t$ are i.i.d. random variables with zero mean and variance 1, and
$\sigma_t^2=\varphi_0+\varphi_1
\varepsilon_{t-1}^2+
\psi_1 \sigma_{t-1}^2$,
with $\varphi_0>0$, $\varphi_1\geq0$, $\psi_1\geq0$,
and $\varphi_1+\psi_1<1$.
The corresponding $h$-step-head
predictive model is
\eqref{AR(k) misspecified multi},
with $\mathbf{x}_{t}(k)=x_t$,
$\bm{\beta}_h(k)= \beta_1^{h}$, and
$\varepsilon_{t,h,k}=\sum_{i=0}^{h-1}\beta_1^{i}
\varepsilon_{t+h-i}$.
After engaging in detailed calculations,
we can represent the right-hand side of \eqref{tMSPE.result} as
\begin{align}
\label{ing24mar1}
\begin{split}
 &  \sum_{i=0}^{h-1}\beta_1^{2i}\sigma_\varepsilon^2+\{\sum_{i=0}^{h-1}\beta_1^{2i}(\varphi_1+\psi_1)^{h-i-1}\}c_1^\star\sigma_\varepsilon^2+2\sum_{s=1}^{h-1}\beta_1^{2s}(\sum_{i=0}^{h-s-1}\beta_1^{2i})\sigma_\varepsilon^2\\
&+2\sum_{s=1}^{h-1}\beta_1^{2s}\{\sum_{i=0}^{h-s-1}\beta_1^{2i}(\varphi_1+\psi_1)^{h-i-1}\}c_1^\star\sigma_\varepsilon^2,
\end{split}
\end{align}
where $\sigma_\varepsilon^2=E(\varepsilon_t^2)$,
$$c_1^\star=\frac{(1-\beta_1^2)\varphi_1(1-\varphi_1\psi_1-\psi_1^2)(E(z_t^4)-1)}{\{1-\beta_1^2(\varphi_1+\psi_1)\}
\{1-(E(z_t^4)\varphi_1^2+2\varphi_1\psi_1+\psi_1^2)\}},$$
and
$1>E(z_t^4)\varphi_1^2+2\varphi_1\psi_1+\psi_1^2$
is ensured by $E(\varepsilon_t^4)<\infty$.
However, when $\varphi_1=\psi_1=0$,
\eqref{ing24mar1} is substantially simplified 
to 
\begin{eqnarray*}
\sum_{i=0}^{h-1}
\beta_1^{2i}\sigma_\varepsilon^2+2\sum_{s=1}^{h-1}\beta_1^{2s}(\sum_{i=0}^{h-s-1}\beta_1^{2i})\sigma_\varepsilon^2,
\end{eqnarray*}
which has been given by Theorem 2 of \cite{Ing2003a}.













%{ Mention the examples constructed in the previous version on the expression of
%2nd-order MSPE and the differences from the iid case}
%$\mathrm{tr}\{\mathbf{R}^{-1}(k)\mathbf{L}^\ast_{s,h}(k)\}$ and $\mathrm{tr}\{\mathbf{R}^{-1}(k)\tilde{\mathbf{L}}_{s,h}(k)\}$ vanish for each $0\leq s\leq h-1$, and the second order term of the MSPE of $\hat{x}_{n+h}(k)$ given in \eqref{tMSPE.result} reduces to

%simplifies to
%\begin{align}
%\label{MSPE second order first term}
%\begin{split}
 %    & \bm{A}_{h, k}=\mathrm{tr}[\mathbf{R}^{-1}(k)E\{\mathbf{x}_k(k)\mathbf{x}_{k}^\top(k)\}] E(\tilde{\varepsilon}_{k,h}^2)\\
  %   & +2\sum_{s=1}^{h-1}\mathrm{tr}[\mathbf{R}^{-1}(k)E\{\mathbf{x}_k(k)\mathbf{x}_{k+s}^\top(k)\}]
   %   E(\tilde{\varepsilon}_{k,h}\tilde{\varepsilon}_{k+s,h})
%\end{split}
%\end{align}

%, where for random vector $\mathbf{x}$, $\mathrm{cov}(\mathbf{x})=E[\{\mathbf{x}-E(\mathbf{x})\}\{\mathbf{x}-E(\mathbf{x})\}^\top]$.
% $\sigma^2=E(\varepsilon_1^2)$.
%When $h=1$, \eqref{MSPE second order first term} further simplifies to the well-known expression, $k\sigma^2$; see, e.g., \cite{Fuller1981}, \cite{Ing2003} and \cite{Kunitomo1985a}. %See Section \ref{sec:3.2} for more implications of Theorem \ref{tMSPE}.
%When
%the right-hand side of \eqref{tMSPE.result} becomes
%\begin{align}
%\label{ing121}
%\begin{split}
%\bm{A}_{h, k}+\bm{C}_{h, k}=
%\mathrm{tr}\{\mathbf{R}^{-1}(k)\mathbf{L}_{0,h}(k)\}+2\sum_{s=1}^{h-1}\mathrm{tr}\{\mathbf{R}^{-1}(k)\mathbf{L}_{s,h}(k)\},
%\end{split}
%\end{align}
%which can be deduced from,


%\begin{remark}\label{remark Ing 2019}



 %also the result \eqref{tMSPE.result} %(and \eqref{cor subset.1})
%in  their paper. However, in order to prove their theorem,
%by directly assuming %based on a set of high-level assumptions including
%\eqref{t.fisher.inverse moment}.
%However, it is unknown from their as we can see from the proofs of Theorems \ref{tfisher} and \ref{t SV negative bound},
%it
Another notable feature of Theorem \ref{tMSPE} is that its moment condition
\eqref{tMSPE.moment} is substantially weaker than those in the existing results on
the ``same-realization'' MSPE; see, e.g., \cite{Chi2021}, \cite{Fuller1981}, \cite{Ryan2013}, \cite{Hsu2019},
\cite{Ing2003}, and \cite{Kunitomo1985a},
where the existence of at least finite $(8+\delta)$th moments
of $\varepsilon_t$ is required.
In fact,
\eqref{ptMSPE.1} in Section \ref{appB} of the Supplementary Material reveals that
\eqref{tMSPE.moment} can hardly be weakened in the case of same-realization prediction.
Moreover, we illustrate via a numerical example in Section \ref{sec:4.1} %Supplement
that when $\varepsilon_t$ only has a finite $(4+\delta)$th moment,
our asymptotic expression
%$g_{h}(k)$, in \eqref{ing122}
on the right-hand side of
\eqref{tMSPE.result}, denoted by
$g_h(k)$,
is about 16 times as large as
the numerical approximation
of \eqref{ing2}, denoted by
$g_{n,h}(k)$ (see \eqref{ing123}), even when $n$ grows up to 5,000.
In sharp contrast,
$g_{n,h}(k)/g_{h}(k)$ is very close to 1 for all $n \geq 500$
as long as \eqref{tMSPE.moment} is satisfied.
%Moreover, as the discussion in Section 2.2 of \cite{Hsu2019}, their assumptions essentially require the boundedness of the $10+\delta$ moment of $\varepsilon_t$, which is much more restrictive than \eqref{tMSPE.moment} used in Theorem \ref{tMSPE}.
%\end{remark}
%\begin{remark}\label{remark betahat MSE and bias}
%As a consequence of the proof of Theorem \ref{tMSPE}, one can obtain moment bounds of the LS estimator $\hat{\bm{\beta}}_{n,h}(k)$. That is, for any $q>0$,
%\begin{equation}\label{beta moment}
%E\lVert\hat{\bm{\beta}}_{n,h}(k)-\bm{\beta}_h(k)\rVert^q=O(n^{-\frac{q}{2}}),
%\end{equation}
%provided \eqref{AR mean part}--\eqref{stationary char poly}, (CH)(i), \eqref{zt density}--\eqref{zt density 2} (or (CH)(ii) and \eqref{smoothness condition zt}), and
%\begin{equation}\label{beta MSE moment}
%\sup\limits_{-\infty<t<\infty}E|\varepsilon_t|^{2q+\delta}<\infty,
%\end{equation}
%for some $\delta>0$. When $q=2$, \eqref{beta moment} reduces to the MSE of $\hat{\bm{\beta}}_{n}(k)$. Moreover, under the assumptions required for \eqref{beta moment} with $q=2$ in \eqref{beta MSE moment}, one can derive the asymptotic bias of $\hat{\bm{\beta}}_{n}(k)$,
%\begin{equation}\label{beta moment 2}
%E\{\hat{\bm{\beta}}_{n,h}(k)\}=\bm{\beta}_h(k)+O(n^{-1}).
%\end{equation}
%It is worth mentioning that while these moment properties of the LS estimator $\hat{\bm{\beta}}_{n,h}(k)$ were established decades ago for the case of independent errors (see, e.g., \cite{Bhansali1991} and \cite{Kunitomo1985a}), the above results seem to be the first extensions to the case of conditional heteroscedastic errors.
%\end{remark}

If the focus
is instead on
the mean squared error of
$\hat{\bm{\beta}}_{n,h}$, namely
%(e.g., the MSE of the LS estimator),
%the moment properties of the LS estimator $\hat{\bm{\beta}}_{n,h}$, for instance,
$E\lVert 
\hat{\bm{\beta}}_{n,h}-\bm{\beta}_h\rVert^2$, the required moment
condition on $\varepsilon_t$ can indeed be reduced to $4+\delta$.
%For an analysis of the moment properties of $\hat{\bm{\beta}}_{n,h}$ in the case of independent errors,see \cite{Bhansali1991}. 
Moreover, if prediction
accuracy is measured by the mean absolute prediction error (MAPE) rather than the MSPE, the
moment condition on $\varepsilon_t$
could potentially be further weakened. A detailed investigation of MAPE-based prediction, however, is beyond the scope of the
present paper and is left for future work.
The remainder of this section is dedicated to the discussion of 
\eqref{tMSPE.moment2} and \eqref{tMSPE.conditional moment}.
%in the rest of this section. %, the technical assumptions imposed on the innovations in Theorem \ref{tMSPE},

%can be
%are fulfilled by a wide range of conditional heteroscedastic processes, suggesting the broad applicability of result \eqref{tMSPE.result}.
\begin{proposition}\label{prop3.1}
Suppose that $\{\varepsilon_t^2\}$ admits a martingale representation,
\begin{equation}\label{h.varepsilon square}
\varepsilon_t^2-E(\varepsilon_t^2)=\sum_{s=0}^{\infty}\mathbf{a}^\top_s\bm{w}_{t,s},
\end{equation}
where $\mathbf{a}_s$ are $r$-dimensional, $r\geq 1$, real vectors satisfying
\begin{equation}\label{h.coefficient a}
\lVert \mathbf{a}_s\rVert=O\{(s+1)^{-\varsigma}\},\,\,\mbox{for some}\,\,\varsigma>3/2,
\end{equation}
and for any $t$,
$\bm{w}_{t,t-j}$ is
$\mathcal{F}_{j}$-measurable
and obeys
\begin{equation}\label{wt moment0}
E(\bm{w}_{t,t-j}|\mathcal{F}_{j-1})=\mathbf{0}\,\,\mbox{a.s.}
\end{equation}
Moreover, assume for some $\delta^{*}>0$, 
\begin{equation}\label{wt moment}
\sup\limits_{-\infty<t,s<\infty,}E\lVert \bm{w}_{t,s}\rVert^{3+\delta^\star}<\infty,
\end{equation}
where
$\bm{w}_{t,s}=\mathbf{0}$ if $s<0$.
Then, \eqref{tMSPE.moment2} and \eqref{tMSPE.conditional moment} hold. %with $\delta^\star=\delta$.
\end{proposition}

The proof of Proposition \ref{prop3.1} is given in
Section \ref{appB} of the Supplementary Material. In the following examples,
we argue that
\eqref{h.varepsilon square}--\eqref{wt moment}
are easily met in common practice.

\begin{example}\label{remark3.1}
Consider again the asymmetric power GARCH model, which is
\eqref{varepsilon basic form} with $\sigma_t$
satisfying \eqref{apGARCH}. %with $z_t$ being symmetric.
When $\mu=2$ in \eqref{apGARCH}, the model can be expressed as,
\begin{equation}\label{ap GARCH alter 2}
\varepsilon_t^2=\varphi_{0}+\sum_{i=1}^{p^\prime\vee q^\prime}\{\varphi_{i}E(|z_1|-\lambda_i z_1)^2+\psi_{i}\}\varepsilon_{t-i}^2+w_{1,t}-\sum_{j=1}^{q^\prime}\psi_{j}w_{1,t-j}
+\sum_{l=1}^{p^\prime}4\varphi_{l}\lambda_lw_{2,t-l},
\end{equation}
where $w_{1,t}=\varepsilon_{t}^2-\sigma_{t}^2$, $w_{2,t}=\varepsilon_{t}^2I_{\{\varepsilon_t<0\}}-\frac{1}{2}\varepsilon_{t}^2$, and $\varphi_{i}$, $\lambda_i$, and $\psi_{j}$, are set to 0 when $i>p^\prime$ and $j>q^\prime$.
 As discussed in \cite{Ding1993},
\eqref{ap GARCH alter 2} includes the GARCH model (\cite{Bollerslev1986}) as well as the GJR-GARCH model (\cite{Glosten1993})
as special cases.
Let $\mathcal{F}_{t}=\sigma(z_{t},z_{t-1},\ldots)$
and assume that $z_t$ are symmetric.
Then, 
\eqref{h.varepsilon square}--\eqref{wt moment0} hold with $\bm{w}_{t,s}=\bm{w}_{t-s}=(w_{1,t-s},w_{2,t-s})^\top$ and $\mathbf{a}_s=(a_{1s}, a_{2s})^\top$,
where $a_{1s}$ and $a_{2s}$, respectively, satisfy
\begin{equation*}\label{ap GARCH alter ceof 1}
\sum_{s=0}^{\infty}a_{1s}z^{s}=\frac{1-\sum_{i=1}^{q^\prime}\psi_{i}z^{i}}
{1-\sum_{j=1}^{p^\prime\vee q^\prime}\{\varphi_{j}E(|z_1|-\lambda_j z_1)^2+\psi_{j}\}z^{j}}
\end{equation*}
and
\begin{equation*}\label{ap GARCH alter ceof 2}
\sum_{s=0}^{\infty}a_{2s}z^{s}=\frac{\sum_{i=1}^{p^\prime}4\varphi_{i}\lambda_iz^{i}}
{1-\sum_{j=1}^{p^\prime\vee q^\prime}\{\varphi_{j}E(|z_1|-\lambda_j z_1)^2+\psi_{j}\}z^{j}}.
\end{equation*}
See Appendix B of \cite{Ding1993} and Theorem 3.1 of \cite{Ling2002} for more details.
Moreover, if for some $\delta_1>0$,
\begin{equation}\label{varepsilon 1 moment}
E|\varepsilon_1|^{6+\delta_1}<\infty,
\end{equation}
then \eqref{wt moment} follows from \eqref{varepsilon 1 moment} and the i.i.d. assumption on $\{z_t\}$.
%Consequently, \eqref{tMSPE.moment2} and \eqref{tMSPE.conditional moment} hold,
%in view of Proposition \ref{prop3.1}.
%Note that %\eqref{apGARCH stationary} with $\mu=2$, which is ensured by
%\eqref{varepsilon 1 moment} and \eqref{ap GARCH alter 2} together imply that
%\eqref{h.varepsilon square} holds with $\lVert \mathbf{a}_s\rVert\leq \tilde{c}_1\exp(-\tilde{c}_2s)$, for some $0<\tilde{c}_1<\tilde{c}_2<\infty$
%(see ).
%See also Section 3 of \cite{Huang2022} for a discussion of the relation between \eqref{h.varepsilon square} and \eqref{ap GARCH alter 2} when $\{\varepsilon_t\}$ follows a GJR-GARCH process.
\end{example}

\begin{example}\label{remark3.2}
Suppose that $\{\varepsilon_t\}$ follows the SV model 
\eqref{varepsilon basic form} and
\eqref{SV model}, and obeys $E|z_1|^{6+\delta_1}<\infty$, for some $\delta_1>0$. Let $\mathcal{F}_{t}=\sigma(z_t,z_{t-1},\ldots,v_{t+1},v_t,\ldots)$. Then,
it can be shown that \eqref{h.varepsilon square}--\eqref{wt moment} hold with $\bm{w}_{t,0}=\varepsilon_t^2-\sigma_t^2$,
\begin{equation}\label{SV wts}
\bm{w}_{t,s}=
\begin{cases}
  \tilde{b}_s^{-1}\{e^{\tilde{b}_sv_{t-s+1}}-E(e^{\tilde{b}_sv_{t-s+1}})\}\prod_{u=s+1}^{\infty}e^{\tilde{b}_uv_{t-u+1}}, & \mbox{if } \tilde{b}_s\neq0, \\
  v_{t-s+1}\prod_{u=s+1}^{\infty}e^{\tilde{b}_uv_{t-u+1}}, & \mbox{if } \tilde{b}_s=0,
\end{cases}
\end{equation}
for $s\geq1$, $\mathbf{a}_0=1$, $\mathbf{a}_1=\tilde{b}_1e^{\tilde{b}_0^\ast}$, and
\begin{equation}\label{SV as}
\mathbf{a}_s=\tilde{b}_se^{\tilde{b}_0^\ast}\prod_{i=1}^{s-1}E(e^{\tilde{b}_iv_{t-i+1}}),
\end{equation}
for $s\geq2$, where $\tilde{b}^\ast_0=\tilde{a}_0/(1-\sum_{i=1}^{\tilde{p}}\tilde{a}_i)$ and $\sum_{i=0}^{\infty}\tilde{b}_{i+1}z^i=(1-\sum_{i=1}^{\tilde{p}}\tilde{a}_iz^i)^{-1}$. Indeed, $\lVert \mathbf{a}_s\rVert\leq \tilde{c}_1\exp(-\tilde{c}_2s)$, for some $0<\tilde{c}_1<\tilde{c}_2<\infty$, $\sup_{-\infty<t<\infty}E\lVert \bm{w}_{t,0}\rVert^{3+\delta}<\infty$, for some $\delta>0$, and $\sup_{-\infty<t<\infty,s\geq1}E\lVert \bm{w}_{t,s}\rVert^{\tilde{\theta}}<\infty$, for any $\tilde{\theta}>0$.
%Therefore, \eqref{tMSPE.moment2} and \eqref{tMSPE.conditional moment} follows.
\end{example}


Before concluding this section, we note that
\eqref{tMSPE.moment2} and \eqref{tMSPE.conditional moment}
are also satisfied by other types of weakly dependent processes, such as those whose serial dependencies are characterized using functional dependence measures
(\cite{Wu2005Main}) or strong mixing conditions. For additional details, refer to Section \ref{app:A}.

\subsection{Best Subset Selection}\label{sec:3.2}
A natural generalization of \eqref{AR(k) misspecified multi} is the subset AR model for $h$-step prediction,
\begin{equation}\label{xt+h J}
x_{t+h}=\bm{\beta}_h^\top(J)\mathbf{x}_{t}(J)+\varepsilon_{t,h,J},
\end{equation}
where $J\subset\mathbb{N}$, $\mathbf{x}_{t}(J)=(x_{t+1-j},j\in J)^\top$, and
$\bm{\beta}_h(J)=\arg\min_{\mathbf{c}\in\mathbb{R}^{\sharp(J)}}E\{x_{t+h}-\mathbf{c}^\top\mathbf{x}_{t}(J)\}^2$. 
%Note that if a specific number $m$ is in the set $J$, it means that the $m$-th lagged variable is included in the predictive model \eqref{xt+h J}. 
With model \eqref{xt+h J}, the $h$-step LS predictor
for $x_{n+h}$ is
\begin{equation}\label{hat xn+h J}
\hat{x}_{n+h}(J)=\hat{\bm{\beta}}_{n,h}^\top(J)\mathbf{x}_n(J),
\end{equation}
where
\begin{equation*}%\label{betahat J}
\hat{\bm{\beta}}_{n,h}(J)=\hat{\mathbf{R}}^{-1}_{n,h}(J)\frac{1}{n-h-\tilde{d}+1}\sum_{j=\tilde{d}}^{n-h}\mathbf{x}_j(J)x_{j+h},
\end{equation*}
with
\begin{equation*}%\label{Rhat J}
\hat{\mathbf{R}}_{n,h}(J)=\frac{1}{n-h-\tilde{d}+1}\sum_{j=\tilde{d}}^{n-h}\mathbf{x}_j(J)\mathbf{x}^\top_j(J),
\end{equation*}
and  $\tilde{d} \equiv$ the largest number in $J$.
%Recall $\bm{A}_{h, k}, \bm{B}_{h, k}$, and $\bm{C}_{h, k}$ in \eqref{tMSPE.result}.
The following
corollary shows that
the limit of 
$
n\{E\{x_{n+h}-\hat{x}_{n+h}(J)\}^2-E(\varepsilon_{n,h,J}^2)\}$
exists and has an expression 
akin to the right-hand side of
\eqref{tMSPE.result}.


%The following corollary
%is a counterpart of \eqref{ing122} in model \eqref{xt+h J},
%whose proof is similar to that of Theorem \ref{tMSPE}
%and is omitted here.
 %can be obtained in the setting of \eqref{xt+h J} and \eqref{hat xn+h J}.

%It is desirable to know whether a conclusion analogous to \eqref{tMSPE.result} can be obtained. The answer is confirmed by the next corollary.

\begin{corollary}\label{cor subset}
Under the same assumptions as in Theorem \ref{tMSPE}(a) (GARCH-type errors) or Theorem \ref{tMSPE}(b) (SV-type errors),
the limit in \eqref{ing4} exists and has the following expression:
\begin{align}\label{cor subset.1}
\begin{split}
g_{h}(J)=\mathrm{tr}\{\mathbf{R}^{-1}(J)\mathbf{L}_{0,h}(J)\}+2\sum_{s=1}^{h-1}\mathrm{tr}\{\mathbf{R}^{-1}(J)\mathbf{L}_{s,h}(J)\},
\end{split}
\end{align}
where
$\mathbf{R}(J)=E\{\hat{\mathbf{R}}_{n,h}(J)\}$ and $\mathbf{L}_{s,h}(J)=E\{\mathbf{x}_{1}(J)\mathbf{x}_{1+s}^\top(J)\varepsilon_{1,h,J}\varepsilon_{1+s,h,J}\}$.
\end{corollary}

The derivation of Corollary \ref{cor subset} relies on Theorem \ref{tfisher} and follows the approach used in the proof of Theorem \ref{tMSPE}, which is omitted here for conciseness.
In fact, in the special case where
$\varepsilon_t$ are independent errors,
\eqref{cor subset.1} has been established
in Theorem 3.2 of \cite{Hsu2019}.
These authors further suggested estimating
$g_{h}(J)$ using
\begin{eqnarray}
\label{ing125}
\hat{g}_h(J)=\mathrm{tr}\{\hat{\mathbf{R}}_{n,h}^{-1}(J)\hat{\mathbf{L}}_{0,h}(J)\}+2\sum_{s=1}^{h-1}
\mathrm{tr}\{\hat{\mathbf{R}}_{n,h}^{-1}(J)\hat{\mathbf{L}}_{s,h}(J)\},
\end{eqnarray}
where
\begin{eqnarray*}
\hat{\mathbf{L}}_{s,h}(J)=(n-h-\tilde{d}-s+1)^{-1}
\sum_{t=\tilde{d}}^{n-h-s}\mathbf{x}_t(J)\mathbf{x}_{t+s}^\top(J)\hat{\varepsilon}_{t,h,J}\hat{\varepsilon}_{t+s,h,J},
\end{eqnarray*}
with $\hat{\varepsilon}_{t,h,J}=x_{t+h}-\hat{x}_{t+h}(J)$ and $\hat{x}_{t+h}(J)=\hat{\bm{\beta}}_{n,h}^\top(J)\mathbf{x}_t(J)
$.
They also demonstrated that
\begin{eqnarray}
\label{ing129}
\hat{g}_{h}(J)=g_h(J) +o_p(1),
\end{eqnarray}
and
\begin{eqnarray}
\label{ing128}
\hat{\sigma}^2_{h}(J)=f_h(J) +O_p(n^{-1/2}),
\end{eqnarray}
where
\begin{eqnarray}
\label{ing124}
\hat{\sigma}^2_{h}(J)=\frac{1}{n-h-\tilde{d}+1}\sum_{t=\tilde{d}}^{n-h}\hat{\varepsilon}_{t,h,J}^2.
\end{eqnarray}
Equation \eqref{ing128} asserts that
$\hat{\sigma}^2_{h}(J)$ is a $\sqrt{n}$-consistent estimate of
$f_h(J)=E(\varepsilon_{t,h,J}^2)$.





% can be proved in the same manner as that in the proof of Theorem \ref{tMSPE}; details are omitted.
%After obtaining %establishing
%an asymptotic expression for the MSPE of the LS predictor in Corollary \ref{cor subset}, we are led to the question of how to select the best predictive model from a finite family of candidate models in the presence of conditional heteroscedasticity. More specifically,
 
Consider a finite set ${\cal J}$ of candidate subset AR models,
\begin{equation*}
x_{t+h}=\bm{\beta}_h^\top(J)\mathbf{x}_{t}(J)+\varepsilon_{t,h,J},\ J\in\mathcal{J},
\end{equation*}
where the number of candidate models is finite,
and the cardinality $\sharp{(J)}$
of each $J \in {\cal J}$ is uniformly bounded above by a finite constant.
Based on
\eqref{ing129} and \eqref{ing128},
\cite{Hsu2019} proposed selecting the best predictive model
using MRIC, defined in \eqref{MRIC0}, with
$\hat{g}_{h}(J)$
and $\hat{\sigma}_{h}(J)$
given by
\eqref{ing125}
and \eqref{ing124}, respectively.
The penalty
$C_n$ satisfies
\begin{eqnarray}
\label{ing126}
\frac{C_n}{n^{1/2}}\to\infty \,\,\mbox{and}\,\, \frac{C_n}{n} \to 0.
\end{eqnarray}
%and
%$\hat{g}_h(J)$ and
%$\hat{\sigma}^2_{h}(J)$
%in  \eqref{MRIC0}
%are
%\eqref{ing125} and \eqref{ing124}, respectively,
%with $\tilde{d}$ 
%replaced by $\bar{d} \equiv$
%the largest number in $\bigcup_{J\in\mathcal{J}}J$.
%Define $f_h(J_l)=E(\varepsilon_{t,h,J_l}^2)$. 
%Let
%\begin{eqnarray*}
%M_1(h)=\{i:1\leq i\leq %K,f_h(J_i)=\min\limits_{1\leq l\leq K}f_h(J_l)\}.
%\end{eqnarray*}
%Then,
%\begin{eqnarray}
%\label{M1}
%{\cal C}_{M_1(h)}\equiv\{J_l: l \in M_1(h)\}
%\end{eqnarray}
%is the set of candidate models having the %smallest population MSPE.
%Let
%\begin{eqnarray*}
%M_2(h)=\{i:i\in M_1(h),g_h(J_i)=\min\limits_{l\in %M_1(h)}g_h(J_l)\}.
%\end{eqnarray*}
%Our goal is to choose a model belonging to %${\cal C}_{M_2(h)}\equiv\{J_l: l \in M_2(h)\}$, where
%\begin{eqnarray}
%\label{ing3}
%{\cal C}_{M_2(h)}\equiv\{J_l: l \in M_2(h)\},
%\end{eqnarray}
%which
%consists of the candidates having the best (in terms of the second-order MSPE)
%predictive power
%among ${\cal C}_{M_1(h)}$.
%We are interested in selecting $1 \leq \hat{l} \leq K$ such that
%$\hat{x}_{n+h}(J_{\hat{l}})$ has the smallest MSPE when $x_t$ is driven by conditional heteroscedastic errors. %See Remark \ref{remark subset} for the definitions of $\bm{\beta}_h(J_l)$, $\mathbf{x}_{t}(J_l)$, $\hat{x}_{n+h}(J_l)$ and other notations used below.
%Recently, \cite{Hsu2019} proposed using
%MRIC to (asymptotically) achieve this goal under the assumption of independent noise.
%\begin{eqnarray}
%\label{ing124}
%\hat{\sigma}^2_{h}(J_l)=\frac{1}{n-h-\bar{d}+1}\sum_{t=\bar{d}}^{n-h}\hat{\varepsilon}_{t,h,J_l}^2
%\end{eqnarray}
%is a natural estimate of $f_{h}(J_l)$,
%$\check{g}_h(J_l)$ is a consistent estimate of $g_h(J_l)$, and
% $\{C_n\}$ is a sequence of positive numbers satisfying
%In the case of independent noise, it is argued in \cite{Hsu2019}
%that $g_{h}(J)=g_{h}(J)$ is also valid and
%\begin{eqnarray}
%\label{ing125}
%\hat{g}_h(J_l)\equiv\mathrm{tr}\%{\hat{\mathbf{R}}_{n,h}^{-1}(J_l)\hat{\mathbf{L}}_{0,h}(J_l)\}+2\sum_{s=1}^{h-1}
%\mathrm{tr}\{\hat{\mathbf{R}}_{n,h}^{-1}(J_l)\hat{\mathbf{L}}_{s,h}(J_l)\}
%\end{eqnarray}
%is a consistent estimate of $g_{h}(J_l)$,
%where
%\begin{eqnarray*}
%\hat{\mathbf{L}}_{s,h}(J_l)=(n-h-\bar{d}-s+1)^{-1}
%\sum_{t=\bar{d}}^{n-h-s}\mathbf{x}_t(J_l)\mathbf{x}_{t+s}^\top(J_l)\hat{\varepsilon}_{t,h,J_l}\hat{\varepsilon}_{t+s,h,J_l}.
%\end{eqnarray*}
%noting that $g_{h}(J_l)$ in the case of independent errors
%has the same expression as the one obtained under assumption (CH) (see \eqref{cor subset.1}).
Define
\begin{eqnarray*}
\label{ing127}
\hat{J}(h)=\arg\min_{J\in\mathcal{J}}\textrm{MRIC}_{h}(J).
\end{eqnarray*}
% defined in \eqref{ing125}.
\cite{Hsu2019}
demonstrated that
\eqref{ing3} holds with
$\hat{J}=\hat{J}(h)$,
even when all $J$ are misspecified and ${\cal J}_{1}(h)$, defined in \eqref{M1}, contains multiple elements.

It is worth noting that this nuanced model selection problem was also examined by
\cite{Kilian2006}, who showed that BIC-like criteria consistently select the model with the fewest parameters among those
in ${\cal J}_{1}(h)$.
However,
as pointed out by
\cite{Findley91} and
\cite{Hsu2019}, when both $J_1$ and $J_2$ are misspecified, 
the ranking of their second-order MSPE terms
$g_h(J_1)$ and $g_h(J_2)$ may not
correspond to the comparison of their cardinalities. 
This raises the possibility that
the model with the fewest parameters in ${\cal J}_{1}(h)$ may not belong to 
${\cal J}_{2}(h)$, 
as defined in \eqref{M2}. 
Essentially, achieving \eqref{ing3} via conventional criteria such as AIC, BIC, or HQ is problematic, since these criteria penalize only model size and fail to account for second-order MSPE. Although \cite{Hsu2019} only
established the consistency of MRIC under the assumption of independent errors, 
Theorem \ref{t MRIC} below
shows that MRIC remains asymptotically valid as long as the expression for
$g_{h}(J)$ in
\eqref{cor subset.1} holds.
 This includes a wide range of conditionally heteroscedastic processes, provided suitable moment conditions on
$\{\varepsilon_t\}$ are met.











%The result \eqref{cor subset.1} indicates that the MRIC proposed by \cite{Hsu2019} can be an answer. The MRIC is of the form,
%and

 %and $g_h(J_l)$ is defined in \eqref{ing4}.
%The following theorem shows that MRIC still achieves asymptotically efficiency in the presence of conditional heteroscedasticity.

\begin{theorem}\label{t MRIC}
Suppose that \eqref{AR mean part}--\eqref{stationary char poly} hold, and assume
\begin{equation}\label{tMRIC.moment}
\sup\limits_{-\infty<t<\infty}E|\varepsilon_t|^{4}<\infty.
\end{equation}
Moreover,
suppose that the expression for
$g_h(J)$ in \eqref{cor subset.1} holds,
and the following conditions are satisfied:
\begin{equation}\label{pt.MRIC.1.3}
\sup\limits_{-\infty<s_1,s_2<\infty}E|\frac{1}{\sqrt{n}}\sum_{t=1}^{n}(\varepsilon_{t+s_1}\varepsilon_{t+s_2}-E(\varepsilon_{t+s_1}\varepsilon_{t+s_2}))|^2= O(1),
\end{equation}
and 
\begin{equation}\label{pt.MRIC.2.6}
\begin{split}
&\sup\limits_{-\infty<s_1,s_2,s_3,s_4<\infty}E|\frac{1}{n}\sum_{t=1}^{n}(\varepsilon_{t+s_1}\varepsilon_{t+s_2}\varepsilon_{t+s_3}\varepsilon_{t+s_4}-E(\varepsilon_{t+s_1}\varepsilon_{t+s_2}\varepsilon_{t+s_3}\varepsilon_{t+s_4}))|\\
=&o(1).
\end{split}
\end{equation}
Then, the model,$\hat{J}(h)$,
defined as the minimizer of the MRIC criterion with
$\hat{g}_{h}(J)$,
$\hat{\sigma}^{2}_{h}(J)$, and $C_n$
given by
\eqref{ing125}, \eqref{ing124}, and
\eqref{ing126}, respectively,
satisfies
\eqref{ing3}. 
\end{theorem}

Assumptions \eqref{pt.MRIC.1.3} and \eqref{pt.MRIC.2.6}, similar to \eqref{tMSPE.moment2} and \eqref{tMSPE.conditional moment},
can be verified for
a wide range of weakly dependent processes
that exhibit certain functional dependence structures, satisfy mixing conditions, or meet condition \eqref{h.varepsilon square}.
In Section \ref{appA00}, we provide a detailed discussion on these sufficient conditions and illustrate how GARCH and SV processes fulfill \eqref{pt.MRIC.1.3} and \eqref{pt.MRIC.2.6}.
%Finally, we highlight that pursuing 
%\eqref{ing3} using MRIC
%may be of limited significance without the support of
%Corollary \ref{cor subset}, which establishes (under suitable conditions) that
%the second-order MSPE of candidate model $J$
%for predicting $x_{n+h}$ converges to the limiting value
%given on the right-hand side of
%\eqref{cor subset.1}.

\subsection{Further Discussion on \eqref{tMSPE.moment2} and \eqref{tMSPE.conditional moment}}\label{app:A}
In this section, we illustrate that \eqref{tMSPE.moment2} and \eqref{tMSPE.conditional moment} also hold for the following two classes of weakly dependent processes.

{\bf A.} {\it 
Strictly stationary processes with dependencies characterized by functional dependence measures}.
Assume $\{\varepsilon_t\}$ is a strictly stationary process of the form
\begin{equation}\label{propA.1.1}
\varepsilon_t=\tilde{f}(\ldots,\tilde{\mathbf{z}}_{t-1},\tilde{\mathbf{z}}_{t}),
\end{equation}
where $\tilde{\mathbf{z}}_t$ are i.i.d. random variables, and $\tilde{f}$ is a measurable function such that $\varepsilon_t$ is well-defined. Let $\{\mathbf{z}^\star_t\}$ be an i.i.d. copy of $\{\tilde{\mathbf{z}}_t\}$, and define
$\varepsilon_i^\star=\tilde{f}(\ldots,\tilde{\mathbf{z}}_{-1},\mathbf{z}^\star_0,\tilde{\mathbf{z}}_1,\ldots,\tilde{\mathbf{z}}_i)$.
The functional dependence measure of
$\{\varepsilon_t\}$ is then given by
$\delta^\prime_\xi(i)=(E|\varepsilon_i-\varepsilon_i^\star|^\xi)^{1/\xi}$, $\xi>0$.
For more details, see \cite{Wu2005Main}.




\begin{proposition}\label{propA.1}
%Let $\delta$ be defined as in Theorem \ref{tMSPE}.
Assume \eqref{propA.1.1},
%$\{\varepsilon_t\}$ is a strictly stationary process of the form
%\begin{equation}\label{propA.1.1}
%\varepsilon_t=\tilde{g}(\ldots,\tilde{\mathbf{z}}_{t-1},\tilde{\mathbf{z}}_{t}),
%\end{equation}
%where $\tilde{\mathbf{z}}_t$ are i.i.d. random variables, and $\tilde{g}$ is a measurable function such that $\varepsilon_t$ is well-defined. Let $\{\mathbf{z}^\star_t\}$ be an i.i.d. copy of $\{\tilde{\mathbf{z}}_t\}$. 
\begin{equation}\label{propA.1.3}
E|\varepsilon_1|^{6+{\delta}}<\infty,
\end{equation}
and
\begin{equation}\label{propA.1.2}
\delta^\prime_{6+{\delta}}(n)=O(n^{-\nu}),
\end{equation}
for some ${\delta}>0$ and $\nu>2$. Then, \eqref{tMSPE.moment2},
with $\delta^\star={\delta}/2$,
and \eqref{tMSPE.conditional moment} follow. 
\end{proposition}

{\bf B.} {\it Strong mixing processes}.
Define the strong mixing coefficient
$$\alpha(m)=\sup\limits_{-\infty<t<\infty}\sup\limits_{A\in\mathcal{F}_{-\infty}^t,B\in\mathcal{F}_{t+m}^\infty}|P(A\cap B)-P(A)P(B)|,$$
where $\mathcal{F}_j^l=\sigma(\varepsilon_t,j\leq t\leq l)$. We say that $\{\varepsilon_t\}$ is strong mixing if $\alpha(m)\to0$ as $m\to\infty$.

\begin{proposition}\label{propA.2}
%Let $\delta$ be defined as in Theorem \ref{tMSPE}.
Assume that $\{\varepsilon_t\}$ is a strong mixing process satisfying \eqref{tMSPE.moment} 
%\begin{equation}\label{propA.2.2}
%\sup\limits_{-\infty<t<\infty}E|\varepsilon_t|^{6+\tilde{\delta}}<\infty,
%\end{equation}
and 
\begin{equation}\label{propA.2.3}
\alpha(n)=O(n^{-{\nu}}),
\end{equation}
for some $\nu\geq3$. Then, \eqref{tMSPE.moment2},
with $\delta^\star={\delta}/2$, and \eqref{tMSPE.conditional moment} hold.
\end{proposition}

The proofs of Propositions \ref{propA.1} and \ref{propA.2} are provided in Section \ref{supp3.1} of the Supplementary Material. %Similar to the coefficient assumptions
%in Proposition \ref{prop3.1}, assumptions \eqref{propA.1.2} and \eqref{propA.2.1} ensure weak dependence in $\{\varepsilon_t\}$. Additionally, assumptions \eqref{propA.1.4} and \eqref{propA.2.3} 
%guarantee the rate specified
%in \eqref{tMSPE.conditional moment}. 
Together with Proposition \ref{prop3.1}, these results establish that
a broad class of weakly dependent processes satisfies
assumptions \eqref{tMSPE.moment2} and \eqref{tMSPE.conditional moment}.


\subsection{Further Discussion on \eqref{pt.MRIC.1.3} and \eqref{pt.MRIC.2.6}}\label{appA00}
In this section, we demonstrate that assumptions \eqref{pt.MRIC.1.3} and \eqref{pt.MRIC.2.6} 
are satisfied by the processes presented in Section
\ref{app:A}.

%However, as the corresponding technical arguments are substantially more involved, we do not pursue them in this paper.

%Specifically, by assuming $\varepsilon_{t_1}^2\varepsilon_{t_2}^2 - E(\varepsilon_{t_1}^2\varepsilon_{t_2}^2)$ exhibits a martingale structure similar but more complex than \eqref{h.varepsilon square}, one can show that \eqref{pt.MRIC.1.3} and \eqref{pt.MRIC.2.6} hold and the assumed martingale structure accommodate GARCH and SV processes. However, considering the corresponding proof is complex and lengthy, we do not pursue it here. Instead, we discuss processes functional dependence condition or mixing condition

%This makes the corresponding proof complex and lengthy. The details are not included here but can be found in an extended version of this paper, 
%\cite{Huang2025}.

\begin{proposition}\label{propB.1}
Assume \eqref{propA.1.1},
\begin{equation}\label{propB.1.1}
E|\varepsilon_1|^{4+\delta}<\infty,
\end{equation}
and
\begin{equation}\label{propB.1.2}
\Theta_{4+\delta}<\infty,
\end{equation}
for some ${\delta}>0$, where $\Theta_{\xi}=
\sum_{i=0}^{\infty}\delta^\prime_{\xi}(i)$. Then, \eqref{pt.MRIC.1.3} and \eqref{pt.MRIC.2.6} are valid.
\end{proposition}

\begin{proposition}\label{propB.2}
Assume $\{\varepsilon_t\}$ is a strong mixing process with 
\begin{equation}\label{propB.2.1}
\sup\limits_{-\infty<t<\infty}E|\varepsilon_t|^{4+{\delta}}<\infty,
\end{equation}
for some ${\delta}>0$. Moreover, suppose that
\begin{equation}\label{App.2}
\sum_{i=1}^{\infty}\alpha(i)^{\frac{{\delta}}{4+{\delta}}}<\infty.
\end{equation}
Then, \eqref{pt.MRIC.1.3} and \eqref{pt.MRIC.2.6} hold.
\end{proposition}

The proofs of Propositions \ref{propB.1} and \ref{propB.2} are given in Section \ref{Ex3-S} of the Supplementary Material.

\begin{remark}\label{rm.mixing}
%When ${\varepsilon_t}$ is uniform mixing, the condition on the corresponding mixing coefficients can be relaxed to satisfy \eqref{pt.MRIC.1.3} and \eqref{pt.MRIC.2.6}. To see this, 
Define the uniform mixing coefficient as
$$\phi(m)=\sup\limits_{-\infty<t<\infty}\sup\limits_{A\in\mathcal{F}_{-\infty}^t,B\in\mathcal{F}_{t+m}^\infty,P(A)>0}|P(B|A)-P(B)|.$$
The process $\{\varepsilon_t\}$ is said to be uniform mixing if $\phi(m) \to 0$ as $m \to \infty$. If $\{\varepsilon_t\}$ is uniform mixing and satisfies \eqref{propB.2.1} along with
\begin{equation}\label{r.mixing.1} 
\sum_{i=1}^{\infty} \phi(i)^{\frac{2+{\delta}}{4+{\delta}}} < \infty, 
\end{equation}
then \eqref{pt.MRIC.1.3} and \eqref{pt.MRIC.2.6} hold (see \cite{Andrews1988M} and \cite{Kim1994M}). This implies that when $\{\varepsilon_t\}$ satisfies uniform mixing instead of strong mixing, the required decay rate of the mixing coefficients can be relaxed from \eqref{App.2} to \eqref{r.mixing.1}.
\end{remark}

%Proofs of Proposition \ref{propB.1}--\ref{propB.3} are provided in Section \ref{Ex3-S} of the Supplementary Material. To satisfy \eqref{pt.MRIC.2.6}, Proposition \ref{propB.1} only requires the existence of the fourth moment of $\varepsilon_t$, but assumes that ${\varepsilon_t}$ is strictly stationary and driven by i.i.d. errors. In contrast, Propositions \ref{propB.2} and \ref{propB.3} do not assume that ${\varepsilon_t}$ is strictly stationary or driven by i.i.d. errors. Instead, they require the existence of the $4 + \delta$ moment of $\varepsilon_t$. This difference arises because strong mixing processes use uniform integrability arguments to establish \eqref{pt.MRIC.2.6}, whereas the process specified in Proposition
%\ref{propB.3}
%utilizes H$\ddot{\mathrm{o}}$lder's inequality for the same purpose; see Section \ref{Ex3-S} of the Supplementary Material for more details.
%Assumption \eqref{App.1} or \eqref{App.2} is used to ensure that $\{\varepsilon_t\}$ exhibits weak dependence. 

The following examples further illustrate the flexibility of assumptions \eqref{pt.MRIC.1.3} and \eqref{pt.MRIC.2.6}.

%A discussion of \eqref{h.varepsilon square} has been provided in Section \ref{sec:3.1}.
%We now offer two examples showing that
%\eqref{t MRIC.1}--\eqref{t MRIC.4} are satisfied in common practice.

\begin{example}
\label{ex3.1}
 Assume that $\{\varepsilon_t\}$ %(see \eqref{varepsilon basic form})
is a GARCH($p^\prime,q^\prime$) process,
\eqref{varepsilon basic form} and
\eqref{apGARCH}
with
$\mu=2$ and $\lambda_i=0$,
satisfying \eqref{propB.2.1}. Moreover, suppose that the density of $z_t$ is continuous with respect to the Lebesgue measure and positive on $(-\infty,\infty)$. Then, by Proposition 12 of \cite{Carrasco2002}, $\{\varepsilon_t\}$ is a strong mixing process with exponentially decaying coefficients, ensuring that \eqref{App.2} holds. Consequently,
\eqref{pt.MRIC.1.3} and \eqref{pt.MRIC.2.6} are satisfied by Proposition \ref{propB.2}.

%\begin{equation}\label{remark t.MRIC.1}
%E|\varepsilon_1|^{4}<\infty.
%\end{equation}
%Then, by Theorem 3.1 of \cite{Ling2002}, \eqref{propA.1.1} holds with $\tilde{\mathbf{z}}_t=z_t$ and 
%\begin{equation*}
%\tilde{g}(\ldots,z_{t-1},z_{t})=z_t\{\varphi_0+\sum_{j=0}^{\infty}\tilde{\mathbf{c}}^\top(\prod_{i=1}^{j}\tilde{\mathbf{A}}_{t-i})\bm{\xi}_{t-j}\}^{1/2},
%\end{equation*}
%where
%$\tilde{\mathbf{c}}=(\varphi_1,\ldots,\varphi_{p^\prime},\psi_1,\ldots,\psi_{q^\prime})^\top$,
%$\bm{\xi}_{t}=(\varphi_0z_t^2,0,\ldots,0,\varphi_0,0,\ldots,0)$ is a $(p^\prime+q^\prime)$-dimensional vector, 
%with $\varphi_0z_t^2$ and $\varphi_0$ being its
%first and $(p^\prime+1)$th components, respectively, and
%\NiceMatrixOptions{cell-space-limits=10pt}
%\begin{equation*}
%\tilde{\mathbf{A}}_t=
%\setlength{\arraycolsep}{10pt}
%\begin{pNiceArray}{ccc|ccc}
%  \varphi_1z_t^2 & \cdots & \varphi_{p^\prime}z_t^2 & \psi_1z_t^2 & \cdots & \psi_{q^\prime}z_t^2 \\
%  \Block{1-2}{\mathbf{I}_{(p^\prime-1)\times(p^\prime-1)}} && \mathbf{O}_{(p^\prime-1)\times1} & \Block{1-3}{\mathbf{O}_{(p^\prime-1)\times q^\prime}} \\
%  \hline
%  \varphi_1 & \cdots & \varphi_{p^\prime}& \psi_1& \cdots & \psi_{q^\prime} \\
%  \Block{1-3}{\mathbf{O}_{(q^\prime-1)\times p^\prime}} &&& \Block{1-2}{\mathbf{I}_{(q^\prime-1)\times(q^\prime-1)}} && \mathbf{O}_{(q^\prime-1)\times1}
%\end{pNiceArray}.
%\end{equation*}
%Note that for any positive integers $r$ and $s$,
%$\mathbf{I}_{r \times r}$ denotes the $r \times r$ identity matrix, and $\mathbf{O}_{r \times s}$ denotes the $r \times s$ zero matrix. Moreover, by Proposition 3 of \cite{Wu2005on}, \eqref{propB.1.2} follows.
%Consequently,
%\eqref{pt.MRIC.1.3} and \eqref{pt.MRIC.2.6}
%are ensured by Proposition \ref{propB.1}.
\end{example}
%We see that $\varepsilon_{t_1}^2\varepsilon_{t_2}^2-E(\varepsilon_{t_1}^2\varepsilon_{t_2}^2)$ can be decomposed into two components.
%The first component is a linear sum of martingale sequence $\{h_t\}$ with constant coefficients $\mathbf{b}^{(1)}_{t_1,t_2,s}$, where $h_t$ is a quartic %function in $\varepsilon_t$ and $\sigma_t$. This component arises from the presence of terms of form $w_i^2-E(w_i^2)$ in the decomposition of %$\varepsilon_{t_1}^2\varepsilon_{t_2}^2-E(\varepsilon_{t_1}^2\varepsilon_{t_2}^2)$; The second component is a linear sum of martingale sequence $\{w_t\}$ %with random coefficients $\mathbf{b}^{(2)}_{t_1,t_2,s}$, where $w_t$ is a quadratic function in $\varepsilon_t$ and $\sigma_t$. This term is due to the %existence of terms of form $w_iw_j-E(w_iw_j)$, $i\neq j$, in the decomposition of %$\varepsilon_{t_1}^2\varepsilon_{t_2}^2-E(\varepsilon_{t_1}^2\varepsilon_{t_2}^2)$.
\begin{example}
\label{ex3.2}
Assume that
$\{\varepsilon_t\}$ is a SV process satisfying \eqref{varepsilon basic form} and \eqref{SV model},
in which $\{z_t\}$ is a sequence of i.i.d. random variables with $E|z_1|^{4+\delta}<\infty$. Since $\sigma_t$ is log-normally distributed, all its finite moments exist. This property, combined with $E|z_1|^{4+\delta}<\infty$ and the independence of $\{\sigma_t\}$ and $\{z_t\}$, implies \eqref{propB.2.1}. Moreover, since $\{\sigma_t\}$ and $\{z_t\}$ are independent and $\{z_t\}$ is a sequence of i.i.d. random variables, Theorem 5.2(a) of \cite{Bradley2005} implies that it suffices for \eqref{App.2} to verify that $\{\sigma_t\}$ has an exponentially decaying mixing coefficient. Given that $\{\log\sigma_t^2\}$ satisfies the AR($\tilde{p}$) model in \eqref{SV model} with $1-\sum_{i=1}^{\tilde{p}}\tilde{a}_iz^i\neq0$ for all $|z|\leq1$, Theorem 1 of \cite{Mokkadem1988}
establishes that 
$\{\log\sigma_t^2\}$ has exponentially decaying mixing coefficients, implying the same for
$\{\sigma_t\}$. Consequently, Proposition
\ref{propB.2} ensures that
\eqref{pt.MRIC.1.3} and \eqref{pt.MRIC.2.6} are fulfilled.


%Since $1-\sum_{i=1}^{\tilde{p}}\tilde{a}_iz^i\neq0$ for all $|z|\leq1$, it follows that \eqref{propA.1.1} holds with $\tilde{\mathbf{z}}_t=(z_t,v_t)^\top$ and 
%\begin{equation*}
%\tilde{g}(\ldots,\tilde{\mathbf{z}}_{t-1},\tilde{\mathbf{z}}_{t})=z_t\exp(\frac{1}{2}\tilde{b}_0^\star+\frac{1}{2}\sum_{i=0}^{\infty}\tilde{b}_{i+1}v_{t-i}),
%\end{equation*}
%where $\tilde{b}_0^\star$ and $\tilde{b}_{i}$s are defined in Remark \ref{remark3.2}. Moreover, Since $\sigma_t$ is log-normally distributed, all its finite moments exist. This property, combined with $E|z_1|^{4}<\infty$ and the independence of $\{\sigma_t\}$ and $\{z_t\}$, implies \eqref{propB.1.1}.
%In addition, \eqref{propB.1.2} is ensured by \begin{equation}\label{S7.2.1} 
%\sum_{i=0}^{\infty}(E|\sigma_i-\sigma_i^\star|^4)^{1/4}<\infty, 
%\end{equation} 
%where $\sigma_i^\star$ is $\sigma_i$ with $v_0$ therein replaced by its independent copy. %$v_0^\star$. 
%As a result, Proposition
%\ref{propB.1} guarantees the validity of
%\eqref{pt.MRIC.1.3} and \eqref{pt.MRIC.2.6}.


%To demonstrate \eqref{S7.2.1}, note first that
%for all $i\geq0$,
%\begin{equation}\label{S7.2.2}
%\begin{split}
%    & E|\sigma_i-\sigma_i^\star|^4=E|e^{\frac{1}{2}w_i}-e^{\frac{1}{2}w_i^\star}|^4 \\
%   =   & E|\frac{1}{2}e^{\frac{1}{2}\tilde{w}_i}|^4||w_i-w_i^\star|^4\leq (E|\frac{1}{2}e^{\frac{1}{2}\tilde{w}_i}|^8)^{\frac{1}{2}}(E|w_i-w_i^\star|^8)^{\frac{1}{2}},
%\end{split}
%\end{equation}
%where $w_i=\log(\sigma_i^2)$, $w_i^\star=\log({\sigma_i^\star}^2)$, and $\tilde{w}_i$ satisfies $|\tilde{w}_i-w_i^\star|\leq |w_i-w_i^\star|$. Since for all $i\geq0$, all moments of $\sigma_i$ and $\sigma_i^\star$ exist, we have 
%$$\sup_{0\leq i<\infty}E|e^{\frac{1}{2}\tilde{w}_i}|^8<\infty.$$ 
%Moreover, 
%since $w_i$ follows the AR($\tilde{p}$) model, \eqref{SV model}, with $1-\sum_{i=1}^{\tilde{p}}\tilde{a}_iz^i\neq0$ for all $|z|\leq1$, 
%by Examples 1 and 2 of \cite{Wu2011},
%$$\sum_{i=0}^{\infty}(E|w_i-w_i^\star|^8)^{\frac{1}{8}}<\infty.$$
%Combining this result with \eqref{S7.2.2} establishes \eqref{S7.2.1}.  
%See Section S3 of the Supplementary Material for further details.
\end{example}
Finally, we note that while \eqref{pt.MRIC.1.3} and \eqref{pt.MRIC.2.6} can also be verified for the process considered in Proposition \ref{prop3.1}, doing so requires additional assumptions on
$\varepsilon_{t_1}^2\varepsilon_{t_2}^2 - E(\varepsilon_{t_1}^2\varepsilon_{t_2}^2)$,
for all $-\infty < t_1, t_2 < \infty$. Given the substantial technical complexity involved, we do not pursue this direction further in the present paper.




%In particular, by assuming that this quantity follows a martingale structure analogous to, though more intricate than, the formulation in \eqref{h.varepsilon square}, one can establish the validity of \eqref{pt.MRIC.1.3} and \eqref{pt.MRIC.2.6}. 
%Such a martingale structure can accommodate GARCH and SV processes. 
%As the corresponding technical arguments are substantially more involved, we do not pursue them further in this paper.




\section{Simulation Studies}\label{sec:4}
\subsection{Numerical Illustrations of Theorem \ref{tMSPE}}\label{sec:4.1}
This subsection provides three numerical examples to illustrate Theorem \ref{tMSPE}.
\begin{example}\label{ex4.1}
We generate $M=5000$ realizations from the following AR(2) model,
\begin{eqnarray}
\label{ing133}
x_t=-0.5x_{t-2}+\varepsilon_t,
\end{eqnarray}
where
$\varepsilon_t$ obeys
\eqref{varepsilon basic form} with $\{z_t\}$ being a sequence of
i.i.d. $N(0,1)$ random variables
and $\{\sigma_t\}$ being either a GARCH(1,1) process,
\begin{equation}\label{ing135}
\sigma_t^2=0.4+0.2\varepsilon_{t-1}^2+0.55\sigma_{t-1}^2,
\end{equation}
or a SV(1) (\eqref{SV model} with $\tilde{p}=1$) process,
\begin{equation}\label{ing136}
(1-0.98B)\log(\sigma_t^2)=0.01+v_t,
\end{equation}
in which
$\{v_t\}$ is a sequence of i.i.d. $N(0,0.04)$ random variables.
To evaluate predictive performance under model misspecification,
we adopt the working AR(1) model in 
\eqref{AR(k) misspecified multi} with $k=1$ and
approximate \eqref{ing2}
using
\begin{eqnarray}
\label{ing123}
g_{n,h}(k)=\frac{1}{M}\sum_{l=1}^{M}n\{\tilde{g}^{(l)}_{1,n,h}(k)+\tilde{g}^{(l)}_{2,n,h}(k)\},
\end{eqnarray}
where
$\tilde{g}^{(l)}_{1,n,h}(k)=(x^{(l)}_{n+h}-\hat{x}^{(l)}_{n+h}(k)-\varepsilon^{(l)}_{n,h,k})^2$,  and $\tilde{g}^{(l)}_{2,n,h}(k)=2(x^{(l)}_{n+h}-\hat{x}^{(l)}_{n+h}(k)-\varepsilon^{(l)}_{n,h,k})(\varepsilon^{(l)}_{n,h,k}-\tilde{\varepsilon}^{(l)}_{n,h})$. Here, $x^{(l)}_{n+h}$, $\hat{x}^{(l)}_{n+h}(k)$, $\varepsilon^{(l)}_{n,h,k}$, and $\tilde{\varepsilon}^{(l)}_{n,h}$
are the respective realizations of
$x_{n+h}$, $\hat{x}_{n+h}(k)$, $\varepsilon_{n,h,k}$, and $\tilde{\varepsilon}_{n,h}$
in the $l$-th simulation.
Note that $M^{-1}\sum_{l=1}^{M}n\tilde{g}^{(l)}_{1,n,h}(k)$ and $M^{-1}\sum_{l=1}^{M}n\tilde{g}^{(l)}_{2,n,h}(k)$
approximate the first and second expectations on the right-hand side of \eqref{ptMSPE.1} in Section \ref{appB} of the Supplementary Material.
Their sum yields an empirical approximation to
\eqref{ing2}.
We then compute $g_{h}(k)$ (the right-hand side of \eqref{tMSPE.result}) based on \eqref{ing133} and \eqref{ing135} (or \eqref{ing136}),
and assess the closeness between $g_{h}(k)$ and $g_{n,h}(k)$ by
the ratio $R_{n,h}=g_{n,h}(k)/g_{h}(k)$,
where %Two examples are considered as follows.
$n=500, 2000$ and $h=1,2,3,4,5$.
We report
$g_{h}(k)$ and $R_{n,h}$ in Table \ref{tab1}.
%The results of $R_{n,h}$ of this example are reported in Table \ref{tab1}.
%The corresponding multi-step prediction results are reported in Table \ref{tab1}..
%Suppose that
\end{example}

\begin{example}\label{ex4.2}
In this example,
we generate $M=5000$ realizations from
the MA(1) model,
\begin{eqnarray}
\label{ing134}
x_t=\varepsilon_t-0.8\varepsilon_{t-1},
\end{eqnarray}
instead of model \eqref{ing133}.
On the other hand, $\varepsilon_t$, the working model,
and the values of $n$ and $h$ in this example are the same as those in Example \ref{ex4.1}.
%but the fitted model is \eqref{AR(k) misspecified multi} with $k=1$. Suppose that \eqref{varepsilon basic form} holds with $z_t$ being i.i.d. $N(0,1)$ and $\sigma_t$ obeying either %\eqref{GARCH(1,1)} or \eqref{SV(1)}.
Table \ref{tab2} summarizes the corresponding  $g_{h}(k)$
and $R_{n,h}$.
\end{example}

\begin{table}
\caption{The values of $g_{h}(k)$ and $R_{n,h}$, with $k=1$, $n=500, 2000$, and $h=1,\ldots, 5$, in Example \ref{ex4.1}}
\label{tab1}
\centering
\begin{tabular}{ccccccccccc} \hline
    GARCH(1,1) & $h=1$ & 2 & 3 & 4 & 5 \\ \hline
    $g_{h}(1)$ & 3.454 & 2.425 & 3.760 & 3.702 & 3.986 \\ \hline
   $R_{500,h}$ & 0.814 & 0.910 & 0.795 & 0.928 & 1.058 \\
   $R_{2000,h}$ & 1.043 & 1.021 & 0.920 & 0.961 & 1.034 \\ \hline \hline
   SV(1) & $h=1$ & 2 & 3 & 4 & 5 \\ \hline
    $g_{h}(1)$ & 9.680 & 7.119 & 11.686 & 12.061 & 13.546 \\ \hline
   $R_{500,h}$ & 0.925 & 0.690 & 0.800 & 0.793 & 0.644 \\
   $R_{2000,h}$ & 1.082 & 0.945 & 1.041 & 0.902 & 1.055 \\ \hline
  \end{tabular}
\end{table}

\begin{table}
\caption{The values of $g_{h}(k)$ and $R_{n,h}$, with $k=1$, $n=500, 2000$, and $h=1,\ldots, 5$, in Example \ref{ex4.2}}
\label{tab2}
\centering
\begin{tabular}{ccccccccccc} \hline
    GARCH(1,1) & $h=1$ & 2 & 3 & 4 & 5 \\ \hline
    $g_{h}(1)$ & 2.964 & 5.808 & 5.324 & 4.961 & 4.689 \\ \hline
   $R_{500,h}$ & 0.679 & 0.984 & 0.976 & 0.953 & 0.929 \\
   $R_{2000,h}$ & 1.077 & 0.979 & 0.975 & 0.981 & 0.976 \\ \hline \hline
   SV(1) & $h=1$ & 2 & 3 & 4 & 5 \\ \hline
    $g_{h}(1)$ & 7.484 & 17.448 & 17.113 & 16.791 & 16.481 \\ \hline
   $R_{500,h}$ & 1.284 & 0.768 & 0.614 & 0.731 & 0.865 \\
   $R_{2000,h}$ & 0.979 & 0.948 & 0.902 & 0.997 & 0.947 \\ \hline
  \end{tabular}
\end{table}

It follows from \eqref{ing133}--\eqref{ing136}, \eqref{ing134}
and the normality of $\{z_t\}$ and $\{v_t\}$
that the assumptions in Theorem \ref{tMSPE}
are fulfilled by Examples \ref{ex4.1} and \ref{ex4.2}.
For the SV(1) error \eqref{ing136},
Tables \ref{tab1} and \ref{tab2}
reveal that whereas
the values of $R_{500,h}$, oscillating
between 0.614 and 1.284, are relatively distant from 1,
$R_{n,h}$ becomes very close to 1 as $n$ increases to 2000.
The behavior of $R_{2000,h}$ under the GARCH(1, 1) error \eqref{ing135}
is similar to that under the SV(1) error, although
$R_{500, h}$ under the former error is generally closer to 1 than the latter.
%which seems not satisfactory enough. Nevertheless, in a majority of cases (12 out of 20), $R_{500,h}$ still lies within an acceptable range (0.8--1.058). On the other hand, $g_{h}(1)$'s are well approximated %predicted
%by $g_{n,h}(1)$ when $n=2000$. Note that values of $R_{2000,h}$ fall between 0.902 and 1.082 for all cases listed in Tables \ref{tab1} and \ref{tab2}.
Therefore, we conclude that these numerical
results are aligned with
the asymptotic results displayed
in Theorem \ref{tMSPE}.


To examine whether finite-sample performance depends on the sign or
magnitude of the coefficients in Examples
\ref{ex4.1} and \ref{ex4.2}, Section
\ref{sec:add sim ex4.1 4.2}
of
the Supplementary Material evaluates $R_{n,h}$ for several representative
parameter values. The results show that $R_{n,h}$ remains close to~1
(typically within the range $(0.9,1.1)$), indicating that the finite-sample
behavior is largely insensitive to the choice of AR and MA coefficients.
In addition, Section
\ref{sec:add EGARCH HAR} extends this analysis to EGARCH and HAR errors, confirming that similar finite-sample behavior persists under more
persistent volatility dynamics. Section
\ref{sec:heavy tail errors}
further considers fat-tailed
(Student-$t$) innovations and shows that, although heavier tails increase
variability, the overall convergence pattern of $R_{n,h}$ remains
consistent with the asymptotic theory.
%is aligned with
 %Examples \ref{ex4.1} and \ref{ex4.2} demonstrate that the finite-sample behaviors of the multi-step MSPE of the LS predictor concur with the asymptotic results derived .
 
%We illustrate how sensitive \eqref{tMSPE.result} is to the moment assumption \eqref{tMSPE.moment} through next example. 
The following example illustrates the sensitivity of \eqref{tMSPE.result} to the moment condition \eqref{tMSPE.moment}.

\begin{example}\label{ex.S2.1}
We generate $M=5000$ realizations from the following
AR(1) model,
% As in Examples \ref{ex4.1} and \ref{ex4.2}, all results in this section
\begin{equation*}
\begin{split}
     & x_t=-0.5x_{t-1}+\varepsilon_t, \\
     & \varepsilon_t=\sigma_tz_t,\ \sigma_t^2=\varphi_0+\varphi_1\varepsilon_{t-1}^2+\psi_1\sigma_{t-1}^2,
\end{split}
\end{equation*}
where $\{z_t\}$ is a sequence of
i.i.d. $N(0,1)$ random variables.
We perform prediction using
model \eqref{AR(k) misspecified multi} with $k=1$. Therefore, the working model is correctly specified.
Letting $(\varphi_0,\varphi_1,\psi_1)=(0.4,0.5,0.2)$ and $(0.4,0.2,0.55)$,
we report $g_{h}(k)$ and $R_{n,h}$, with $k=1$, $h=1$,
and $n=500, 2000, 5000$, in Table \ref{tabS1}.
\end{example}

\begin{table}
\caption{The values of $g_{h}(k)$ and $R_{n,h}$, with $h=1$, $k=1$, and $n=500, 2000, 5000$, in Example \ref{ex.S2.1}}
\label{tabS1}
\centering
\begin{tabular}{cccc} \hline
    \multicolumn{2}{c}{$(\varphi_0,\varphi_1,\psi_1)=(0.4,0.5,0.2)$} &  \multicolumn{2}{c}{$(\varphi_0,\varphi_1,\psi_1)=(0.4,0.2,0.55)$}\\ \hline
   $g_{1}(1)$ & 105.579 &  & 2.571\\ \hline
   $R_{500,1}$ & 0.041 & & 0.968 \\
   $R_{2000,1}$ & 0.053 &  & 0.961 \\
   $R_{5000,1}$ & 0.063 &   & 1.056 \\ \hline
  \end{tabular}
\end{table}

When $(\varphi_0,\varphi_1,\psi_1)=(0.4,0.5,0.2)$, $\varepsilon_t$
has only a finite $(4+\bar{\delta}_1)$-th moment, where $\bar{\delta}_1$ is some small positive number
(see Theorem 4 of \cite{N1990} for details). Table \ref{tabS1} shows that even
though $n$ increases to $5000$, $R_{n,1}$
still
deviates substantially from 1.
On the other hand, when $(\varphi_0,\varphi_1,\psi_1)=(0.4,0.2,0.55)$,
$\varepsilon_t$ has a
finite $(6+\bar{\delta}_2)$-th moment, for some small $\bar{\delta}_2>0$,
and hence \eqref{tMSPE.moment} follows.
In this latter case, the values of $R_{n,1}$, falling between 0.961--1.056,
are very close to 1 as long as $n\geq 500$.
These results indicate that the moment condition
\eqref{tMSPE.moment}
imposed in Theorem \ref{tMSPE} is essential and unlikely to be further relaxed.



%These results suggest that the moment condition \eqref{tMSPE.moment}
%used in Theorem \ref{tMSPE} can hardly be relaxed. %to $\sup_{-\infty<t<\infty}E|\varepsilon_t|^{4+\delta}<\infty$.

\subsection{MRIC for Subset Selection}\label{sec:4.2}
In this subsection, we first illustrate the finite sample performance of MRIC (see \eqref{MRIC0} and Section \ref{sec:3.2})
through the following data-generating process,
\begin{equation*}
\begin{split}
\label{ing137}
     & x_t=0.4x_{t-3}+\varepsilon_t, \\
     & \varepsilon_t=\sigma_tz_t,\ \sigma_t^2=0.4+0.2\varepsilon_{t-1}^2+0.55\sigma_{t-1}^2,  \\
\end{split}
\end{equation*}
where $\{z_t\}$ is a sequence of i.i.d. $N(0,1)$ random variables.
%Model \eqref{ing137}
%implies that $\{x_t\}$ and $\{\varepsilon_t\}$ are strictly stationary,
%and for some $\delta_1>0$, $E|\varepsilon_1|^{6+\delta_1}<\infty$,
%ensuring that $\hat{g}_h(J)$ (see \eqref{ing125})
%in MRIC is a reliable estimate of
%$g_{h}(J)=g_h(J)$, as argued in Section \ref{sec:3.1}.
These model specifications imply that
the assumptions of Theorem \ref{t MRIC} follow.
We are interested in performing $h$-step-ahead forecast, $h=1, 2, 3$,
based on two predictive models, $J_1=\{1\}$ and $J_2=\{2\}$, both are misspecified.
It can be shown that
$f_h(J_1)=f_h(J_2)$ and $g_h(J_1)>g_h(J_2)$ for $h=1$, $f_h(J_1)>f_h(J_2)$ for $h=2$, and
$f_h(J_1)<f_h(J_2)$ for $h=3$.
Therefore, for $h=1$, ${\cal J}_{1}(h)=\{J_1, J_2\}$ and ${\cal J}_{2}(h)=\{J_2\}$;
for $h=2$, ${\cal J}_{1}(h)={\cal J}_{2}(h)=\{J_2\}$;
for $h=3$, ${\cal J}_{1}(h)={\cal J}_{2}(h)=\{J_1\}$,
where ${\cal J}_{1}(h)$
and ${\cal J}_{2}(h)$
are defined in \eqref{M1}
and \eqref{M2}, respectively.
%Their predictors of $x_{n+h}$ are
%\begin{equation}\label{xt+h J1 J2}
%x_{t+h}=\bm{\beta}_h^\top(J_l)\mathbf{x}_{t}(J_l)+\varepsilon_{t,h,J_l}=\beta_h(\{l\})x_{t-l}+\varepsilon_{t,h,\{l\}},\ l=1,2.
%\end{equation}
%Denote (I) and (II) as the candidate models corresponding to $J_1$ and $J_2$ defined in \eqref{xt+h J1 J2}, respectively.
To implement MRIC,
we set $C_n=n^{0.6}$ as suggested in \cite{Hsu2019},
where $n$ is set to
500, 1000, 2000, 3000
in our study. While \eqref{ing126} allows choosing $C_n=n^{\nu}$ for any $\nu\in(0.5,1)$, we have found that increasing $\nu$ makes the MRIC less effective in identifying the model with the smallest $f_h$ value in finite samples. Thus, we choose a relatively small penalty parameter, $\nu=0.6$. 
 %and perform it to select the best predictive model between (I) and (II) based on $1000$ replications for and .
%($f_h(J_l)$ and $g_h(J_l)$ are the first- and second-order terms of the MSPE of $\hat{x}_{n+h}(J_l)$ defined in \eqref{M1} and \eqref{ing4}, respectively).
%Therefore, the best predictive model is (I) ((II)) for $h=3$ ($h=1,2$).
For comparison, we also use
$$\textrm{AIC}(J)=\log\{\hat{\sigma}^2_{h}(J)\}+\frac{2\sharp(J)}{n}$$ and
$$\textrm{BIC}(J)=\log\{\hat{\sigma}^2_{h}(J)\}+\frac{\sharp(J)\log n}{n}$$
to choose between $J_1$ and $J_2$.
The frequency, in 1000 simulations, of
each criterion to choose a candidate in ${\cal J}_{2}(h)$ is reported
%The frequencies of MRIC, AIC, and BIC choosing the better candidate model
in Table \ref{tab3}.
Note that since $\sharp(J_1)=\sharp(J_2)$, the model selection results
of AIC and BIC are exactly the same.
\begin{table}
\caption{Frequency, in 1,000 simulations, of choosing candidates belonging to ${\cal J}_{2}(h)$}
\label{tab3}
\centering
\begin{tabular}{cccc}
   $h$ & $n$ & AIC/BIC & MRIC \\ \hline
   1 & 500 & 476 & 695 \\
    & 1000 & 509 & 752  \\
    & 2000 & 502 & 790  \\
    & 3000 & 487 & 826  \\ \hline
   2 & 500 & 1000 & 999  \\
    & 1000 & 1000 & 1000  \\
    & 2000 & 1000 & 1000  \\
    & 3000 & 1000 & 1000  \\ \hline
   3 & 500 & 1000 & 997  \\
    & 1000 & 1000 & 1000  \\
    & 2000 & 1000 & 1000  \\
    & 3000 & 1000 & 1000  \\ \hline
  \end{tabular}
\end{table}

For $h=2$ and 3,
${\cal J}_{1}(h)$ only contains a single model,
and Table \ref{tab3} reveals that all criteria correctly select this model nearly 100\% of the time,
even when $n=500$.
This result is expected, as the difference
$|f_h(J_1)-f_h(J_2)|\approx0.305$ is bounded away from 0,
making it straightforward to identify the superior predictive model using consistent estimators,
$\hat{\sigma}^2_{h}(J_1)$ and
$\hat{\sigma}^2_{h}(J_2)$,
of $f_h(J_1)$ and $f_h(J_2)$;
see \eqref{ing128}.
In contrast, for $h=1$, we have
$f_h(J_1)=f_h(J_2)$,
so the distinction must rely on their second-order MSPEs.
Due to model misspecification,
$g_1(J_1)$ and $g_1(J_2)$
are not determined solely by model complexity, making it difficult for AIC and BIC to favor the better predictive model $J_2$. 
 %difficulty in identifying the best predictive model increases %becomes apparent
%because a comparison between $g_h(J_1)$ and $g_h(J_2)$ is now required.
%, whose penalties are proportional to the numbers of parameters in the candidate models,
%fail to distinguish
%$J_2$ from $J_1$,
%which satisfy
%$g_1(J_1)>g_{1}(J_2)$.
Indeed, Table \ref{tab3} shows that
these criteria select 
$J_2$ only about 50\% of the time, regardless of $n$.
%such a challenging case ($|f_1(J_1)-f_1(J_2) |=0$ and $|g_1(J_1)-g_1(J_2)|\approx0.224$),
By contrast, MRIC selects $J_2$ approximately 70\% of the time when $n=500$,
with the rate rising to about
80\% as $n$
increases to 2000.
These results highlight the advantage of MRIC in identifying the better predictive model under challenging conditions, where the under lying process exhibits conditional heteroscedasticity and the competing models are both misspecified yet yield the same population (first-order) MSPE.
%in terms of the values of $f_h(\cdot)$.

%In contrast, AIC/BIC can not decide which model is the best, and the corresponding frequencies of selecting (I) and (II) are about equal for all $n\geq500$. This shows the superiority of MRIC to other criteria that do not take into account the information of the second-order MSPE, $g_h(J_l)$, like AIC and BIC.

%\section{Concluding Remarks}\label{sec:5}
%In this work, we develop negative moment bounds for autocovariance matrices of stationary processes driven by conditional heteroscedastic errors. These results are then applied to establish the MSE and moment bounds of the LS estimator and an asymptotic expression of the multi-step MSPE of the LS predictor. Since none of the above results have been obtained for the case of conditional heteroscedastic errors before, this paper attempts to fill the lacuna %in the time series literature
%in the most exhaustive way. Note that our model framework is fairly general because it includes the two most popular classes of conditional heteroscedastic models, the GARCH- and SV-type models. Moreover, the model for forecasting is allowed to be misspecified. In addition, the results of the MSPEs are further applied to the scenario of best subset selection for AR models with conditional heteroscedastic errors, which is also an important and novel result.

%A future related topic worth investigating is to develop asymptotic expressions of the MSPEs of the LS predictor for nonstationary processes driven by conditional heteroscedastic innovations. A comparison of the behaviors of the MSPEs of the LS predictor under stationary and nonstationary processes in the presence of conditional heteroscedasticity should be intriguing.

%In Section \ref{sec:S9.2} of the Supplementary Material, we also compare the performance of MRIC, AIC, and BIC in correctly specified cases. Our results show that MRIC continues to outperform the other two criteria in such cases.
Next, we compare the performance of MRIC, AIC, and BIC in scenarios where some candidate models are correctly specified. 
Consider
the following two data-generating processes (DGPs):
%  \begin{equation*}
%\begin{split}
%  (I)\hspace{1cm} x_t=  & 0.8x_{t-1}-0.5x_{t-2}+0.35x_{t-4}+\varepsilon_t, \\
%     \varepsilon_t= &\sigma_tz_t,\ \sigma_t^2=0.4+0.2\varepsilon_{t-1}^2+0.55\sigma_{t-1}^2,  \\
%\end{split}
%\end{equation*}
\begin{itemize}
  \item [(I)]
  \begin{equation*}
\begin{split}
     & x_t=0.8x_{t-1}-0.5x_{t-2}+0.35x_{t-4}+\varepsilon_t, \\
     & \varepsilon_t=\sigma_tz_t,\ \sigma_t^2=0.4+0.2\varepsilon_{t-1}^2+0.55\sigma_{t-1}^2,  \\
\end{split}
\end{equation*}
  \item [(II)]
  \begin{equation*}
\begin{split}
     & x_t=0.55x_{t-2}-0.4x_{t-5}+\varepsilon_t, \\
     & \varepsilon_t=\sigma_tz_t,\ (1-0.98B)\log(\sigma_t^2)=0.01+v_t,  \\
\end{split}
\end{equation*}
\end{itemize}
where $\{z_t\}$ and $\{v_t\}$ are sequences of i.i.d. $N(0,1)$ and $N(0,0.04)$ random variables, respectively.
Let ${\cal J}=2^{\{1,2,3,4,5\}}$ be 
the power set of $\{1,2,3,4,5\}$,
representing the collection of candidate subset AR models for one-step-ahead prediction.
Then,
both DGPs (I) and (II),
referred to as the true models, are
contained in ${\cal J}$.
In addition, any candidate model that contains the true model yields the same population MSPE as the true model. 
However, among these models, only the true model achieves the minimum second-order MSPE,
a result that follows from 
Theorem \ref{tMSPE}.
Hence, we aim to identify the true model using MRIC, AIC, and BIC.
We perform 1,000 simulations and
record the frequency with which each criterion correctly identifies the true model
for sample sizes $n=200,\ 500,$ and $1000$. 
The $C_n$ in MRIC is still set to $n^{0.6}$ in these experiments.

\begin{table}
\caption{Frequency, in 1,000 simulations, of selecting the true model when the data is generated from DGP {\rm (I)}.}
\label{tabS9.1}
\centering
\begin{tabular}{cccc}
  $n$ & AIC & BIC & MRIC \\ \hline
    200 & 596 & 889 & 877 \\
    500 & 555 & 917 & 991  \\
    1000 & 558 & 945 & 1000 \\  \hline
  \end{tabular}
\end{table}

\begin{table}
\caption{Frequency, in 1,000 simulations, of selecting the true model when the data is generated from DGP {\rm (II)}.}
\label{tabS9.2}
\centering
\begin{tabular}{cccc}
  $n$ & AIC & BIC & MRIC \\ \hline
    200 & 416 & 835 & 950 \\
    500 & 375 & 842 & 989  \\
    1000 & 369 & 828 & 996 \\  \hline
  \end{tabular}
\end{table}
Under DGP (I), BIC slightly outperforms MRIC when 
$n=200$. However, as 
$n$ increases to 500 and beyond, 
MRIC quickly surpasses BIC in performance. 
In contrast, AIC consistently underperforms, 
showing minimal improvement with increasing sample size.
Under DGP (II), MRIC is the most effective method across all sample sizes, consistently identifying the true model more frequently than AIC or BIC. 
Section \ref{sec:add sub sel}
extends this comparison to settings with alternative error structures, including EGARCH, HAR, and IGARCH errors, as well as fat-tailed innovations.
The results indicate that the relative performance patterns of AIC, BIC, and MRIC remain broadly consistent across these settings, while model selection becomes notably more challenging under IGARCH errors, where variances do not exist.





%Notably, comparing the misspecified and correctly specified cases shows that the serial correlation in model errors induced by model misspecification increases the sample size $n$ required by MRIC to identify the optimal model.

\section{Proof of Theorem \ref{tfisher}}\label{sec:5}
\begin{proof}%[Proof of Theorem \ref{tfisher}]
The proof is divided into four steps to clarify the structure of
the argument. 
Throughout the remainder of this paper, we use
$C, C_1$, and $C_2$ to denote generic positive
constants independent of both $n$ and $k$, whose values may vary from place to place.




\vspace{0.3cm}
{\sc step 1 (initial simplification):}
Under \eqref{AR mean part}--\eqref{stationary char poly}, it can be shown that $x_t$ has the following AR($\infty$) representation,
\begin{equation}\label{ar infty}
x_t=\sum_{i=1}^{\infty}\beta_ix_{t-i}+\varepsilon_t,
\end{equation}
where $\sum_{i=1}^{\infty}|\beta_i|<\infty$. To see this, note that the $\beta_i$'s are the Taylor coefficients of the function $1/\alpha(z)$, where $\alpha(z) =
\sum_{i\geq0}\alpha_iz^i$. The summability of $\{\beta_i\}$ follows from the Wiener theorem (see Theorem 5.2(ii) of page 245 of \cite{Zygmund2002} and its corollary on the next page). This, together with Proposition 3.1.1 of \cite{Brockwell1991}, implies that \eqref{ar infty} holds true. Define
$$\mathbf{A}=\begin{pmatrix}
    1 & -\beta_1 & \cdots & -\beta_{k-1} \\
    0 & 1 & \ddots & \vdots \\
    \vdots & \ddots & \ddots & -\beta_1 \\
    0 & \cdots & 0 & 1
  \end{pmatrix},$$
and
\begin{equation}\label{pt.fisher.1}
\bm{\Phi}_j=\mathbf{A}\mathbf{x}_j(k)=\begin{pmatrix}
\varepsilon_j\\
\vdots\\
\varepsilon_{j-k+1}
\end{pmatrix}
+\bm{\eta}_{j,k},
\end{equation}
where the components of $\bm{\eta}_{j, k}$
are linear combinations of
$\{\varepsilon_{j-k},\varepsilon_{j-k-1},\ldots\}$
with absolutely summable coefficients,
and the dependence of $\bm{\Phi}_j$ on $k$ is suppressed to simplify the notation. 
Similar to (2.8) and (2.9) of \cite{Ing2003}, we can utilize $\lambda_{\min}^{-1}(\sum_{j=k}^{n-1}\mathbf{x}_j(k)\mathbf{x}_j^\top(k))\leq\lambda_{\max}(\mathbf{A}^\top\mathbf{A})\lambda_{\min}^{-1}(\sum_{j=k}^{n-1}\bm{\Phi}_j\bm{\Phi}_j^\top) \leq 
C \lambda_{\min}^{-1}(\sum_{j=k}^{n-1}\bm{\Phi}_j\bm{\Phi}_j^\top)$, along with the convexity of $x^{-q}, x>0$, to obtain
\begin{equation}\label{pt.fisher.2.0}
E\{\lambda_{\min}^{-q}(\hat{\mathbf{R}}_n(k))\}\leq \frac{C}{k}\sum_{j=0}^{k-1}\frac{1}{C_n}\sum_{s=0}^{C_n-1}(dk)^q E\{\lambda_{\min}^{-q}(\sum_{i=1}^{dk}\bm{\Phi}_{(i+sdk)k+j}\bm{\Phi}_{(i+sdk)k+j}^\top)\},
\end{equation}
where $d$ is some positive integer to be specified later,
$\lambda_{{\rm max}}(\mathbf{M})$
denotes the maximum eigenvalue of matrix $\mathbf{M}$,
$C_n=\lfloor\lfloor (n-k)/k\rfloor/(dk)\rfloor$, and $\lfloor z\rfloor$ denotes the largest integer less than or equal to $z$. By \eqref{pt.fisher.2.0}, proving \eqref{t.fisher.inverse moment}  reduces to establishing the following moment bound:
for every $q>0$ and integer $1\leq k <\infty$,
\begin{equation}\label{pt.fisher.2}
E\{\lambda_{\min}^{-q}(\sum_{i=1+l_0}^{dk+l_0}\bm{\Phi}_{ik+j}\bm{\Phi}_{ik+j}^\top)\}\leq  C_{1}\exp({C_{2}k\log k}),
\end{equation}
 uniformly over $0\leq j\leq k-1$ and $0\leq l_0\leq\lfloor (n-k)/k\rfloor-dk$.
Although 
$k$ is regarded as fixed,
we highlight its appearance in \eqref{pt.fisher.2} 
because the bound remains valid even if 
$k$ grows slowly with 
$n$, as discussed in Remark \ref{rm2.2} below.
Since the argument is identical for all $j$ and $l_0$, we
focus on the case
$j=0$ and $l_0=0$ in the remainder of the proof. 

Let $(q+3)/2<l<\infty\ \textrm{and}\ 1+(2/q)<\tilde{\theta}<\infty$ 
%\begin{equation}\label{l tilde theta}
%(q+3)/2<l<\infty\ \textrm{and}\ 1+(2/q)<\tilde{\theta}<\infty
%\end{equation}
be given. Then, for any $\tilde{M}>0$, we have
%\begin{eqnarray}\label{pt.fisher.3}
% \nonumber % Remove numbering (before each equation)
%    && E\{\lambda_{\min}^{-q}(\sum_{i=1}^{dk}\bm{\Phi}_{ik}\bm{\Phi}_{ik}^\top)\} =   \int_{0}^{\infty}P\{\lambda_{\min}(\sum_{i=1}^{dk}\bm{\Phi}_{ik}\bm{\Phi}_{ik}^\top)<u^{-\frac{1}{q}}\}du \nonumber\\
%   &\leq& \tilde{M}+\int_{\tilde{M}}^{\infty}P\{\inf\limits_{\lVert \bm{y}\rVert=1}\sum_{i=1}^{dk}(\bm{\Phi}_{ik}^\top\bm{y})^2<u^{-\frac{1}{q}},
%   \max\limits_{1\leq j\leq dk^2}\sigma_j^2<u^\theta,\sum_{i=1}^{dk}\lVert\bm{\Phi}_{ik}\rVert^2<\frac{u^{\frac{2l}{q}}}{k}\}du \\
% && +\int_{\tilde{M}}^{\infty}P(\max\limits_{1\leq j\leq dk^2}\sigma_j^2\geq u^\theta)du+\int_{\tilde{M}}^{\infty}P(\sum_{i=1}^{dk}\lVert\bm{\Phi}_{ik}\rVert^2\geq\frac{u^{\frac{2l}{q}}}{k})du \nonumber\\
%   &:=& \tilde{M}+(I)+(II)+(III), \nonumber
%\end{eqnarray}
\begin{equation}\label{pt.fisher.3}
\begin{split}
     & E\{\lambda_{\min}^{-q}(\sum_{i=1}^{dk}\bm{\Phi}_{ik}\bm{\Phi}_{ik}^\top)\} =   \int_{0}^{\infty}P\{\lambda_{\min}(\sum_{i=1}^{dk}\bm{\Phi}_{ik}\bm{\Phi}_{ik}^\top)<u^{-\frac{1}{q}}\}\mathrm{d}u \\
    \leq & \tilde{M}+\int_{\tilde{M}}^{\infty}P\{\inf\limits_{\lVert \bm{y}\rVert=1}\sum_{i=1}^{dk}(\bm{\Phi}_{ik}^\top\bm{y})^2<u^{-\frac{1}{q}},
   \max\limits_{1\leq j\leq dk^2}\sigma_j^2<u^{\tilde{\theta}},\max\limits_{1\leq i\leq dk}\lVert\bm{\Phi}_{ik}\rVert<\frac{u^{\frac{l}{q}}}{\sqrt{k}}\}\mathrm{d}u \\
     & +\int_{\tilde{M}}^{\infty}P(\max\limits_{1\leq j\leq dk^2}\sigma_j^2\geq u^{\tilde{\theta}})du+\int_{\tilde{M}}^{\infty}P({\max\limits_{1\leq i\leq dk}\lVert\bm{\Phi}_{ik}\rVert\geq \frac{u^{\frac{l}{q}}}{\sqrt{k}}})\mathrm{d}u \\
    := & \tilde{M}+{\rm (I)+(II)+(III)}.
\end{split}
\end{equation}
%In what follows, we choose $d$ and $\tilde{M}$ such that the terms (I), (II), and (III) are bounded. 
Since \eqref{varepsilon basic form} is assumed and $E(\varepsilon_t^2)$ is a finite constant, 
it is straightforward to show
that 
for any $d\geq 1$ and  $\tilde{M}>ck^q$,
where $c$ is any positive constant, 
\begin{equation}\label{pt.fisher.4}
{\rm (II)}\leq C dk^2\int_{\tilde{M}}^{\infty}u^{-{\tilde{\theta}}}\mathrm{d}u\leq C,
\end{equation}
and
\begin{equation}\label{pt.fisher.5}
{\rm (III)}\leq C dk^3\int_{\tilde{M}}^{\infty}u^{-\frac{2l}{q}}\mathrm{d}u\leq C.
\end{equation}

To deal with (I), we apply Lemma 5.2 of \cite{Eldar2012S}, a theory of covering numbers of the sphere, to obtain
\begin{equation}\label{pt.fisher.Qj}
\begin{split}
&P\bigg\{
\inf\limits_{\lVert \bm{y}\rVert=1}\sum_{i=1}^{dk}(\bm{\Phi}_{ik}^\top\bm{y})^2<u^{-\frac{1}{q}},\max\limits_{1\leq j\leq dk^2}\sigma_j^2<u^{\tilde{\theta}},{\max\limits_{1\leq i\leq dk}\lVert\bm{\Phi}_{ik}\rVert<\frac{u^{\frac{l}{q}}}{\sqrt{k}}}\bigg\} \\
\leq&
\sum_{v=1}^{m^\ast}P\{Q_v(u)\},
\end{split}
\end{equation}
where  $m^\ast$
is a positive integer satisfying
$m^{\ast} \leq {(1+2u^{(l+1/2)q^{-1}}/\sqrt{k})^k}$ and
\begin{equation*}
Q_v(u)=\bigcap_{i=1}^{dk}\{\inf\limits_{\bm{y}\in G_v}|\bm{y}^\top\bm{\Phi}_{ik}|<u^{-\frac{1}{2q}},\lVert\bm{\Phi}_{ik}\rVert\leq\frac{u^{\frac{l}{q}}}{\sqrt{k}},{\max\limits_{0\leq j\leq k-1}\sigma^2_{ik-j}}<u^{\tilde{\theta}}\},
\end{equation*}
{in which
$G_v, v=1,\ldots, m^\ast$,
are subsets of the $(k-1)$-sphere, $\mathbb{S}_{k-1}$,
satisfying $\bigcup_{v=1}^{m^{\ast}}G_v=\mathbb{S}_{k-1}$ and for any $\bm{y}_1,\bm{y}_2\in G_v$, $\lVert \bm{y}_1-\bm{y}_2\rVert\leq 2\sqrt{k}u^{-(l+1/2)q^{-1}}$.} 
Let $\bm{l}_v=(l_{v,1},\ldots ,l_{v,k})^{\top} \in G_v, 1\leq v\leq m^\ast$,
be arbitrarily chosen. Then, {similar to (2.12) of \cite{Ing2003}}, it holds that
\begin{equation*}
Q_v(u)\subseteq\bigcap_{i=1}^{dk}\{|\bm{l}_v^\top\bm{\Phi}_{ik}|\leq3u^{-\frac{1}{2q}},{\max\limits_{0\leq j\leq k-1}\sigma^2_{ik-j}}<u^{\tilde{\theta}}\}\equiv\bigcap_{i=1}^{dk}E_{v,i}(u).
\end{equation*}
and hence
\begin{equation}\label{pt.fisher.6}
{\rm (I)} \leq 
\int_{\tilde{M}}^{\infty}
\sum_{v=1}^{m^\ast} P\{Q_v(u)\} du {\leq} \int_{\tilde{M}}^{\infty}
\sum_{v=1}^{m^\ast}E(\prod_{i=1}^{dk}I_{E_{v,i}(u)})du.
\end{equation}
Owing to the intricate dependence on the
$\{\sigma_i\}$,
our treatment of term (I) {\it departs markedly } from that in
\cite{Ing2003}.

\vspace{0.3cm}
{\sc step 2 (A construction
of the conditional pdf with infinite-dimensional conditioning):}
To derive an upper bound for the expectation on the right-hand side of \eqref{pt.fisher.6},
 we first express the conditional probabilities
\begin{eqnarray*}
P(E_{v,i}(u)|\varepsilon_s,s\leq (i-1)k), \,\,i=dk,\ldots, 1,
\end{eqnarray*}
using the conditional pdf of 
$\bm{\varepsilon}_i=(\varepsilon_{ik},\ldots,
\varepsilon_{(i-1)k+1})$ given $\sigma(\varepsilon_j,j\leq (i-1)k)$, which we construct in this step.

Since $\sigma_i$
is $\sigma(\varepsilon_{i-1},\varepsilon_{i-2},\ldots)$-measurable,
Theorem 1.4.5 of \cite{Chow1997}
ensures the existence of a measurable function
$\tilde{\sigma}_{i}:\mathbb{R}^\infty\to [c_0, \infty)
$ such that
$\sigma_i=\tilde{\sigma}_{i}(\varepsilon_{i-1},\varepsilon_{i-2},\ldots)$ a.s., where $c_0$ is defined in
\eqref{ing101}.
%a regular conditional distribution for $\bm{\varepsilon}=\bm{\varepsilon}_{dk^2,k}=(\varepsilon_{dk^2}, \ldots,\varepsilon_{dk^2-k+1})$ given $\sigma(\varepsilon_j,j\leq dk^2-k)$ exists. 
%For each $i\in\mathbb{Z}$, let $\tilde{\sigma}_{i}:\mathbb{R}^\infty\to\mathbb{R}_+$ be a measurable function with respect to $\sigma(\varepsilon_{i-1},\varepsilon_{i-2},\ldots)$ satisfying $\tilde{\sigma}_{i}(\varepsilon_{i-1},\varepsilon_{i-2},\ldots)=\sigma_{i}$ almost surely. 
Given $t\in\mathbb{Z}$, $m\in\mathbb{N}$, and an $m$-dimensional vector $\mathbf{s}_{t,m}=(s_t,\ldots,s_{t-m+1})$, define, for
$t-m+1 \leq j \leq t$,
\begin{equation}\label{tilde sigma def}
\tilde{\sigma}^{(t,m)}_j(\mathbf{s}_{t,m})=
\begin{cases}
  \tilde{\sigma}_{j}(s_{j-1},\ldots,s_{t-m+1},\varepsilon_{t-m},\varepsilon_{t-m-1},\ldots), & \mbox{if } j=t-m+2,\ldots,t, \\
  \tilde{\sigma}_{t-m+1}(\varepsilon_{t-m},\varepsilon_{t-m-1},\ldots), & \mbox{if } j= t-m+1,
\end{cases}
\end{equation}
and
\begin{equation}\label{ing106}
p^{(t,m)}(\mathbf{s}_{t,m})=\prod_{j=t-m+1}^{t}p^{(t,m)}_j(\mathbf{s}_{t,m}),\,\,{\rm with}\,\, 
p^{(t,m)}_j(\mathbf{s}_{t,m})=\frac{1}{\tilde{\sigma}^{(t,m)}_j(\mathbf{s}_{t,m})}\phi_j(\frac{s_{j}}{\tilde{\sigma}^{(t,m)}_j(\mathbf{s}_{t,m})}).
\end{equation}
In Lemma \ref{lm.condi pdf} below, we show that for any real numbers $v_{ik},\ldots,v_{ik-k+1}$, the conditional distribution of $\bm{\varepsilon}_i$ given $\sigma(\varepsilon_j,j\leq (i-1)k)$ admits the  representation
\begin{align}\label{condi dist 3}
\begin{split}
    & P(\varepsilon_{ik}<v_{ik},\ldots,\varepsilon_{(i-1)k+1}<v_{(i-1)k+1} |\varepsilon_j,j\leq (i-1)k)\\ =& \int_{-\infty}^{v_{(i-1)k+1}}\cdots\int_{-\infty}^{v_{ik}}p^{(ik,k)}(\mathbf{s})\mathrm{d}\mathbf{s} \,\,\,\, \textrm{a.s.,}
\end{split}
\end{align}
with $\mathbf{s}=\mathbf{s}_{ik,k}=(s_{ik},\ldots,s_{(i-1)k+1})$.
By Theorem 7.2.2 of \cite{Chow1997},
a regular conditional distribution always exists; hence, for any Borel set 
$A \subset \mathbb{R}^k$,
we may take
$
P(\bm{\varepsilon}_i \in A |\varepsilon_j,j\leq (i-1)k)$
to be that regular version in the sequel.





%Note that $\tilde{\sigma}^{(dk^2,k)}_j(\mathbf{s})$ is $\sigma(\varepsilon_{dk^2-k},\varepsilon_{dk^2-k-1},\ldots)$-measurable and also $\sigma(z_{dk^2-k},\varepsilon_{dk^2-k-1},\varepsilon_{dk^2-k-2},\ldots)$-measurable. 
Now consider a class of rectangles
$$\mathcal{R}:=\{(-\infty,r_1)\times\cdots\times(-\infty,r_k):(r_1,\ldots,r_k)\in\mathbb{Q}^k\}.$$
Recall the probability space $(\Omega, {\cal F}, P)$ introduced at the beginning of Section~\ref{sec:2}. By \eqref{condi dist 3}, for each $R\in\mathcal{R}$, there exists a measurable set $\Omega_R\in\mathcal{F}$ with $P(\Omega_R)=1$ such that, for all $\omega\in\Omega_R$,
\begin{equation}
\label{ing250605}
\begin{split}
    & P(\bm{\varepsilon}_i\in R |\varepsilon_j,j\leq (i-1)k)(\omega) = \bigg(\underset{R}{\int\cdots\int}p^{(ik, k)}(\mathbf{s})\mathrm{d}\mathbf{s}\bigg)(\omega).
\end{split}
\end{equation}
Let $\Omega_0=\bigcap_{R\in\mathcal{R}}\Omega_R$. Since $\mathcal{R}$ is countable, $P(\Omega_0)=1$.
Next, 
let
$\mathcal{B}(\mathbb{R}^k)$ be the Borel $\sigma$-algebra on $\mathbb{R}^k$, and set
$$\mathcal{C}:=\{A\in \mathcal{B}(\mathbb{R}^k):
\forall \omega\in\Omega_0,
P(\bm{\varepsilon}_i\in A|\varepsilon_j,j\leq (i-1)k)(\omega)=\bigg(\underset{A}{\int\cdots\int}p^{(ik, k)}(\mathbf{s})\mathrm{d}\mathbf{s}\bigg)(\omega)
\}.$$
Because for each fixed  $\omega\in\Omega_0$,
the left-hand side defines a probability measure on 
$\mathcal{B}(\mathbb{R}^k)$,
the Lebesgue Monotone Convergence Theorem guarantees that
$\mathcal{C}$ is a $\lambda$-system. Since it contains the $\pi$-system $\mathcal{R}$
by \eqref{ing250605}, 
the $\pi$-$\lambda$ theorem
(Theorem 1.3.2 of 
\cite{Chow1997}) gives
$\mathcal{C}= \sigma(\mathcal{R})=\mathcal{B}(\mathbb{R}^k)$, where 
$\sigma(\mathcal{R})$ denotes the $\sigma$-algebra
generated by $\mathcal{R}$.
%the $\pi$-$\lambda$ theorem implies that $\mathcal{C}=\mathcal{B}(\mathbb{R}^k)$. 
Thus,
\begin{align}
\label{ing250602}
\begin{split}
P(\bm{\varepsilon}_i\in A|\varepsilon_j,j\leq (i-1)k)=\bigg(\underset{A}{\int\cdots\int}p^{(ik, k)}(\mathbf{s})\mathrm{d}\mathbf{s}\bigg),\,\,
\forall A\in \mathcal{B}(\mathbb{R}^k),
\,\,\mbox{a.s.}
\end{split}
\end{align}
Therefore,
we refer to $p^{(ik,k)}(\cdot)$
as the conditional pdf of
 $\bm{\varepsilon}_i$ given $\sigma(\varepsilon_j,j\leq k(i-1))$.
 Notably, this construction avoids specifying any joint density for the infinite-dimensional history
%It is worth noting that this conditional pdf %$(\varepsilon_t,\varepsilon_{t-1},\ldots,\varepsilon_{t-k+1})$ given $\sigma(\varepsilon_{j},j\leq t-k)$ can 
%is derived without requiring a specification of the joint density of the infinite-dimensional vector 
$(\varepsilon_j)_{j\leq k(i-1)}$,
which cannot be defined via Lebesgue measure in infinite-dimensional spaces.


By \eqref{ing250602}, we conclude that
\begin{equation}\label{pt.fisher.7}
\begin{split}
     & P(E_{v,i}(u)|\varepsilon_s,s\leq (i-1)k)=P(\bm{\varepsilon}_i \in A_{ik,k}(u)|\varepsilon_s,s\leq (i-1)k)\\ =&
     \underset{A_{ik, k}(u)}{\int\cdots\int}p^{(ik, k)}(\mathbf{s})\mathrm{d}\mathbf{s}\,\, \textrm{a.s.},
\end{split}
\end{equation}
where
\begin{equation*}%\label{ing108}
\begin{split}
&A_{ik,k}(u)\\
=&\bigg\{\mathbf{s}: \bigg|\sum_{j=0}^{k-1}l_{v,j+1}s_{ik-j}+\bm{l}_v^\top\bm{\eta}_{ik, k}\bigg|
\leq 3u^{-\frac{1}{2q}},\max\limits_{(i-1)k+1\leq j\leq ik}\{\tilde{\sigma}^{(ik,k)}_{j}(\mathbf{s})\}^2<u^{\tilde{\theta}}\bigg\}.
\end{split}
\end{equation*}

\vspace{0.3cm}
{\sc step 3 (Decoupling the conditional pdf using smoothness conditions):}
Although we have the integral representation in 
\eqref{pt.fisher.7},
directly bounding
$P(E_{v,i}(u)|\varepsilon_s,s\leq (i-1)k)$
is intractable, since 
$p^{(ik, k)}(\cdot)$
factors into $k$ multivariable functions with highly entangled arguments.
To address this challenge, we use the smoothness conditions \eqref{zt density}--\eqref{zt density 2}
to bound each factor, $p^{(ik,k)}_j(\cdot)$, of
$p^{(ik, k)}(\cdot)$ by an appropriate univariate function, thereby ``decoupling" the multidimensional dependence.








For notational simplicity, we henceforth write
$p^{(ik,k)}_j(\cdot)$ and $\tilde{\sigma}^{(ik,k)}_j(\cdot)$ as $p_j(\cdot)$ and $\tilde{\sigma}_j(\cdot)$, respectively.
We also require that $u>(c_0/c_2)^{4/(3\tilde{\theta})}\vee e \vee (6\sqrt{k}c_0^{-1})^{2q}$ throughout this step. 
Note that
$$p_j(\mathbf{s})=\frac{1}{\tilde{\sigma}_j(\mathbf{s})}\phi_j(\frac{s_j}{\tilde{\sigma}_j(\mathbf{s})}),$$
where $(i-1)k+1 \leq j\leq ik$
and $\mathbf{s}$ defined after \eqref{condi dist 3}.
A central observation in this step is that, for every $\mathbf{s}\in A_{ik,k}(u)$
and $(i-1)k+1 \leq j \leq ik$,
\begin{equation}\label{p j less bar p j}
p_j(\mathbf{s})\leq\bar{p}_j(s_j),
\end{equation}
where the univariate envelope
$\bar{p}_j$ depends on $u$ and is defined by
\begin{equation}\label{bar p j}
\bar{p}_j(s_j)= 
  \begin{cases}
    \frac{1}{c_0}\phi_j(\frac{s_j}{c_0}) & \mbox{for } |s_j|<\frac{c_0}{c_2}, \\
    \frac{M_{c^\star}}{c_1|s_j|} & \mbox{for } \frac{c_0}{c_2}\leq|s_j|\leq u^{\frac{3{\tilde{\theta}}}{4}}, \\
    \frac{1}{c_0}\phi_j(\frac{s_j}{u^{\frac{{\tilde{\theta}}}{2}}}) & \mbox{for } u^{\frac{3{\tilde{\theta}}}{4}}<|s_j|.
  \end{cases}
\end{equation}
Here $M_{c^\star}$ is $M_{\delta}$
(see \eqref{ing104}) with
$\delta=c^\star:=(1/c_2) \wedge 1$. 
To show \eqref{p j less bar p j}, first note that since
$\tilde{\sigma}_j(\mathbf{s})\geq c_0$,
\eqref{p j less bar p j} holds trivially
when $s_j=0$.
For $s_j\neq 0$ and
$|s_j|<c_0/c_2$, the assumption on $m_t(x)$
guarantees
$m_j(s_j)<c_2|s_j|<c_0$. This, along with \eqref{zt density 2}, indicates that $(1/c)\phi_j(s_j/c)$ is non-increasing in $c$ for all $c\geq c_0$. Given that $\tilde{\sigma}_j(\mathbf{s})\geq c_0$, it follows that $p_j(\mathbf{s})\leq(1/c_0)\phi_j(s_j/c_0)$. 
For $c_0/c_2
\leq|s_j|\leq u^{3{\tilde{\theta}}/4}$,
the conclusion follows by combining the bound
$p_j(\mathbf{s}) \leq (1/m_j(s_j))\phi_j(s_j/m_j(s_j))$
(from 
\eqref{zt density 0}),
with
\eqref{zt density}, the first relation in \eqref{ing104}, and the assumption on $m_t(x)$.
%$(1/\tilde{\sigma}_j(\mathbf{s}))g_j(s_j/\tilde{\sigma}_j(\mathbf{s}))
%Using this inequality with , we arrive at the desired conclusion. 
For $u^{3{\tilde{\theta}}/4}<|s_j|$, note that
if $\mathbf{s}\in A_{ik,k}(u)$, then
$\max\limits_{(i-1)k+1\leq j\leq ik}\tilde{\sigma}^2_{j}(\mathbf{s})<u^{\tilde{\theta}}$.
Together with \eqref{zt density},
this gives $\phi_j(s_j/\tilde{\sigma}_j(\mathbf{s}))\leq \phi_j(s_j/u^{\frac{{\tilde{\theta}}}{2}})$, and
since $\tilde{\sigma}_j(\mathbf{s})\geq c_0$,
the desired bound follows.


Let $k^\ast=\arg\max_{1\leq j\leq k}|l_{v,k+1-j}|$. Then, by \eqref{p j less bar p j}, the right-hand side of \eqref{pt.fisher.7} is bounded by
\begin{equation}\label{pt.fisher.7.1}
\underset{\bar{A}_{ik, k}(u)}{\int\cdots\int}\prod_{j=(i-1)k+k^{\ast}+1}^{ik}\bar{p}_j(s_j)\prod_{j=(i-1)k+1}^{(i-1)k+k^{\ast}}p_j(\mathbf{s})\mathrm{d}\mathbf{s},
\end{equation}
where 
\begin{equation*}\label{ing108}
\begin{split}
\bar{A}_{ik,k}(u)
=\bigg\{\mathbf{s}: \bigg|\sum_{j=0}^{k-1}l_{v,j+1}s_{ik-j}+\bm{l}_v^\top\bm{\eta}_{ik, k}\bigg|
\leq 3u^{-\frac{1}{2q}}\bigg\},
\end{split}
\end{equation*}
and we adopt the convention that
$\prod_{i=a}^{b} \cdot$ is equal to 1 whenever
$a>b$; that is, the product
vanishes from \eqref{pt.fisher.7.1}. 
%Since the first product in the integrand of \eqref{pt.fisher.7.1} depends only on $s_{ik},\ldots,s_{(i-1)k+k^{\star}+1}$, whereas the second product depends only on $s_{(i-1)k+k^{\star}},\ldots,s_{(i-1)k+1}$. Thus, 
In the following, we focus on the case
$1<k^{*}<k$, as the proofs for the cases
$k^{*}=1$ and $k^{*}=k$
are similar and simpler.
Equation
\eqref{pt.fisher.7.1}
can now be further rewritten as:
\begin{equation}\label{pt.fisher.7.2}
\int_{-\infty}^{\infty}\cdots\int_{-\infty}^{\infty}\bar{p}^{\prime}(\mathbf{s}^\prime)p^{\prime\prime}(\mathbf{s}^{\prime\prime})
\bigg(\int_{B(\mathbf{s}^{\prime},\mathbf{s}^{\prime\prime},u)}\phi_{(i-1)k+k^{\star}}(z)\mathrm{d}z\bigg)\mathrm{d}\mathbf{s}^{\prime}\mathrm{d}\mathbf{s}^{\prime\prime},
\end{equation}
where $\mathbf{s}^{\prime}=(s_{ik},\ldots,s_{(i-1)k+k^{\star}+1})$, $\mathbf{s}^{\prime\prime}=(s_{(i-1)k+k^{\star}-1},\ldots,s_{(i-1)k+1})$, 
$$\bar{p}^{\prime}(\mathbf{s}^\prime)=\prod_{j=(i-1)k+k^\star+1}^{ik}\bar{p}_j(s_j),\ p^{\prime\prime}(\mathbf{s}^{\prime\prime})=\prod_{j=(i-1)k+1}^{(i-1)k+k^\star-1}p_j(\mathbf{s}),$$
and $B(\mathbf{s}^{\prime},\mathbf{s}^{\prime\prime},u)$ is the interval of length 
$$\delta_B(\mathbf{s}^{\prime\prime},u)=6u^{-1/(2q)}/|l_{k+1-k^{\ast}}\tilde{\sigma}_{(i-1)k+k^{\ast}}(\mathbf{s})|
\leq k^{1/2}6u^{-1/(2q)}/c_0$$
centered at
$$-(\bm{l}_v^\top\bm{\eta}_{ik, k}-\mathop{{\sum_{j=0}^{k-1}}}_{j\neq k-k^\ast}l_{v,j+1}s_{ik-j})/(l_{k+1-k^{\ast}}\tilde{\sigma}_{(i-1)k+k^{\ast}}(\mathbf{s})).$$
Let $B(\mathbf{s}^{\prime},\mathbf{s}^{\prime\prime},u)$ and $\delta_B(\mathbf{s}^{\prime\prime},u)$ be abbreviated as $B(u)$ and $\delta_B(u)$, respectively. 
Since $B(u)$ may not be centered at 0,
three scenarios arise: (a) $B(u)\subset(0,\infty)$; (b) $B(u)\subset (-\infty,0)$; (c) $B(u)$ contains 0. In scenario (a), it follows from \eqref{zt density} that
\begin{equation}\label{pt.fisher.7.2.1}
\int_{B(u)}\phi_{(i-1)k+k^{\star}}(z)\mathrm{d}z\leq \int_{0}^{\delta_B(u)}\phi_{(i-1)k+k^{\star}}(z)\mathrm{d}z\leq \int_{-\delta_B(u)}^{\delta_B(u)}\phi_{(i-1)k+k^{\star}}(z)\mathrm{d}z.
\end{equation}
Similarly, in scenario (b),
\begin{equation}\label{pt.fisher.7.2.2}
\int_{B(u)}\phi_{(i-1)k+k^{\star}}(z)\mathrm{d}z\leq \int_{-\delta_B(u)}^{0}\phi_{(i-1)k+k^{\star}}(z)\mathrm{d}z\leq \int_{-\delta_B(u)}^{\delta_B(u)}\phi_{(i-1)k+k^{\star}}(z)\mathrm{d}z.
\end{equation}
In scenario (c), the bound is immediate:
\begin{equation}\label{pt.fisher.7.2.3}
\int_{B(u)}\phi_{(i-1)k+k^{\star}}(z)\mathrm{d}z\leq \int_{-\delta_B(u)}^{\delta_B(u)}\phi_{(i-1)k+k^{\star}}(z)\mathrm{d}z.
\end{equation}
As a consequence of \eqref{pt.fisher.7.2.1}--\eqref{pt.fisher.7.2.3} and the second part of condition \eqref{ing104}, we have
\begin{equation}\label{ing113}
\int_{B(u)}\phi_{(i-1)k+k^{\star}}(z)\mathrm{d}z\leq \bar{C}(k^{1/2}6u^{-1/(2q)}/c_0)^{\bar{\theta}}.
\end{equation}
By \eqref{ing113} and the fact that the arguments 
in $\bar{p}^{\prime}(\cdot)$
and $p^{\prime\prime}(\cdot)$ are decoupled, 
the integral in 
\eqref{pt.fisher.7.2} is bounded by
\begin{align}
\label{ing250610}
\begin{split}
\bar{C}(k^{1/2}6u^{-1/(2q)}/c_0)^{\bar{\theta}} 
\left(\int_{-\infty}^{\infty}\cdots\int_{-\infty}^{\infty}
 \bar{p}^{\prime}(\mathbf{s}^\prime)\mathrm{d}\mathbf{s}^{\prime} \right)
 \left(\int_{-\infty}^{\infty}\cdots\int_{-\infty}^{\infty}
p^{\prime\prime}(\mathbf{s}^{\prime\prime}) \mathrm{d}\mathbf{s}^{\prime\prime}\right).
\end{split}
\end{align}


According to \eqref{bar p j}, we have
\begin{equation}\label{pt.fisher.7.3}
\begin{split}
     & \int_{-\infty}^{\infty}\cdots\int_{-\infty}^{\infty}
 \bar{p}^{\prime}(\mathbf{s}^\prime)\mathrm{d}\mathbf{s}^{\prime}=
 \prod_{j=(i-1)k+k^\ast+1}^{ik}
\int_{-\infty}^{\infty}\bar{p}_{j}(s_j)ds_{j}
 \\
    \leq & \prod_{j=(i-1)k+k^\ast+1}^{ik}\bigg\{\int_{-\frac{c_0}{c_2}}^{\frac{c_0}{c_2}} \frac{1}{c_0}
   \phi_j(\frac{s_j}{c_0})\mathrm{d}s_j+\frac{M_{c^\star}}{c_1}\int_{\{\frac{c_0}{c_2}\leq |s_j|\leq u^{\frac{3{\tilde{\theta}}}{4}}\}}\frac{1}{|s_j|} \mathrm{d}s_j\\
   &+\frac{1}{c_0}\int_{\{u^{\frac{3{\tilde{\theta}}}{4}}<|s_j|\}} \phi_j(\frac{s_j}{u^{\frac{{\tilde{\theta}}}{2}}})\mathrm{d}s_j \bigg\}\\
   \leq & \prod_{j=(i-1)k+k^\ast+1}^{ik}\bigg\{\tilde{c}_1+\tilde{c}_2\log u+\frac{3}{c_0}\int_{\{u^{\frac{{\tilde{\theta}}}{4}}<|\tilde{s}_j|\}} \tilde{s}_j^2 \phi_j(\tilde{s}_j)d\tilde{s}_j\bigg\} \leq (\tilde{c}_3\log u)^{k},
\end{split}
\end{equation}
where $\tilde{c}_1,\tilde{c}_2$, and $\tilde{c}_3$
are some positive constants independent of $n$ and $k$. In addition, since  $p^{\prime\prime}(\mathbf{s}^{\prime\prime})$ is the conditional pdf of $(\varepsilon_{(i-1)k+k^{*}-1}, \ldots, \varepsilon_{(i-1)k+1})$
given $\{\varepsilon_s, s\leq (i-1)k\}$, we have
\begin{equation}\label{pt.fisher.7.4}
{\int_{-\infty}^{\infty}\cdots\int_{-\infty}^{\infty}
p^{\prime\prime}(\mathbf{s}^{\prime\prime}) \mathrm{d}\mathbf{s}^{\prime\prime}}=1.
\end{equation}
Combining \eqref{pt.fisher.7}, \eqref{pt.fisher.7.1}, \eqref{pt.fisher.7.2}, and \eqref{ing250610}--\eqref{pt.fisher.7.4} yields 
\begin{equation}
\label{pt.fisher.7.5}
P(E_{v,i}(u)|\varepsilon_s,s\leq (i-1)k)\leq \tilde{c}_4{k^{\frac{\bar{\theta}}{2}}}u^{-\frac{\bar{\theta}}{2q}}({\tilde{c}_3}\log u)^k,\ \textrm{a.s.},
\end{equation}
with $\tilde{c}_4=\bar{C}(6/c_{0})^{\bar{\theta}}$. 

\vspace{0.3cm}
{\sc step 4 (Iterated application and final bound):}
By iterating \eqref{pt.fisher.7.5}, we obtain,
for $u{>} (6\sqrt{k}c_0^{-1})^{2q} \vee e \vee (c_0/c_2)^{4/(3\tilde{\theta})}$,
\begin{align*}
\begin{split}
&E(\prod_{i=1}^{dk}I_{E_{v,i}(u)})=
E\{\prod_{i=1}^{dk-1}I_{E_{v,i}(u)}P(E_{v,dk}(u)|\varepsilon_s,s\leq dk^2-k)\}\\
&\leq 
\tilde{c}_4{k^{\frac{\bar{\theta}}{2}}}u^{-\frac{\bar{\theta}}{2q}}({\tilde{c}_3}\log u)^k
E(\prod_{i=1}^{dk-1}I_{E_{v,i}(u)}) 
\leq 
(\tilde{c}_4)^{dk}{k^{\frac{dk\bar{\theta}}{2}}}u^{-\frac{dk\bar{\theta}}{2q}}({\tilde{c}_3}\log u)^{dk^2}.
\end{split}
\end{align*}
%In view of
%\eqref{pt.fisher.6},
%repeating the same argument $dk-1$ times, one gets for $u{>} (6\sqrt{k}c_0^{-1})^{2q} \vee e$,
%\begin{equation*}%\label{pt.fisher.8}
%\begin{split}
%   P\{Q_v(u)\}\leq  & (\tilde{c}_4)^{dk}{k^{\frac{dk\bar{\theta}}{2}}}u^{-\frac{dk\bar{\theta}}{2q}}({\tilde{c}_3}\log u)^{dk^2},
%\end{split}
%\end{equation*}
Hence,
\begin{align}\label{pt.fisher.8}
\begin{split}
   \sum_{v=1}^{m^\ast}
   E(\prod_{i=1}^{dk}I_{E_{v,i}(u)})
   \leq  & (\tilde{c}_4)^{dk}m^\ast {k^{\frac{dk\bar{\theta}}{2}}}u^{-\frac{dk\bar{\theta}}{2q}}({\tilde{c}_3}\log u)^{dk^2}\\
   \leq&  (\tilde{c}_4)^{dk}{k^{-\frac{k}{2}}}u^{\frac{k(l+\frac{1}{2})}{q}}{k^{\frac{dk\bar{\theta}}{2}}}u^{-\frac{dk\bar{\theta}}{2q}}({\tilde{c}_3}\log u)^{dk^2}.
\end{split}
\end{align}
By letting 
\begin{equation}\label{d tilde M}
d>\frac{2l+1+2q}{\bar{\theta}}\ \textrm{and}\ \tilde{M}>\max\{e^{ak\log k},(6\sqrt{k}c_0^{-1})^{2q},e,
(c_0/c_2)^{4/(3\tilde{\theta})}
\},
\end{equation}
where $a=\{2dq(1+\theta)\}/(d\bar{\theta}-2l-1)$ with some small $\theta>0$, 
\eqref{pt.fisher.6}, along with
\eqref{pt.fisher.8}, yields
\begin{align}\label{pt.fisher.9}
\begin{split}
&  {\rm (I)} \leq  \int_{\tilde{M}}^{\infty}\ \sum_{j=1}^{m^\ast}P\{Q_j(u)\} du
% \leq  & C(\tilde{c}_4)^{dk}{k^{-\frac{k}{2}}}{k^{\frac{dk\bar{\theta}}{2}}}e^{-d(1+\theta)k^2\log k}e^{ak\log k}({\tilde{c}_3})^{dk^2}(ak\log k)^{dk^2} \\
   %= & Ce^{dk\log\tilde{c}_4}e^{\{(d\bar{\theta}-1)/2\}k\log k}e^{-d(1+\theta)k^2\log k}e^{ak\log k}e^{dk^2\log\tilde{c}_3}e^{dk^2(\log a+\log k+\log\log k)}\\
    \leq C_{\bm{\xi}_1}e^{-d\theta k^2\log k+\bar{f}(k)},
\end{split}
\end{align}
where $\bar{f}(k)$ satisfies
$\bar{f}(k)=o(k^2\log k)$ as $k\to\infty$ and $C_{\bm{\xi}_1}$ is a
 positive constant depending only on $\bm{\xi}_1=(q,c_0,c_1,c_2,\bar{C},\bar{\theta})$. Combining \eqref{pt.fisher.4}, \eqref{pt.fisher.5}, and \eqref{pt.fisher.9}, it follows that under the choice of $d$ and $\tilde{M}$ given in \eqref{d tilde M}, the quantity $\tilde{M}$+(I)+(II)+(III) in \eqref{pt.fisher.3} is bounded above by $C_{\bm{\xi}_1}\exp({C_{\bm{\xi}_2}k\log k})$, where $C_{\bm{\xi}_2}$
denotes a positive constant depending only on $\bm{\xi}_2=(q,\bar{\theta})$.
This, in turn, implies that \eqref{pt.fisher.2} holds for any positive integer $d$ satisfying the first part of \eqref{d tilde M}.
%Consequently, when $d$ and $\tilde{M}$ obey \eqref{d tilde M}, we have $\tilde{M}+(I)+(II)+(III)$ in \eqref{pt.fisher.3} is bounded by $C_{\bm{\xi}_1}\exp({C_{\bm{\xi}_2}k\log k})$, where $C_{\bm{\xi}_2}$ is defined in \eqref{pt.fisher.2}, which further implies that ensure that
%\eqref{pt.fisher.2} holds with
%$d$ satisfying the first statement in \eqref{d tilde M}.
%Since \eqref{pt.fisher.4} and \eqref{pt.fisher.5} remain valid when $d$ and $\tilde{M}$ are chosen according to \eqref{d tilde M}, it follows from \eqref{pt.fisher.3}--\eqref{pt.fisher.5} and \eqref{pt.fisher.9} that \eqref{pt.fisher.2} holds with
%$d$ satisfying the first statement in \eqref{d tilde M}.
%Consequently,
% \eqref{pt.fisher.3}--\eqref{pt.fisher.5} and
% \eqref{pt.fisher.9} ensure that
%\eqref{pt.fisher.2} holds with
%$d$ satisfying \eqref{d tilde M}.
\end{proof}




\begin{lemma}\label{lm.condi pdf}
Suppose that condition (CH)(i) holds and that $z_t$
admits a density $\phi_t(\cdot)$ with respect to the Lebesgue measure. Then, for any $t\in\mathbb{Z}$, any integer $1\leq k <\infty$, and
%the conditional distribution of $(\varepsilon_t,\varepsilon_{t-1},\ldots,\varepsilon_{t-k+1})$ given $\sigma(\varepsilon_{j},j\leq t-k)$
%can be represented as follows:for 
any real numbers $v_t,\ldots,v_{t-k+1}$, we have
\begin{equation}\label{l.condi pdf.1}
\begin{split}
&P(\varepsilon_t<v_t,\ldots,\varepsilon_{t-k+1}<v_{t-k+1}|\varepsilon_{j},j\leq t-k)\\
=&\int_{-\infty}^{v_{t-k+1}}\cdots\int_{-\infty}^{v_t}p^{(t,k)}(\mathbf{s}_{t,k})\mathrm{d}\mathbf{s}_{t,k}\,\,{\rm a.s.},
\end{split}
\end{equation}
where $\mathbf{s}_{t,k}=(s_{t},\ldots,s_{t-k+1})$ and $p^{(t,k)}(\cdot)$ is defined in \eqref{ing106}.
\end{lemma}

\begin{proof}[Proof of Lemma \ref{lm.condi pdf}]
Let $t \in \mathbb{Z}$ be given.
All subsequent statements involving
conditional probabilities are understood
to hold on $\bigcap_{j=t-k+1}^{t}\{c_0 \leq \sigma_{j} <\infty\}$, which
has probability 1.

We prove \eqref{l.condi pdf.1} 
by induction.
For the case of $k=1$, note that $z_t$ is independent of $\sigma(\varepsilon_{j},j\leq t-1)$ and 
$\sigma_{t}$ is $\sigma(\varepsilon_j,j\leq t-1)$-measurable. It follows that
\begin{equation}\label{p.lm.condi pdf.1}
\begin{split}
    &P(\varepsilon_t<v_t|\varepsilon_{j},j\leq t-1)= P(z_t<\frac{v_t}{\sigma_{t}}|\varepsilon_{j},j\leq t-1) \\
    =&  \int_{-\infty}^{\frac{v_t}{\sigma_{t}}}\phi_t(x)\mathrm{d}x 
   = \int_{-\infty}^{v_{t}}\frac{1}{\sigma_{t}}\phi_{t}(\frac{s_t}{\sigma_{t}})\mathrm{d}s_t
   = \int_{-\infty}^{v_{t}}\frac{1}{\tilde{\sigma}^{(t,1)}_{t}(s_t)}\phi_{t}(\frac{s_t}{\tilde{\sigma}^{(t,1)}_{t}(s_t)})\mathrm{d}s_t,\ \textrm{a.s.},
\end{split}
\end{equation}
where the third equality follows by the change of variable $s_t=\sigma_{t}x$, and the last one
uses the definition of $\tilde{\sigma}^{(t,1)}_{t}(\cdot)$ in \eqref{tilde sigma def}. Equation \eqref{p.lm.condi pdf.1} verifies \eqref{l.condi pdf.1} for $k=1$. 

Now suppose that \eqref{l.condi pdf.1} holds for $k=i$, and consider the case $k=i+1$. Observe that
\begin{equation}\label{p.lm.condi pdf.2}
\begin{split}
    & P(\varepsilon_t<v_t,\ldots,\varepsilon_{t-i}<v_{t-i}|\varepsilon_{j},j\leq t-i-1) \\
    = & E[I_{\{\varepsilon_{t-i}<v_{t-i}\}}\\
    &\hspace{3em}\times P(\varepsilon_t<v_t,\ldots,\varepsilon_{t-i+1}<v_{t-i+1}|\varepsilon_{j},j\leq t-i)|\varepsilon_{j},j\leq t-i-1],\ \textrm{a.s}.
\end{split}
\end{equation}
By the induction hypothesis, we have
\begin{align}\label{p.lm.condi pdf.3}
\begin{split}
&P(\varepsilon_t<v_t,\ldots,\varepsilon_{t-i+1}<v_{t-i+1}|\varepsilon_{j},j\leq t-i)\\
=
& \int_{-\infty}^{v_{t-i+1}}\cdots\int_{-\infty}^{v_t}\prod_{j=t-i+1}^{t}\frac{1}{{\tilde{\sigma}}^{(t,i)}_{j}({\mathbf{s}}_{t,i})}\phi_{j}(\frac{{s}_j}{{\tilde{\sigma}}^{(t,i)}_{j}({\mathbf{s}}_{t,i})})\mathrm{d}{\mathbf{s}}_{t,i},\ \textrm{a.s.}
\end{split}
\end{align}
Note that %$\tilde{\sigma}^{(t,i)}_{j}({\mathbf{s}}_{t,i})$ is measurable with respect to $\sigma(\varepsilon_{t-i},\varepsilon_{t-i-1},\ldots)$, and hence also with respect to  $\sigma(z_{t-i},\varepsilon_{t-i-1},\varepsilon_{t-i-2},\ldots)$, since $\varepsilon_{t-i}=\sigma_{t-i}z_{t-i}$ and $\sigma_{t-i}$ is $\sigma(\varepsilon_{t-i-1},\varepsilon_{t-i-2},\ldots)$-measurable. Specifically, 
for $j=t-i+2,\ldots,t$, we have
$\tilde{\sigma}^{(t,i)}_{j}({\mathbf{s}}_{t,i})=\tilde{\sigma}_j(s_{j-1},\ldots,s_{t-i+1},\varepsilon_{t-i},\varepsilon_{t-i-1},\ldots)=\tilde{\sigma}_j(s_{j-1},\ldots,s_{t-i+1},\sigma_{t-i}z_{t-i},\varepsilon_{t-i-1},\varepsilon_{t-i-2},\ldots)$ and for $j=t-i+1$, $\tilde{\sigma}^{(t,i)}_{j}({\mathbf{s}}_{t,i})=\tilde{\sigma}_j(\varepsilon_{t-i},$ $\varepsilon_{t-i-1},\ldots)=\tilde{\sigma}_j(\sigma_{t-i}z_{t-i},\varepsilon_{t-i-1},\varepsilon_{t-k-2},\ldots)$. 
Since $\sigma_{t-i}$ is measurable with respect to $\sigma(\varepsilon_{t-i-1},\varepsilon_{t-i-2},\ldots)$, it follows that $\tilde{\sigma}^{(t,i)}_{j}({\mathbf{s}}_{t,i})$ is $\sigma(z_{t-i})$-measurable, conditional on $\{\varepsilon_{j},j\leq t-i-1\}$. 
As a result, the integral on the right-hand side of \eqref{p.lm.condi pdf.3} is a measurable function of $z_{t-i}$, given $\{\varepsilon_{j},j\leq t-i-1\}$. 
%and so is the conditional probability on the left-hand side.
Additionally, under the same conditioning, it is easy to see that the indicator function $I_{\{\varepsilon_{t-i}<v_{t-i}\}}$ is a measurable function of 
$z_{t-i}$.

%Define $\tilde{\sigma}^{(t,i+1)}_j(s_{t},\ldots,s_{t-i+1},\sigma_{t-i}x)$
 % as $\tilde{\sigma}^{(t,i+1)}_{j}({\mathbf{s}}_{t,i},\sigma_{t-i}x)$, and note
  Since
$\tilde{\sigma}^{(t,i)}_j(\mathbf{s}_{t,i})$
can be expressed as
$\tilde{\sigma}^{(t,i+1)}_j(\mathbf{s}_{t,i},\sigma_{t-i}z_{t-i})$, 
taking the conditional expectation with respect to $z_{t-i}$, given $\{\varepsilon_{j},j\leq t-i-1\}$, yields
\begin{align}\label{p.lm.condi pdf.4}
\begin{split}
     & E[I_{\{\varepsilon_{t-i}<v_{t-i}\}}\\
     &\hspace{3em}\times P(\varepsilon_t<v_t,\ldots,\varepsilon_{t-i+1}<v_{t-i+1}|\varepsilon_{j},j\leq t-i)|\varepsilon_{j},j\leq t-i-1]\\
     =& \int_{-\infty}^{\frac{v_{t-i}}{\sigma_{t-i}}}\bigg(\int_{-\infty}^{v_{t-i+1}}\cdots\int_{-\infty}^{v_t}\prod_{j=t-i+1}^{t}\frac{1}{{\tilde{\sigma}}^{(t,i+1)}_{j}({\mathbf{s}}_{t,i},\sigma_{t-i}x)}\\
     &\hspace{12em}\times \phi_{j}(\frac{{s}_j}{{\tilde{\sigma}}^{(t,i+1)}_{j}({\mathbf{s}}_{t,i},\sigma_{t-i}x)})\mathrm{d}{\mathbf{s}}_{t,i}\bigg)\phi_{t-i}(x)\mathrm{d}x\\
    = & \int_{-\infty}^{v_{t-i}}\int_{-\infty}^{v_{t-i+1}}\cdots\int_{-\infty}^{v_t}\prod_{j=t-i}^{t}\frac{1}{\tilde{\sigma}^{(t,i+1)}_{j}({\mathbf{s}_{t,i+1}})}\\
    &\hspace{15em}\times \phi_{j}(\frac{{s}_j}{\tilde{\sigma}^{(t,i+1)}_{j}({\mathbf{s}_{t,i+1}})})\mathrm{d}{\mathbf{s}_{t,i+1}},\ \textrm{a.s.,}
\end{split}
\end{align}
where the second equality follows from the change of variables $s_{t-i}=\sigma_{t-i}x$,
and Tonelli's Theorem
(see Theorem 1.7.15 of \cite{Tao21}). This completes the induction step.
\end{proof}

\begin{remark}\label{rm2.2}
Our analysis in the proof of Theorems \ref{tfisher} remains valid when $k=k_n\to\infty$ at a slow rate. Specifically, under the assumptions of Theorem \ref{tfisher}, it can be shown that for any $q>0$ and $k=k_n= o(n^{1/2})$, both
\eqref{pt.fisher.2.0} and \eqref{pt.fisher.2}
continue to hold. Consequently,
\begin{equation}\label{GARCH inverse moment k}
E[\lambda_{\min}^{-q}(\hat{\mathbf{R}}_n(k))]\leq C_{1}k^{q}\exp({C_{2}k\log k}).
\end{equation}
Using \eqref{GARCH inverse moment k} and
following the argument in the proof of Theorem 2 of
\cite{Ing2003}, 
we further conclude the following: If $k=O((\log n)^\epsilon)$ for some $0<\epsilon<1$, then
\eqref{t.fisher.inverse moment} holds for $0<q<q_1$, provided that $\sup_{-\infty<t<\infty}E|\sigma_t|^{2(q_1\vee2)}<\infty$ and $\sup_{-\infty<t<\infty}E|z_t|^{2(q_1\vee2)}<\infty$ for some $q_1>0$. 
\end{remark}



\section{Concluding Remarks}
This paper takes an initial step toward developing a rigorous framework for prediction and model selection in weakly stationary processes with conditional heteroscedasticity, from the perspective of mean squared prediction error (MSPE). A central contribution lies in establishing a negative moment bound for the minimum eigenvalue of the sample autocovariance matrix. This is achieved through a highly strategic construction and analysis of conditional probability densities in infinite-dimensional spaces.
Specifically, we define a conditional density of the finite-dimensional vector 
$(\varepsilon_t,\varepsilon_{t-1},\ldots,\varepsilon_{t-k+1})$ given 
the infinite-dimensional history
$(\varepsilon_{j})_{j\leq t-k}$,
despite the absence of a Lebesgue measure in infinite-dimensional spaces. 
A key technical innovation is decoupling the highly entangled multivariate structure in the conditional density into a product of piecewise smooth univariate envelope functions.
These envelope bounds dominate each component of the joint density, allowing tractable integration and enabling the desired moment control.


The resulting negative moment bounds support an asymptotic MSPE decomposition that explicitly accounts for model complexity, misspecification, and time-varying volatility. This decomposition provides the theoretical basis for extending the MRIC criterion--proposed initially by
\cite{Hsu2019} for independent errors--to broader contexts involving weak dependence and conditional heteroscedasticity. 
Overall, the techniques and results in this paper lay a methodological foundation for extending MSPE-based inference to multivariate, nonstationary, nonlinear, or high-dimensional time series models with conditionally heteroscedastic errors.






%\begin{appendices}
\begin{appendix}
%\setcounter{section}{0}
%\renewcommand\thesection{Appendix \Alph{section}:}
\section{Proof of Theorem \ref{t SV negative bound} and Details on Remark \ref{rm2.1}}\label{appA}
\begin{proof}[Proof of Theorem \ref{t SV negative bound}]
%The proof of this result is similar to the proof of Theorem \ref{tfisher},
%we describe only those portions of the proof that require modification here.
In view of \eqref{ar infty}--\eqref{pt.fisher.3},
it suffices to show that
for some positive integer $d$ and some
positive real numbers $M^\ast$ and $\theta^\ast$,
%First, $E\{\lambda_{\min}^{-q}(\sum_{i=1}^{dk}\bm{\Phi}_{ik}\bm{\Phi}_{ik}^\top)\}$ in \eqref{pt.fisher.3} is now bounded by
\begin{align}\label{pc.SV neg 1}
\begin{split}
& E\{\lambda_{\min}^{-q}(\sum_{i=1}^{dk}\bm{\Phi}_{ik}\bm{\Phi}_{ik}^\top)\}\\
\leq&
M^\ast +\int_{M^\ast}^{\infty}P\big(\inf\limits_{\lVert \bm{y}\rVert=1}\sum_{i=1}^{dk}(\bm{\Phi}_{ik}^\top\bm{y})^2<u^{-\frac{1}{q}},\\
&\hspace{10em}
\min\limits_{1\leq j\leq dk^2}\sigma_j>u^{-\frac{1}{\theta^\ast}},{\max\limits_{1\leq i\leq dk}\lVert\bm{\Phi}_{ik}\rVert<\frac{u^{\frac{l}{q}}}{\sqrt{k}}}\big)du\\
   &+\int_{M^\ast}^{\infty}P(\min\limits_{1\leq j\leq dk^2}\sigma_j\leq u^{-\frac{1}{\theta^\ast}})du+\int_{M^\ast}^{\infty}P({\max\limits_{1\leq i\leq dk}\lVert\bm{\Phi}_{ik}\rVert\geq\frac{u^{\frac{l}{q}}}{\sqrt{k}}})du\\
   \equiv & M^\ast+{\rm (IV)+(V)+(VI)} {\leq C k^{(1+\theta)q}},
\end{split}
\end{align}
where $l$ is defined as in the proof of Theorem \ref{tfisher}, and $\theta$ is some small positive number. 
Note that the event $\{\max\limits_{1\leq j\leq dk^2}\sigma_j^2<u^{\tilde{\theta}}\}$
and $\{\max\limits_{1\leq j\leq dk^2}\sigma_j^2 \geq u^{\tilde{\theta}}\}$
in \eqref{pt.fisher.3} have now been replaced by
$\{\min\limits_{1\leq j\leq dk^2}\sigma_j>u^{-\frac{1}{\theta^\ast}}\}$ and
$\{\min\limits_{1\leq j\leq dk^2}\sigma_j \leq u^{-\frac{1}{\theta^\ast}}\}$,
respectively.
%for some .
{Let 
\begin{equation}\label{ing11118}
\begin{split}
     & \theta^\ast>\frac{2q(1+\theta)}{\theta},\,\,d>\bigg(\frac{2l+1+2q}{\rho(1-\frac{2q}{\theta^\ast})}\bigg)\vee\bigg[\frac{(2l+1+2q)(1+\theta)-1}{\rho\{(1-\frac{2q}{\theta^\star})(1+\theta)-1\}}\bigg], \,\,\mbox{and} \\
     & M^{\ast} >\big(\frac{6}{\eta}\big)^{2q}k^{(1+\theta)q},
\end{split}
\end{equation}
noting that $\eta$ and $\rho$
are defined in \eqref{smoothness condition zt}. By \eqref{ing102} and the assumption that $E(\varepsilon_t^2)$ is a finite constant,
it follows that
\begin{eqnarray}
\label{ing115}
{\rm (V)}\leq C  \,\,\mbox{and}\,\, {\rm (VI)}\leq C.
\end{eqnarray}} 
Therefore, it remains to show
\begin{eqnarray}
\label{ing116}
{\rm (IV)} \leq C.
\end{eqnarray}

By an argument similar to that used to prove \eqref{pt.fisher.Qj}, we obtain
\begin{eqnarray}
\label{ing117}
{\rm (IV)} \leq \int_{M^\ast}^{\infty}\sum_{v=1}^{m^\ast}P\{Q_v^{\ast}(u)\} du,
\end{eqnarray}
where $m^{\ast}$ is defined as in
\eqref{pt.fisher.Qj}, and
%Second, the $Q_j(u)$ in \eqref{pt.fisher.Qj} is replaced by
\begin{equation}\label{pc.SV neg 2}
\begin{split}
    Q_v^\ast(u)= &\bigcap_{i=1}^{dk}\{\inf\limits_{\bm{y}\in G_v}|\bm{y}^\top\bm{\Phi}_{ik}|<u^{-\frac{1}{2q}},\lVert\bm{\Phi}_{ik}\rVert\leq\frac{u^{\frac{l}{q}}}{\sqrt{k}},{\min\limits_{0\leq j\leq k-1}\sigma_{ik-j}}>u^{-\frac{1}{\theta^\ast}}\}  \\
    \subseteq & \bigcap_{i=1}^{dk}\{|\bm{l}_v^\top\bm{\Phi}_{ik}|\leq3u^{-\frac{1}{2q}},{\min\limits_{0\leq j\leq k-1}\sigma_{ik-j}}>u^{-\frac{1}{\theta^\ast}}\}\equiv\bigcap_{i=1}^{dk}E^\ast_{v,i}(u),
\end{split}
\end{equation}
with $G_v$ and $\bm{l}_v$ defined shortly after  \eqref{pt.fisher.Qj}. 
Assumption \eqref{smoothness condition zt} and the
independence between $\{z_t\}$ and $\{\sigma_t\}$
(see (CH)(ii)) yield that for
all $u\geq M^\ast$, %and any $\theta^\ast>0$
all $i=1,\ldots,dk$, and all $1\leq v \leq m^{\ast}$,
\begin{equation*}
\begin{split}
    & E(I_{E^\ast_{v,i}(u)}|z_s,s\leq ik-k,\sigma_t,t\leq ik-k)\\
    =& E[I(\min\limits_{0\leq j\leq k-1}\sigma_{ik-j}>u^{-\frac{1}{\theta^\ast}})E\{I(|\bm{l}_v^\top\bm{\Phi}_{ik}|\leq3u^{-\frac{1}{2q}})|z_s,s\leq ik,s\neq ik-k+k^\ast,\\
    &\sigma_t,t\leq ik\}|z_s,s\leq ik-k,\sigma_t,t\leq ik-k]\\
    \leq& E\{\bar{M}(\sqrt{k}\sigma^{-1}_{ik-k+k^\ast}6u^{-\frac{1}{2q}})^\rho I({\min\limits_{0\leq j\leq k-1}\sigma_{ik-j}}>u^{-\frac{1}{\theta^\ast}})|z_s,s\leq ik-k,\sigma_t,t\leq ik-k\}\\
    \leq& \bar{M}(6\sqrt{k})^\rho u^{(\frac{1}{\theta^\ast}-\frac{1}{2q})\rho},\ \textrm{a.s.},
\end{split}
\end{equation*}
which, together with \eqref{ing11118}, \eqref{ing117}, and \eqref{pc.SV neg 2}, gives
\begin{equation*}%\label{pc.SV neg 3}
  {\rm (IV)}\leq \bar{M}^{dk}6^{\rho dk} \int_{M^{\ast}}^{\infty} m^{\ast} {k^{\frac{\rho dk}{2}}}u^{-\rho d k (\frac{1}{2q}-\frac{1}{\theta^\ast})} du
  \leq C_{\bm{\xi}_3},
\end{equation*}
where $C_{\bm{\xi}_3}$ is a positive constant depending only on $\bm{\xi}_3=(q,\rho,\eta,\bar{M})$. 
Hence \eqref{ing116} follows.
%provided $q<\theta^\ast/2$ and
%$$d>\frac{2k(l+\frac{1}{2})+2q}{k\rho(1-\frac{2q}{\theta^\ast})}.$$
%Since $\theta^\ast$ can be arbitrary large, \eqref{pc.SV neg 3} is valid for any $q>0$. Moreover,
%it is not difficult to verify that $(V)<\infty$ and $(VI)<\infty$ for any $\theta^\ast>0$ and $q>0$ under (CH)(ii) and \eqref{t.fisher.moment}.
\end{proof}

In the rest of this section, we provide details on Remark \ref{rm2.1}. First, we show that Theorem \ref{tfisher} remains valid when \eqref{zt density}--\eqref{zt density 2} are replaced by \eqref{zt densityS}--\eqref{zt density 2S} given below. Next, we verify that \eqref{zt densityS}--\eqref{zt density 2S} include the mixture of normal density functions as special cases.

\begin{theorem}\label{tfisherS}
{Assume \eqref{AR mean part}--\eqref{stationary char poly} and {\rm (CH)(i)}}. Suppose for all $t \in \mathbb{Z}$, there exist positive constants $\bar{c}_1$, $\bar{c}_2$, and $\bar{c}_3$ such that
\begin{align}\label{zt densityS}
\begin{split}
& \phi_t(y)\leq \phi_t(x),\ \textrm{for all}\  \bar{c}_1\leq x< y<\infty,\\
& \phi_t(-y)\leq \phi_{t}(-x),\ \textrm{for all}\  \bar{c}_2\leq x< y<\infty,
\end{split}
\end{align}
and
\begin{eqnarray}
\label{ing104S}
\sup_{-\infty<t,x<\infty}\phi_t(x)\leq \bar{c}_3.
\end{eqnarray}
Moreover, for any $t\in \mathbb{Z}$ and $x\neq0$, there exist a positive number $m_t(x)$ and positive constants $\bar{c}_4$ and $\bar{c}_5$ obeying $\bar{c}_4|x| < m_t(x)< \bar{c}_5|x|$, such that
\begin{align}
\begin{split}
\label{zt density 0S}
& \zeta_{t,x}\{m_t(x)\}=\sup_{c>0} \zeta_{t,x}(c),
\end{split}
\end{align}
and 
\begin{equation}
\label{zt density 2S}
\begin{split}
\zeta_{t,x}(c) \,\,\mbox{is non-increasing for} \,\,
c>\bar{c}_5|x|.
\end{split}
\end{equation}
Then,
\eqref{t.fisher.inverse moment} follows.
\end{theorem}

\begin{proof}[Proof of Theorem \ref{tfisherS}]
{
We continue using the notation as in the proof of
Theorem \ref{tfisher}.
In particular, we
abbreviate 
$p^{(ik,k)}_j(\cdot)$ and $\tilde{\sigma}^{(ik,k)}_j(\cdot)$ as $p_j(\cdot)$ and $\tilde{\sigma}_j(\cdot)$, respectively,
for
$j=(i-1)k+1,\ldots,ik$.
To complete this proof, it suffices to replace two earlier bounds, \eqref{p j less bar p j}
and \eqref{ing113} from the proof of Theorem \ref{tfisher},
with the following two new bounds:
\begin{equation}\label{ptfisherS.1}
p_j(\mathbf{s})\leq\tilde{p}_j(s_j), \,\,\mbox{for}\,\,u> (\bar{c}_1\vee\bar{c}_2)^{4/\tilde{\theta}} \vee
(c_0/\bar{c}_5)^{4/(3\tilde{\theta})} \,\,\mbox{and}\,\,
 \mathbf{s}\in A_{ik,k}(u),
%\ \textrm{for all}\ j=(i-1)k+1,\ldots,ik\ \textrm{and}\ \mathbf{s}\in A_{ik,k}(u),
\end{equation}
and 
\begin{equation}\label{ptfisherS.2}
\int_{B(\mathbf{s}^{\prime},\mathbf{s}^{\prime\prime},u)}\phi_{(i-1)k+k^{\star}}(z)\mathrm{d}z\leq \bar{c}_3(k^{1/2}6u^{-1/(2q)}/c_0),\,\,\mbox{for}\,\, u>0,
\end{equation}
respectively.
Here, $\tilde{p}_j(s_j)$ is defined as
\begin{equation*}
\tilde{p}_j(s_j)= 
  \begin{cases}
    \frac{1}{c_0}\phi_j(\frac{s_j}{c_0}) & \mbox{for } |s_j|<\frac{c_0}{\bar{c}_5}, \\
    \frac{\bar{c}_3}{\bar{c}_4|s_j|} & \mbox{for } \frac{c_0}{\bar{c}_5}\leq|s_j|\leq u^{\frac{3\tilde{\theta}}{4}}, \\
    \frac{1}{c_0}\phi_j(\frac{s_j}{u^{\frac{\tilde{\theta}}{2}}}) & \mbox{for } u^{\frac{3\tilde{\theta}}{4}}<|s_j|.
  \end{cases}
\end{equation*}
Note that
$\tilde{\theta}$ is defined right before \eqref{pt.fisher.3},
and for a given $1\leq i \leq dk$,
\eqref{ptfisherS.1}
is required to hold
for all
$j=(i-1)k+1,\ldots,ik$.
%hold under \eqref{zt densityS}--\eqref{zt density 2S}.
%Then, the remainder of the proof is the same as the proof of Theorem \ref{tfisher}. 
%Note that except for $\tilde{p}_j(\cdot)$, $\bar{c}_3$, $\bar{c}_4$, and $\bar{c}_5$, all other notations in \eqref{ptfisherS.1} and \eqref{ptfisherS.2} are consistent with those used in the proof of Theorem \ref{tfisher}. 

 
 Since we have shown in the proof of 
 Theorem \ref{tfisher} that the length of $B(\mathbf{s}^{\prime},\mathbf{s}^{\prime\prime},u)$ is bounded by $k^{1/2}6u^{-1/(2q)}/c_0$,
\eqref{ptfisherS.2} follows directly from
\eqref{ing104S}.
To establish
\eqref{ptfisherS.1}, first observe that the inequality holds trivially 
when $s_j=0$.
For $s_j \neq 0$ and
$|s_j|<c_0/\bar{c}_5$, \eqref{zt density 2S} implies that $(1/c)\phi_j(s_j/c)$ is non-increasing in $c$ for all $c\geq c_0$. Since $\tilde{\sigma}_j(\mathbf{s})\geq c_0$, it follows that $p_j(\mathbf{s})\leq(1/c_0)\phi_j(s_j/c_0)$. For the case $c_0/\bar{c}_5\leq|s_j|\leq u^{(3\tilde{\theta})/4}$, using \eqref{ing104S}, \eqref{zt density 0S}, and the assumption on $m_t(x)$, we have 
\begin{equation*}\label{ptfisherS.3}
\frac{1}{\tilde{\sigma}_j(\mathbf{s})}\phi_j(\frac{s_j}{\tilde{\sigma}_j(\mathbf{s})})\leq \frac{1}{m_j(s_j)}\phi_j(\frac{s_j}{m_j(s_j)})\leq \frac{\bar{c}_3}{\bar{c}_4|s_j|},
\end{equation*}
which implies the desired bound. Finally, for $u^{3\tilde{\theta}/4}<|s_j|$, note that if $\mathbf{s}\in A_{ik,k}(u)$, then $\max\limits_{0\leq j\leq k-1}\tilde{\sigma}^2_{dk^2-j}(\mathbf{s})<u^{\tilde{\theta}}$. 
Combined with \eqref{zt densityS}, this yields
\begin{equation*}\label{ptfisherS.4}
\phi_j(\frac{s_j}{\tilde{\sigma}_j(\mathbf{s})})\leq \phi_j(\frac{s_j}{u^{\frac{\tilde{\theta}}{2}}}),
\end{equation*}
provided $u> (\bar{c}_1\vee\bar{c}_2)^{4/\tilde{\theta}}$.
Additionally, since $\tilde{\sigma}_j(\mathbf{s})\geq c_0$,
the bound
\eqref{ptfisherS.1} follows.
With both 
\eqref{ptfisherS.1} and \eqref{ptfisherS.2} verified,
the proof is complete.
}
\end{proof}

Now, assume for all $t \in \mathbb{Z}$,
\begin{equation}\label{gt mixture normal}
\phi_t(x)=\sum_{i=1}^{\bar{K}}\pi_i\phi_{t,i}(x)
\end{equation}
is a $\bar{K}$ ($\bar{K}\in\mathbb{N}$) component normal mixture density function, where
$$\phi_{t,i}(x)=\frac{1}{\sqrt{2\pi}\bar{\sigma}_i}\exp\bigg\{-\frac{1}{2}\bigg(\frac{x-\bar{\mu}_i}{\bar{\sigma}_i}\bigg)^2\bigg\},$$
with the normal mixture parameter $\bar{\bm{\eta}}=(\pi_1,\ldots,\pi_{\bar{K}},\bar{\mu}_1,\ldots,\bar{\mu}_{\bar{K}},\bar{\sigma}_1,\ldots,\bar{\sigma}_{\bar{K}})\in[0,1]^{\bar{K}}\times\mathbb{R}^{\bar{K}}\times(0,\infty)^{\bar{K}}$ obeying
\begin{equation}\label{gt mixture normal 2}
\sum_{i=1}^{\bar{K}}\pi_i=1,\ \sum_{i=1}^{\bar{K}}\pi_i\bar{\mu}_i=0,\ \textrm{and}\ \sum_{i=1}^{\bar{K}}\pi_i(\bar{\mu}_i^2+\bar{\sigma}_i^2)=1.
\end{equation}
Note that the second and the third equation in \eqref{gt mixture normal 2} ensure that $E(z_t)=0$ and $E(z_t^2)=1$. Then, it is easy to see that \eqref{zt densityS} and \eqref{ing104S} hold with $\bar{c}_1=\max\{1,\bar{\mu}_1,\ldots,\bar{\mu}_{\bar{K}}\}$, $\bar{c}_2=\max\{1,-\bar{\mu}_1,\ldots,-\bar{\mu}_{\bar{K}}\}$, and $\bar{c}_3=\sum_{i=1}^{\bar{K}}\pi_i/(\sqrt{2\pi}\bar{\sigma}_i)$. Furthermore, define $\zeta_{t,x,i}(c)=(1/c)\times$\\$\phi_{t,i}(x/c)$. Since for each $i=1,\ldots,\bar{K}$,
\begin{equation}\label{gt mixture normal 3}
\begin{split}
     & \zeta_{t,x,i}(c^+_i x)=\sup_{c>0} \zeta_{t,x,i}(c),\ x>0, \\
     & \zeta_{t,x,i}(c^-_i x)=\sup_{c>0} \zeta_{t,x,i}(c),\ x<0,
\end{split}
\end{equation}
and for $x>0$ ($x<0$),
\begin{equation}\label{gt mixture normal 4}
\begin{split}
     &  \zeta_{t,x,i}(c)\ \textrm{is increasing},\ \mbox{for } c<c^+_i x\ (c<c^-_i x),\\
     &  \zeta_{t,x,i}(c)\ \textrm{is decreasing},\ \mbox{for } c>c^+_i x\ (c>c^-_i x),
\end{split}
\end{equation}
it follows that \eqref{zt density 0S} and \eqref{zt density 2S} hold with
\begin{equation*}
\begin{split}
     & \bar{c}_4=\min\{c^+_1,\ldots,c^+_{\bar{K}},-c^-_1,\ldots,-c^-_{\bar{K}}\}, \\
     & \bar{c}_5=\max\{c^+_1,\ldots,c^+_{\bar{K}},-c^-_1,\ldots,-c^-_{\bar{K}}\},
\end{split}
\end{equation*}
and some positive $m_t(x)$ that depends on 
$x$ and the normal mixture parameter $\bar{\bm{\eta}}$.
Here,
$$c^+_i=\frac{\sqrt{\frac{\bar{\mu}_i^2}{\bar{\sigma}_i^4}+\frac{4}{\bar{\sigma}_i^2}}-\frac{\bar{\mu}_i}{\bar{\sigma}_i^2}}{2}\ \textrm{and}\
c^-_i=\frac{-\sqrt{\frac{\bar{\mu}_i^2}{\bar{\sigma}_i^4}+\frac{4}{\bar{\sigma}_i^2}}-\frac{\bar{\mu}_i}{\bar{\sigma}_i^2}}{2}.$$

\end{appendix}

\begin{supplement}
\stitle{}
\vspace{-1.5ex}
\sdescription{The Supplementary Material contains the proofs of
Theorems \ref{tMSPE} and \ref{t MRIC}, and
Propositions \ref{prop3.1}--\ref{propB.2}. It also includes additional numerical studies on MSPE and model selection.
%Lemmas \ref{lB.1} and \ref{lB.3}, \eqref{ptMSPE.S1-S6}, \eqref{ptMSPE.S7-S9}, \eqref{ptMSPE.9}, and \eqref{pt.MRIC.2.2}--\eqref{pt.MRIC.2.6}.
%It also provides technical details about Remark \ref{rm2.1} as well as additional numerical
%results.
}
\end{supplement}






%\setcounter{page}{1}
\setcounter{section}{0}
\setcounter{equation}{0}
\def\theequation{S\arabic{section}.\arabic{equation}}
\def\thetable{S\arabic{section}.\arabic{table}}
\def\thesection{S\arabic{section}}
\fontsize{12}{14pt plus.8pt minus .6pt}\selectfont
\setcounter{table}{0}
\setcounter{figure}{0}


%%%%%%%%%%%%%%%%%%%%%%%%%%%%%%%%%%%%%%%%%%%%%%%%%%%%%%%%%%%%%%%%%%%%%%%%%%%%%%%%%%%%%%%%

\section{Proofs of Theorem \ref{tMSPE} and Proposition \ref{prop3.1}}\label{appB}
%\section{Appendix A: Proofs of Theorem \ref{tMSPE}}\label{appB}
%\setcounter{equation}{0}
%\renewcommand{\theequation}{A.\arabic{equation}}

{The following lemma %deferred to Section \ref{AppA-S} of the supplement,
is required in the proof of Theorem \ref{tMSPE}.}

%This section includes three lemmas that are used to prove Theorem \ref{tMSPE}.
%Before we prove Theorem \ref{tMSPE}, three lemmas are needed first.

\begin{lemma}\label{lB.1}
Assume \eqref{AR mean part}--\eqref{stationary char poly}. Suppose for some $q_1, q_2\geq2$,
\begin{equation}\label{l1.moment1}
\sup\limits_{-\infty<t<\infty}E|\varepsilon_t|^{2q_1}<\infty,
\end{equation}
and {for any fixed sequences $\{\ubar{m}_{n}\}$ and $\{\bar{m}_{n}\}$ with $1\leq \ubar{m}_{n}\leq \bar{m}_{n}\leq n$,
\begin{equation}\label{l1.moment2}
\sup\limits_{-\infty<s<\infty}E|\frac{1}{\sqrt{\bar{m}_{n}-\ubar{m}_{n}+1}}\sum_{t=\ubar{m}_{n}}^{\bar{m}_{n}}(\varepsilon_{t+s}^2-E(\varepsilon_{t+s}^2))|^{q_2}\leq C.
\end{equation}
Define $q=\min\{q_1, q_2\}$.
Then, (a) for any $k\leq \ubar{m}_{n}\leq\bar{m}_{n}\leq n$,
\begin{equation}\label{l1.2}
E\lVert\frac{1}{\sqrt{\bar{m}_{n}-\ubar{m}_{n}+1}}\sum_{j=\ubar{m}_{n}}^{\bar{m}_{n}}\mathbf{x}_{j}(k)\varepsilon_{j,h,k}\rVert^q\leq Ck^{\frac{q}{2}}.
\end{equation}
(b) Moreover, we have
\begin{equation}\label{l2.result}
\sup\limits_{n\geq h+k+1}E\lVert\hat{\mathbf{R}}_{n,h}(k)-\mathbf{R}(k)\rVert^q\leq C\frac{k^{q}}{(n-h-k+1)^{\frac{q}{2}}}.
\end{equation}
(c) If we further assume that
\begin{equation}\label{inverse sample autocovariance moment}
E\lVert\hat{\mathbf{R}}^{-1}_{n,h}(k)\rVert^{\theta}=O(1),
\end{equation}
for any $\theta>0$, then for any $0<q^\prime<q$,
\begin{equation}\label{l3.result}
E\lVert\hat{\mathbf{R}}^{-1}_{n,h}(k)-\mathbf{R}^{-1}(k)\rVert^{q^\prime}=O(n^{-\frac{q^\prime}{2}}).
\end{equation}
}
\end{lemma}

\begin{proof}[Proof of Lemma \ref{lB.1}]
{We first prove \eqref{l1.2}. By the convexity of $x^{q/2}$, $x>0$, for any $k\leq \ubar{m}_{n}\leq\bar{m}_{n}\leq n$,
\begin{equation}\label{pl1.1}
\begin{split}
     & E\lVert\sum_{j=\ubar{m}_{n}}^{\bar{m}_{n}}\mathbf{x}_{j}(k)\varepsilon_{j,h,k}\rVert^q \leq  k^{\frac{q}{2}-1}\sum_{l=0}^{k-1}E|\sum_{j=\ubar{m}_{n}}^{\bar{m}_{n}}x_{j-l}\varepsilon_{j,h,k}|^q.
\end{split}
\end{equation}
Thus, it suffices for \eqref{l1.2} to prove that
\begin{equation*}
E|\sum_{j=\ubar{m}_{n}}^{\bar{m}_{n}}x_{j-l}\varepsilon_{j,h,k}|^q\leq C(\bar{m}_{n}-\ubar{m}_{n})^{\frac{q}{2}}.
\end{equation*}
By \eqref{AR mean part} and \eqref{AR(k) misspecified multi}, one has for all $0\leq l\leq k-1$,
\begin{equation}\label{pl1.2}
\begin{split}
     & E|\sum_{j=\ubar{m}_{n}}^{\bar{m}_{n}}x_{j-l}\varepsilon_{j,h,k}|^q \\
   =  & E|\sum_{j=\ubar{m}_{n}}^{\bar{m}_{n}}(x_{j-l}\varepsilon_{j,h,k}-E(x_{j-l}\varepsilon_{j,h,k}))|^q \\
   =  & E|\sum_{j=\ubar{m}_{n}}^{\bar{m}_{n}}[\{x_{j-l}(x_{j+h}-\bm{\beta}_h^\top(k)\mathbf{x}_j(k))\}\\
   &\hspace{10em}-E\{x_{j-l}(x_{j+h}-\bm{\beta}_h^\top(k)\mathbf{x}_j(k))\}]|^q \\
     =  & E|\sum_{i=0}^{\infty}\sum_{u=0}^{\infty}\alpha_i\alpha_u\sum_{j=\ubar{m}_{n}}^{\bar{m}_{n}}[\{\varepsilon_{j-l-i}(\varepsilon_{j+h-u}-\bm{\beta}_h^\top(k)\bm{\varepsilon}_{j-u}(k))\}\\
     &\hspace{10em}-E\{\varepsilon_{j-l-i}(\varepsilon_{j+h-u}-\bm{\beta}_h^\top(k)\bm{\varepsilon}_{j-u}(k))\}]|^q \\
     \leq &E|\sum_{i=0}^{\infty}\sum_{u=0}^{\infty}|\alpha_i||\alpha_u||\sum_{j=\ubar{m}_{n}}^{\bar{m}_{n}}[\{\varepsilon_{j-l-i}(\varepsilon_{j+h-u}-\bm{\beta}_h^\top(k)\bm{\varepsilon}_{j-u}(k))\}\\
     &\hspace{10em}-E\{\varepsilon_{j-l-i}(\varepsilon_{j+h-u}-\bm{\beta}_h^\top(k)\bm{\varepsilon}_{j-u}(k))\}]||^q\\
     \leq&\bigg(\sum_{i=0}^{\infty}|\alpha_i|\bigg)^{2q}E|\sum_{i=0}^{\infty}\sum_{u=0}^{\infty}p_{iu}|\sum_{j=\ubar{m}_{n}}^{\bar{m}_{n}}[\{\varepsilon_{j-l-i}(\varepsilon_{j+h-u}-\bm{\beta}_h^\top(k)\bm{\varepsilon}_{j-u}(k))\}\\
     &\hspace{10em}-E\{\varepsilon_{j-l-i}(\varepsilon_{j+h-u}-\bm{\beta}_h^\top(k)\bm{\varepsilon}_{j-u}(k))\}]||^q\\
     \leq&\bigg(\sum_{i=0}^{\infty}|\alpha_i|\bigg)^{2q}\sum_{i=0}^{\infty}\sum_{u=0}^{\infty}p_{iu}E|\sum_{j=\ubar{m}_{n}}^{\bar{m}_{n}}[\{\varepsilon_{j-l-i}(\varepsilon_{j+h-u}-\bm{\beta}_h^\top(k)\bm{\varepsilon}_{j-u}(k))\}\\
     &\hspace{10em}-E\{\varepsilon_{j-l-i}(\varepsilon_{j+h-u}-\bm{\beta}_h^\top(k)\bm{\varepsilon}_{j-u}(k))\}]|^q,
\end{split}
\end{equation}
where $\bm{\varepsilon}_j(k)=(\varepsilon_j,\ldots,\varepsilon_{j-k+1})^\top$ and $p_{iu}=|\alpha_i\alpha_u|/(\sum_{i=0}^{\infty}|\alpha_i|)^2$. In view of \eqref{alpha rate} and \eqref{pl1.2}, it remains to show that
\begin{align}\label{pl1.3}
\begin{split}
       & \sup\limits_{i,u\geq0}E|\sum_{j=\ubar{m}_{n}}^{\bar{m}_{n}}[\{\varepsilon_{j-l-i}(\varepsilon_{j+h-u}-\bm{\beta}_h^\top(k)\bm{\varepsilon}_{j-u}(k))\}\\
&\hspace{5em}-E\{\varepsilon_{j-l-i}(\varepsilon_{j+h-u}-\bm{\beta}_h^\top(k)\bm{\varepsilon}_{j-u}(k))\}]|^q \leq C(\bar{m}_{n}-\ubar{m}_{n})^{\frac{q}{2}}.
\end{split}
\end{align}
It is easily seen that this relation holds if
\begin{equation}\label{pl1.4}
\sup\limits_{i,j\in\mathbb{Z}}E|\sum_{t=\ubar{m}_{n}}^{\bar{m}_{n}}(\varepsilon_{t+i}\varepsilon_{t+j}-E\varepsilon_{t+i}\varepsilon_{t+j})|^q\leq C(\bar{m}_{n}-\ubar{m}_{n})^{\frac{q}{2}}.
\end{equation}
Note that when $i=j$, \eqref{pl1.4} follows directly from \eqref{l1.moment2}. For the case $i<j$, 
it follows from Burkholder's inequality, Minkowski's inequality, the Cauchy--Schwarz inequality, and \eqref{l1.moment1} that
\begin{equation}\label{pl1.5}
\begin{split}
     & E|\sum_{t=\ubar{m}_{n}}^{\bar{m}_{n}}(\varepsilon_{t+i}\varepsilon_{t+j}-E\varepsilon_{t+i}\varepsilon_{t+j})|^q\\
     = &E|\sum_{t=\ubar{m}_{n}}^{\bar{m}_{n}}\varepsilon_{t+i}\varepsilon_{t+j}|^q
      \leq CE|\sum_{t=\ubar{m}_{n}}^{\bar{m}_{n}}(\varepsilon_{t+i}\varepsilon_{t+j})^2|^{\frac{q}{2}}\\
     \leq& C\{\sum_{t=\ubar{m}_{n}}^{\bar{m}_{n}}(E|\varepsilon_{t+i}|^{2q})^{\frac{1}{q}}(E|\varepsilon_{t+j}|^{2q})^{\frac{1}{q}}\}^{\frac{q}{2}}\leq C(\bar{m}_{n}-\ubar{m}_{n})^{\frac{q}{2}}.
\end{split}
\end{equation}
The proof of \eqref{l1.2} is now complete. The proof of \eqref{l2.result} is analogous to the proof of \eqref{l1.2}, and is thus omitted. To prove \eqref{l3.result}, note that 
\eqref{AR mean part}--\eqref{stationary char poly} indicate that the spectral density function $f(\lambda)$ of $x_t$ obeys $f(\lambda)\leq \bar{f}$ for some $0<\bar{f}<\infty$, where $-\pi\leq\lambda\leq\pi$. This fact, together with Proposition 4.5.3 of \cite{Brockwell1991}, ensures that 
\begin{equation}\label{R invese 2norm 2}
\sup\limits_{k\geq1}\lVert\mathbf{R}^{-1}(k)\rVert\leq C.
\end{equation}
Then, it follows from \eqref{l2.result}, \eqref{inverse sample autocovariance moment}, and \eqref{R invese 2norm 2} that for all large $n$,
\begin{align}\label{pl3.1}
\begin{split}
   &E\lVert\hat{\mathbf{R}}^{-1}_{n,h}(k)-\mathbf{R}^{-1}(k)\rVert^{q^\prime} \\
    = &   E\lVert\hat{\mathbf{R}}^{-1}_{n,h}(k)\rVert^{q^\prime}\lVert\hat{\mathbf{R}}_{n,h}(k)-\mathbf{R}_n(k)\rVert^{q^\prime}\lVert\mathbf{R}^{-1}(k)\rVert^{q^\prime}\\
    \leq& C\{E\lVert\hat{\mathbf{R}}^{-1}_{n,h}(k)\rVert^{\frac{qq^\prime}{q-q^\prime}}\}^{\frac{q-q^\prime}{q}}\{E\lVert\hat{\mathbf{R}}_{n,h}(k)-\mathbf{R}(k)\rVert^{q}\}^{\frac{q^\prime}{q}}=O(n^{-\frac{q^\prime}{2}}).
\end{split}
\end{align}
Thus, \eqref{l3.result} follows.}
\end{proof}

With the help of Lemma \ref{lB.1}, we are now in a position to prove Theorem \ref{tMSPE}.

\begin{proof}[Proof of Theorem \ref{tMSPE}] We only prove (a) because
the proof of (b) is almost identical. Observe that
\begin{equation}\label{ptMSPE.1}
\begin{split}
     & n[E\{x_{n+h}-\hat{x}_{n+h}(k)\}^2-E(\varepsilon_{n,h,k}^2)] \\
    = & \frac{n}{n-h-k+1}E[\mathbf{x}^\top_n(k)\hat{\mathbf{R}}_{n,h}^{-1}(k)\{\frac{1}{\sqrt{n-h-k+1}}\sum_{i=k}^{n-h}\mathbf{x}_i(k)\varepsilon_{i,h,k}\}]^2\\
    &-\frac{2n}{n-h-k+1}E[\varepsilon_{n,h,k}\mathbf{x}^\top_n(k)\hat{\mathbf{R}}_{n,h}^{-1}(k)\{\sum_{i=k}^{n-h}\mathbf{x}_i(k)\varepsilon_{i,h,k}\}]\\
    &\equiv {\rm (I)+(II)}.
\end{split}
\end{equation}
{ Our proof is divided into three steps.\\

\noindent
\underline{Step 1}: \textit{Prove 
  \begin{equation}\label{ptMSPE.1.1}
     {\rm (I)}
   =  E(Y_n)+o(1),
\end{equation}
where
$$Y_n=\frac{1}{n-h-k+1}
\big(\sum_{i=k}^{n-h}\mathbf{x}^{\top}_i(k)\varepsilon_{i,h,k}\big)
\mathbf{R}^{-1}(k)\big(\sum_{i=k}^{n-h}\mathbf{x}_i(k)\varepsilon_{i,h,k}\big).$$}

Let $\{l_n\}$ be a sequence of positive integers satisfying
\begin{equation}\label{ln}
l_n\to\infty\ \ \textrm{and}\ \ l_n=o(n^{\frac{1}{2}}).
\end{equation}
Define
$$Z_n=\{\mathbf{x}^\top_n(k)\hat{\mathbf{R}}_{n,h}^{-1}(k)(\frac{1}{\sqrt{n-h-k+1}}\sum_{i=k}^{n-h}\mathbf{x}_i(k)\varepsilon_{i,h,k})\}^2$$
and
$$W_n=E(Z_n|\mathcal{F}_{n-l_n}).$$
Since
\begin{equation}\label{ptMSPE.1.1.00}
E(Z_n)=E(W_n)=E(W_n-Y_n)+E(Y_n),
\end{equation}
it suffices for \eqref{ptMSPE.1.1} to show that
\begin{align*}\label{ptMSPE.1.1.01}
\begin{split}
&{\rm (C1)}: \{W_n-Y_n\}\ \textrm{is uniformly integrable},\\ 
&{\rm (D1)}: W_n-Y_n=o_p(1).
\end{split}
\end{align*}

\noindent
\underline{Step 1.1} \textit{Proof of {\rm(C1)}}\\

First, note that \eqref{AR mean part}, \eqref{tMSPE.moment}, and the
absolute summability of $\{\alpha_i\}$ imply 
\begin{equation}\label{xt moment}
\sup\limits_{-\infty<t<\infty} E\lVert\mathbf{x}_t(k)\rVert^{6+\delta}<\infty.
\end{equation}
Moreover, by Theorem \ref{tfisher}, we have 
\begin{equation}\label{inverse sample autocovariance moment 2}
E\lVert\hat{\mathbf{R}}^{-1}_{n,h}(k)\rVert^{q}=O(1),
\end{equation}
for any $q>0$. Additionally, by applying \eqref{tMSPE.moment}, \eqref{tMSPE.moment2}, and Lemma \ref{lB.1}(a), one obtains
\begin{equation}\label{sum x varepsilon moment}
E\lVert\frac{1}{\sqrt{n-h-k+1}}\sum_{i=k}^{n-h}\mathbf{x}_i(k)\varepsilon_{i,h,k}\rVert^{3+\tilde{\delta}}= O(1),
\end{equation}
where $\tilde{\delta}=\min\{\delta/2,\delta^\star\}$. By using \eqref{xt moment}--\eqref{sum x varepsilon moment} and H$\ddot{\mathrm{o}}$lder's inequality, we have for any $0<\bar{\delta}<\tilde{\delta}/3$,
\begin{equation}\label{ptMSPE.1.1.02}
\begin{split}
 E|Z_n|^{1+\bar{\delta}}   \leq  & E\lVert\mathbf{x}^\top_n(k)\rVert^{2(1+\bar{\delta})}E\lVert\hat{\mathbf{R}}_{n,h}^{-1}(k)\rVert^{2(1+\bar{\delta})}\\
    &\times E\lVert\frac{1}{\sqrt{n-h-k+1}}\sum_{i=k}^{n-h}\mathbf{x}_i(k)\varepsilon_{i,h,k}\rVert^{2(1+\bar{\delta})}\\
    \leq & (E\lVert\mathbf{x}^\top_n(k)\rVert^{2(1+\bar{\delta})(3+\epsilon)})^{\frac{1}{3+\epsilon}}
    (E\lVert\hat{\mathbf{R}}_{n,h}^{-1}(k)\rVert^{2(1+\bar{\delta})(\frac{3+\epsilon}{\epsilon})})^{\frac{\epsilon}{3+\epsilon}}\\
    &\times
    (E\lVert\frac{1}{\sqrt{n-h-k+1}}\sum_{i=k}^{n-h}\mathbf{x}_i(k)\varepsilon_{i,h,k}\rVert^{2(1+\bar{\delta})(\frac{3+\epsilon}{2})})^{\frac{2}{3+\epsilon}}<\infty,
\end{split}
\end{equation}
where $\epsilon>0$ satisfies $3\bar{\delta}+\epsilon+\epsilon\bar{\delta}<\tilde{\delta}$. Similarly, using \eqref{sum x varepsilon moment} and the fact that
\begin{equation}\label{R invese 2norm}
\sup\limits_{k\geq1}\lVert\mathbf{R}^{-1}(k)\rVert\leq C,
\end{equation}
which is proved above \eqref{R invese 2norm 2}, we obtain for any $0<\bar{\delta}\leq1/2+\tilde{\delta}/2$,
\begin{equation}\label{ptMSPE.1.1.03}
E|Y_n|^{1+\bar{\delta}}<\infty.
\end{equation}
Furthermore, using conditional Jensen's inequality, we have, for any $\bar{\delta}>0$,
\begin{equation}\label{ptMSPE.1.1.04}
E|W_n|^{1+\bar{\delta}}\leq E|Z_n|^{1+\bar{\delta}}.
\end{equation}
Consequently, (C1) is ensured by the fact that 
$E|W_n-Y_n|^{1+\bar{\delta}}<\infty$ for some $\bar{\delta}>0$, which in turn
follows immediately from
\eqref{ptMSPE.1.1.02}, \eqref{ptMSPE.1.1.03}, and \eqref{ptMSPE.1.1.04}.\\

\noindent\noindent
\underline{Step 1.2} \textit{Proof of {\rm (D1)}}\\

First,  note that for all large $n$,
\begin{align}\label{ptMSPE.1.1.1}
\begin{split}
    W_n =& E[M_{2n}^\top\hat{\mathbf{R}}_{n,h}^{-1}(k)\mathbf{x}_n(k)\mathbf{x}^\top_n(k)
     \hat{\mathbf{R}}_{n,h}^{-1}(k)M_{1n}|\mathcal{F}_{n-l_n}]\\
     &+M_{3n}^\top E[\hat{\mathbf{R}}_{n,h}^{-1}(k)\mathbf{x}_n(k)\mathbf{x}^\top_n(k)
     \hat{\mathbf{R}}_{n,h}^{-1}(k)M_{2n}|\mathcal{F}_{n-l_n}]\\
     &+M_{3n}^\top E[\{\hat{\mathbf{R}}_{n,h}^{-1}(k)-\tilde{\mathbf{R}}_n^{-1}(k)\}\mathbf{x}_n(k)\mathbf{x}^\top_n(k)
     \hat{\mathbf{R}}_{n,h}^{-1}(k)|\mathcal{F}_{n-l_n}]M_{3n}\\
     &+M_{3n}^\top\tilde{\mathbf{R}}_n^{-1}(k)E[\mathbf{x}_n(k)\mathbf{x}^\top_n(k)
     \{\hat{\mathbf{R}}_{n,h}^{-1}(k)-\tilde{\mathbf{R}}_n^{-1}(k)\}|\mathcal{F}_{n-l_n}]M_{3n}\\
     &+M_{3n}^\top\tilde{\mathbf{R}}_n^{-1}(k)
     E\{\mathbf{x}_n(k)\mathbf{x}^\top_n(k)-\mathbf{R}(k)|\mathcal{F}_{n-l_n}\}\tilde{\mathbf{R}}_n^{-1}(k)M_{3n}\\
     &+M_{3n}^\top\tilde{\mathbf{R}}_n^{-1}(k)\mathbf{R}(k)\tilde{\mathbf{R}}_n^{-1}(k)M_{3n}\\
     :=& (S1)+(S2)+(S3)+(S4)+(S5)+(S6),
\end{split}
\end{align}
where $M_{1n}=(n-h-k+1)^{-1/2}\sum_{i=k}^{n-h}\mathbf{x}_i(k)\varepsilon_{i,h,k}$, $M_{2n}=(n-h-k+1)^{-1/2}\times$
$\sum_{i=n-h-l_n+1}^{n-h}\mathbf{x}_i(k)\varepsilon_{i,h,k}$, $M_{3n}=(n-h-k+1)^{-1/2}\sum_{i=k}^{n-h-l_n}\mathbf{x}_i(k)\varepsilon_{i,h,k}$ and
$\tilde{\mathbf{R}}_n(k)=(n-l_n-k+1)^{-1}\sum_{j=k}^{n-l_n}\mathbf{x}_j(k)\mathbf{x}_j^\top(k)$. In the following, we will show that
\begin{equation}\label{ptMSPE.S1-S6}
\begin{split}
     & (S1)=o_p(1),\ (S2)=o_p(1),\ (S3)=o_p(1),\ (S4)=o_p(1),   \\
     & (S5)=o_p(1),\ (S6)=Y_n+o_p(1).
\end{split}
\end{equation}
To show $(S1)=o_p(1)$, observe that by the conditional H$\ddot{\mathrm{o}}$lder inequality, one has
\begin{align}\label{ptMSPE.1.1.2}
\begin{split}
   |(S1)| \leq & \{E(\lVert M_{2n}\rVert^{3+\tilde{\delta}}|\mathcal{F}_{n-l_n})\}^{\frac{1}{3+\tilde{\delta}}}
    \{E(\lVert M_{1n}\rVert^{3+\tilde{\delta}}|\mathcal{F}_{n-l_n})\}^{\frac{1}{3+\tilde{\delta}}}\\
    &\times [E\{\lVert\hat{\mathbf{R}}_{n,h}^{-1}(k)\rVert^{\frac{6+2\tilde{\delta}}{\tilde{\delta}}}|\mathcal{F}_{n-l_n}\}]^{\frac{\tilde{\delta}}{3+\tilde{\delta}}}
    [E\{\lVert\mathbf{x}_n(k)\rVert^{6+2\tilde{\delta}}|\mathcal{F}_{n-l_n}\}]^{\frac{1}{3+\tilde{\delta}}}.
\end{split}
\end{align}
By virtue of Lemma \ref{lB.1}(a), \eqref{ln}, \eqref{xt moment}, and \eqref{inverse sample autocovariance moment 2}, it follows that
\begin{equation*}
\begin{split}
     & E(\lVert M_{1n}\rVert^{3+\tilde{\delta}}|\mathcal{F}_{n-l_n})=O_p(1),\ E(\lVert M_{2n}\rVert^{3+\tilde{\delta}}|\mathcal{F}_{n-l_n})=o_p(1),\\
     & E\{\lVert\hat{\mathbf{R}}_{n,h}^{-1}(k)\rVert^{\frac{6+2\tilde{\delta}}{\tilde{\delta}}}|\mathcal{F}_{n-l_n}\}=O_p(1),\ E\{\lVert\mathbf{x}_n(k)\rVert^{6+2\tilde{\delta}}|\mathcal{F}_{n-l_n}\}=O_p(1).
\end{split}
\end{equation*}
Combining these with \eqref{ptMSPE.1.1.2} leads to
\begin{equation}\label{ptMSPE.1.1.3}
(S1)=o_p(1).
\end{equation}
Similarly, it can be shown that
\begin{equation}\label{ptMSPE.1.1.4}
(S2)=o_p(1).
\end{equation}

To deal with $(S3)$, observe that for all large $n$,
\begin{equation}\label{ptMSPE.1.1.5}
\begin{split}
  & \lVert\hat{\mathbf{R}}_{n,h}^{-1}(k)-\tilde{\mathbf{R}}_n^{-1}(k)\rVert \\
  \leq & \lVert\hat{\mathbf{R}}_{n,h}^{-1}(k)\rVert\lVert\tilde{\mathbf{R}}_n^{-1}(k)\rVert\{\lVert\frac{1}{n-h-k+1}\sum_{j=n-l_n+1}^{n-h}\mathbf{x}_j(k)\mathbf{x}_j^\top(k)\rVert\\
    +&\lVert d_n\sum_{j=k}^{n-l_n}\mathbf{x}_j(k)\mathbf{x}_j^\top(k)\rVert\},
\end{split}
\end{equation}
where $d_n=1/(n-h-k+1)-1/(n-l_n-k+1)$
By an argument similar to that used in the proof of Theorem \ref{tfisher}, it can be shown that
\begin{equation}\label{ptMSPE.1.1.6}
E\lVert\tilde{\mathbf{R}}_n^{-1}(k)\rVert^{q}=O(1),
\end{equation}
for any $q>0$. Moreover, by \eqref{xt moment} and Minkowski's inequality, we have
\begin{equation}\label{ptMSPE.1.1.6.2}
\begin{split}
     & E\lVert\frac{1}{n-h-k+1}\sum_{j=n-l_n+1}^{n-h}\mathbf{x}_j(k)\mathbf{x}_j^\top(k)\rVert^{3+\tilde{\delta}} \\
    \leq & (n-h-k+1)^{-(3+\tilde{\delta})}E(\sum_{j=n-l_n+1}^{n-h}\lVert\mathbf{x}_j(k)\rVert^2)^{3+\tilde{\delta}}\\
    \leq & (n-h-k+1)^{-(3+\tilde{\delta})}\{\sum_{j=n-l_n+1}^{n-h}(E\lVert\mathbf{x}_j(k)\rVert^{6+\delta})^{\frac{1}{3+\tilde{\delta}}}\}^{3+\tilde{\delta}}\\
    =&O\{(l_n/n)^{3+\tilde{\delta}}\}.
\end{split}
\end{equation}
Analogously, it can be shown that
\begin{align}\label{ptMSPE.1.1.6.3}
\begin{split}
     & E\lVert d_n\sum_{j=k}^{n-l_n}\mathbf{x}_j(k)\mathbf{x}_j^\top(k)\rVert^{3+\tilde{\delta}}
     =O\{(l_n/n)^{3+\tilde{\delta}}\}.
\end{split}
\end{align}
Thus, it follows from \eqref{ln}, \eqref{inverse sample autocovariance moment 2}, \eqref{ptMSPE.1.1.5}--\eqref{ptMSPE.1.1.6.3} and H$\ddot{\mathrm{o}}$lder's inequality that
\begin{equation}\label{ptMSPE.1.1.7}
E\lVert\hat{\mathbf{R}}_{n,h}^{-1}(k)-\tilde{\mathbf{R}}_n^{-1}(k)\rVert^{\gamma_1}=O\{(l_n/n)^{\gamma_1}\},
\end{equation}
for $0<\gamma_1<3+\tilde{\delta}$. Using \eqref{ptMSPE.1.1.7} and an argument similar to that used in the proof of \eqref{ptMSPE.1.1.3}, one obtains
\begin{equation}\label{ptMSPE.1.1.8.1}
E[\{\hat{\mathbf{R}}_{n,h}^{-1}(k)-\tilde{\mathbf{R}}_n^{-1}(k)\}\mathbf{x}_n(k)\mathbf{x}^\top_n(k)
     \hat{\mathbf{R}}_{n,h}^{-1}(k)|\mathcal{F}_{n-l_n}]=o_p(1).
\end{equation}
A direct consequence of Lemma \ref{lB.1}(a) is
\begin{equation}\label{M 3n}
\lVert M_{3n}\rVert=O_p(1).
\end{equation}
Combining this with \eqref{ptMSPE.1.1.8.1} gives
\begin{equation}\label{ptMSPE.1.1.8}
(S3)=o_p(1).
\end{equation}
Similarly, we can derive that
\begin{equation}\label{ptMSPE.1.1.9}
(S4)=o_p(1).
\end{equation}

Recall that $\bm{\varepsilon}_t(k)=(\varepsilon_t,\ldots,\varepsilon_{t-k+1})^\top$. By \eqref{AR mean part} and Minkowski's inequality, one has
\begin{align}\label{ptMSPE.1.1.10}
\begin{split}
     & E\lVert E\{\mathbf{x}_n(k)\mathbf{x}_n^\top(k)-\mathbf{R}(k)|\mathcal{F}_{n-l_n}\}\rVert \\
     \leq& \sum_{i=0}^{\infty}\sum_{j=0}^{\infty}|\alpha_i||\alpha_j|[E\lVert E\{\bm{\varepsilon}_{n-i}(k)\bm{\varepsilon}_{n-j}^\top(k)\\
     &\hspace{10em}-E(\bm{\varepsilon}_{n-i}(k)\bm{\varepsilon}_{n-j}^\top(k))|\mathcal{F}_{n-l_n}\}\rVert].
\end{split}
\end{align}
We decompose the right-hand side of \eqref{ptMSPE.1.1.10} into three terms,
\begin{align}\label{ptMSPE.1.1.11}
\begin{split}
     & \sum_{i=0}^{\infty}\sum_{j=0}^{\infty}|\alpha_i||\alpha_j|[E\lVert E\{\bm{\varepsilon}_{n-i}(k)\bm{\varepsilon}_{n-j}^\top(k)\\
     &\hspace{10em}-E(\bm{\varepsilon}_{n-i}(k)\bm{\varepsilon}_{n-j}^\top(k))|\mathcal{F}_{n-l_n}\}\rVert] \\
     =& \sum_{i=0}^{\lfloor l_n/2\rfloor}\sum_{j=0}^{\lfloor l_n/2\rfloor}|\alpha_i||\alpha_j|[E\lVert E\{\bm{\varepsilon}_{n-i}(k)\bm{\varepsilon}_{n-j}^\top(k)\\
     &\hspace{10em}-E(\bm{\varepsilon}_{n-i}(k)\bm{\varepsilon}_{n-j}^\top(k))|\mathcal{F}_{n-l_n}\}\rVert]\\
     &+2\sum_{i=0}^{\lfloor l_n/2\rfloor}\sum_{j=\lfloor l_n/2\rfloor+1}^{\infty}|\alpha_i||\alpha_j|[E\lVert E\{\bm{\varepsilon}_{n-i}(k)\bm{\varepsilon}_{n-j}^\top(k)\\
     &\hspace{10em}-E(\bm{\varepsilon}_{n-i}(k)\bm{\varepsilon}_{n-j}^\top(k))|\mathcal{F}_{n-l_n}\}\rVert]\\
     &+\sum_{i=\lfloor l_n/2\rfloor+1}^{\infty}\sum_{j=\lfloor l_n/2\rfloor+1}^{\infty}|\alpha_i||\alpha_j|[E\lVert E\{\bm{\varepsilon}_{n-i}(k)\bm{\varepsilon}_{n-j}^\top(k)\\
     &\hspace{10em}-E(\bm{\varepsilon}_{n-i}(k)\bm{\varepsilon}_{n-j}^\top(k))|\mathcal{F}_{n-l_n}\}\rVert]\\
     :=&(E1)+2(E2)+(E3).
\end{split}
\end{align}
For $0\leq i,j\leq\lfloor l_n/2\rfloor$, we have
\begin{equation}\label{ptMSPE.1.1.12}
\begin{split}
    & E\lVert E\{\bm{\varepsilon}_{n-i}(k)\bm{\varepsilon}_{n-j}^\top(k)-E(\bm{\varepsilon}_{n-i}(k)\bm{\varepsilon}_{n-j}^\top(k))|\mathcal{F}_{n-l_n}\}\rVert \\
    \leq & E\lVert E\{\bm{\varepsilon}_{n-i}(k)\bm{\varepsilon}_{n-j}^\top(k)-E(\bm{\varepsilon}_{n-i}(k)\bm{\varepsilon}_{n-j}^\top(k))|\mathcal{F}_{n-l_n}\}\rVert_F\\
    = & E(\sum_{s=0}^{k-1}\sum_{t=0}^{k-1}[E\{\varepsilon_{n-i-s}\varepsilon_{n-j-t}-E(\varepsilon_{n-i-s}\varepsilon_{n-j-t})|\mathcal{F}_{n-l_n}\}]^2)^{\frac{1}{2}}\\
    \leq & \sum_{s=0}^{k-1}\sum_{t=0}^{k-1}E| E\{\varepsilon_{n-i-s}\varepsilon_{n-j-t}-E(\varepsilon_{n-i-s}\varepsilon_{n-j-t})|\mathcal{F}_{n-l_n}\}|,
\end{split}
\end{equation}
where $\lVert \cdot\rVert_F$ denotes the Frobenius norm. The first inequality follows from the fact that for any real matrix $\mathbf{M}$, $\lVert \mathbf{M}\rVert\leq \lVert\mathbf{M}\rVert_F$, and the second inequality follows from the fact that for any real numbers $x,y\geq0$,
\begin{equation*}
\sqrt{x+y}\leq \sqrt{x}+\sqrt{y}.
\end{equation*}
Note that for $i+s\neq j+t$ and $\max\{n-i-s,n-j-t\}>n-l_n$,
\begin{equation}\label{ptMSPE.1.1.12.1}
E\{\varepsilon_{n-i-s}\varepsilon_{n-j-t}-E(\varepsilon_{n-i-s}\varepsilon_{n-j-t})|\mathcal{F}_{n-l_n}\}=0\ \textrm{a.s.}.
\end{equation}
Moreover, it follows from \eqref{tMSPE.conditional moment} and \eqref{ln} that
\begin{equation}\label{ptMSPE.1.1.12.2}
E|E\{\varepsilon_{n-i-s}^2-E(\varepsilon_{n-i-s}^2)|\mathcal{F}_{n-l_n}\}|=o(1).
\end{equation}
Thus, by \eqref{ptMSPE.1.1.12}--\eqref{ptMSPE.1.1.12.2}, we obtain
$$E\lVert E\{\bm{\varepsilon}_{n-i}(k)\bm{\varepsilon}_{n-j}^\top(k)-E(\bm{\varepsilon}_{n-i}(k)\bm{\varepsilon}_{n-j}^\top(k))|\mathcal{F}_{n-l_n}\}\rVert=o(1).$$
Combining this with the absolute summability of $\{\alpha_i\}$ yields
\begin{equation}\label{E1}
(E1)=o(1).
\end{equation}
To address $(E2)$ and $(E3)$, note that \eqref{tMSPE.moment} ensures
$$\sup\limits_{n\geq0} E\lVert E\{\bm{\varepsilon}_{n-i}(k)\bm{\varepsilon}_{n-j}^\top(k)-E(\bm{\varepsilon}_{n-i}(k)\bm{\varepsilon}_{n-j}^\top(k))|\mathcal{F}_{n-l_n}\}\rVert<\infty,$$
for any $0\leq i,j<\infty$. This, together with \eqref{ln} and the assumption that $\alpha_i=O(i^{-\iota})$ with $\iota>3/2$, implies  
\begin{equation}\label{ptMSPE.1.1.13}
(E2)=o(1)
\end{equation}
and
\begin{equation}\label{ptMSPE.1.1.14}
(E3)=o(1).
\end{equation}
Combining \eqref{ptMSPE.1.1.10}, \eqref{ptMSPE.1.1.11} and \eqref{E1}--\eqref{ptMSPE.1.1.14} yields
\begin{equation}\label{ptMSPE.1.1.15}
E\lVert E\{\mathbf{x}_n(k)\mathbf{x}_n^\top(k)-\mathbf{R}(k)|\mathcal{F}_{n-l_n}\}\rVert =o(1).
\end{equation}
Hence, it follows from \eqref{ptMSPE.1.1.6}, \eqref{M 3n}, and \eqref{ptMSPE.1.1.15} that
\begin{equation}\label{ptMSPE.1.1.16}
(S5)=o_p(1).
\end{equation}
By an argument similar to that used in the proof of \eqref{l3.result}, one can obtain
\begin{equation}\label{ptMSPE.1.1.17}
E\lVert\tilde{\mathbf{R}}^{-1}_n(k)-\mathbf{R}^{-1}(k)\rVert^{\gamma_2}=O(n^{-\frac{\gamma_2}{2}}),
\end{equation}
for $0<\gamma_2< 3+\tilde{\delta}$. It can be derived from Lemma \ref{lB.1}(a), \eqref{ln}, \eqref{R invese 2norm}, and \eqref{ptMSPE.1.1.17} that
\begin{equation}\label{ptMSPE.1.1.18}
      (S6)=M_{1n}^\top\mathbf{R}^{-1}(k)M_{1n}+o_p(1)=Y_n+o_p(1).
\end{equation}
Finally, (D1)  is an immediate consequence
of \eqref{ptMSPE.1.1.1} and \eqref{ptMSPE.S1-S6}.\\

\noindent
\underline{Step 2}: \textit{Prove 
\begin{equation}\label{ptMSPE.1.2}
{\rm (II)}
=-2E(V_n)+o(1),
\end{equation}
where
$$V_n=
\varepsilon_{n,h,k}\mathbf{x}^\top_n(k)\mathbf{R}^{-1}(k)\{\sum_{i=k}^{n-h}\mathbf{x}_i(k)\varepsilon_{i,h,k}\}.$$}

Similar to
the proof of
\eqref{ptMSPE.1.1}, it suffices
for \eqref{ptMSPE.1.2} to show that
\begin{align*}\label{ptMSPE.1.2.0}
\begin{split}
&{\rm (C2)}: \{\tilde{W}_n\}\ \textrm{is uniformly integrable},\\ 
&{\rm (D2)}: \tilde{W}_n=o_p(1),
\end{split}
\end{align*}
where
$$\tilde{W}_n=E[\varepsilon_{n,h,k}\mathbf{x}^\top_n(k)\{\hat{\mathbf{R}}_{n,h}^{-1}(k)-\mathbf{R}^{-1}(k)\}\sum_{i=k}^{n-h}\mathbf{x}_i(k)\varepsilon_{i,h,k}
|\mathcal{F}_{n-l_n}].$$

\noindent
\underline{Step 2.1} \textit{Proof of {\rm(C2)}}\\

First, \eqref{xt moment} and the fact that
\begin{equation}\label{beta_h(k) finite}
\lVert\bm{\beta}_h(k)\rVert<\infty
\end{equation}
imply that
\begin{equation}\label{varepsilon t,h,k moment}
\sup\limits_{-\infty<t<\infty}E|\varepsilon_{t,h,k}|^{6+\delta}<\infty.
\end{equation}
Moreover, from Lemma \ref{lB.1}(c), we have for all $0<\gamma< 3+\tilde{\delta}$,
\begin{equation}\label{Rhatinv-Rinv moment}
E\lVert\hat{\mathbf{R}}^{-1}_{n,h}(k)-\mathbf{R}^{-1}(k)\rVert^{\gamma}=O(n^{-\frac{\gamma}{2}}).
\end{equation}
By \eqref{xt moment}, \eqref{sum x varepsilon moment}, \eqref{varepsilon t,h,k moment}, \eqref{Rhatinv-Rinv moment},
and an argument similar to that used in the proof of \eqref{ptMSPE.1.1.04}, (C2) follows.\\


\noindent
\underline{Step 2.2} \textit{Proof of {\rm(D2)}}\\



For all large $n$, we have
\begin{align}\label{ptMSPE.1.2.1}
\begin{split}
   \tilde{W}_n=  & E[\varepsilon_{n,h,k}\mathbf{x}^\top_n(k)\{\hat{\mathbf{R}}_{n,h}^{-1}(k)-\mathbf{R}^{-1}(k)\}\sum_{i=n-h-l_n+1}^{n-h}\mathbf{x}_i(k)\varepsilon_{i,h,k}|\mathcal{F}_{n-l_n}] \\
   &+E[\varepsilon_{n,h,k}\mathbf{x}^\top_n(k)\{\hat{\mathbf{R}}_{n,h}^{-1}(k)-\tilde{\mathbf{R}}_n^{-1}(k)\}|\mathcal{F}_{n-l_n}]\{\sum_{i=k}^{n-h-l_n}\mathbf{x}_i(k)\varepsilon_{i,h,k}\} \\
    &+E\{\varepsilon_{n,h,k}\mathbf{x}^\top_n(k)|\mathcal{F}_{n-l_n}\}\{\tilde{\mathbf{R}}_n^{-1}(k)-\mathbf{R}^{-1}(k)\}\{\sum_{i=k}^{n-h-l_n}\mathbf{x}_i(k)\varepsilon_{i,h,k}\}\\
    :=&(S7)+(S8)+(S9).
\end{split}
\end{align}
In the following, we will prove that
\begin{equation}\label{ptMSPE.S7-S9}
    (S7)=o_p(1),\ (S8)=o_p(1),\ (S9)=o_p(1).
\end{equation}
By \eqref{ln}, \eqref{xt moment}, \eqref{sum x varepsilon moment}, \eqref{varepsilon t,h,k moment}, \eqref{Rhatinv-Rinv moment} and an argument similar to that used in the proof of \eqref{ptMSPE.1.1.3}, we have
\begin{equation}\label{ptMSPE.1.2.2}
\begin{split}
    (S7)=O_p\{(l_n/n)^{1/2}\}=o_p(1).
\end{split}
\end{equation}
Likewise, \eqref{ln}, \eqref{xt moment}, \eqref{sum x varepsilon moment}, \eqref{ptMSPE.1.1.7} and \eqref{varepsilon t,h,k moment} lead to
\begin{equation}\label{ptMSPE.1.2.3}
(S8)=O_p(l_n/n^{1/2})=o_p(1).
\end{equation}
As for $(S9)$, note that \eqref{AR mean part}, \eqref{AR(k) misspecified multi}, and Minkowski's inequality yield
\begin{equation}\label{ptMSPE.1.2.3.1}
\begin{split}
     & E\lVert E\{\mathbf{x}_n(k)\varepsilon_{n,h,k}|\mathcal{F}_{n-l_n}\}\rVert\\
     =&E\lVert E\{\mathbf{x}_n(k)\varepsilon_{n,h,k}-E(\mathbf{x}_n(k)\varepsilon_{n,h,k})|\mathcal{F}_{n-l_n}\}\rVert\\
    \leq & \sum_{i=0}^{\infty}\sum_{j=0}^{\infty}|\alpha_i||\alpha_j|E\lVert E[\bm{\varepsilon}_{n-i}(k)(\varepsilon_{n+h-j}-\bm{\beta}_j(k)\bm{\varepsilon}_{n-j}(k))\\
    &-E\{\bm{\varepsilon}_{n-i}(k)(\varepsilon_{n+h-j}-\bm{\beta}_j(k)\bm{\varepsilon}_{n-j}(k))\}|\mathcal{F}_{n-l_n}]\rVert.
\end{split}
\end{equation}
Similar to the proof of \eqref{ptMSPE.1.1.15}, the absolute summability of $\{\alpha_i\}$, along with \eqref{ptMSPE.1.1.12.1}, \eqref{ptMSPE.1.1.12.2}, and \eqref{ptMSPE.1.2.3.1}, leads to
\begin{equation}\label{ptMSPE.1.2.4}
E\lVert E\{\mathbf{x}_n(k)\varepsilon_{n,h,k}|\mathcal{F}_{n-l_n}\}\rVert=o(1).
\end{equation}
Hence, it follows from \eqref{M 3n}, \eqref{ptMSPE.1.1.17} and \eqref{ptMSPE.1.2.4} that
\begin{equation}\label{ptMSPE.1.2.5}
   (S9)=o_p(1).
\end{equation}
Consequently, (D2) is ensured by \eqref{ptMSPE.1.2.1} and \eqref{ptMSPE.S7-S9}.\\

\noindent
\underline{Step 3}: \textit{Prove 
\begin{equation}\label{ptMSPE.step 3}
\begin{split}
    & E(Y_n)-2E(V_n) \\
    = & \mathrm{tr}\{\mathbf{R}^{-1}(k)\mathbf{L}_{0,h}(k)\}+2\sum_{s=1}^{h-1}\mathrm{tr}\{\mathbf{R}^{-1}(k)\mathbf{L}_{s,h}(k)\}+o(1).
    \end{split}
\end{equation}
}
The fourth-order stationarity of $\{\varepsilon_t\}$ implies
\begin{equation}\label{ptMSPE.2}
\begin{split}
   E(V_n) % &  \sum_{i=h}^{n-k}E\{\mathbf{x}^\top_k(k)\mathbf{R}^{-1}(k)\mathbf{x}_{k+i}(k)\varepsilon_{k,h,k}\varepsilon_{k+i,h,k}\}\\
     =& \sum_{i=h}^{n-h-k}E\{\mathbf{x}^\top_k(k)\mathbf{R}^{-1}(k)\mathbf{x}_{k+i}(k)\varepsilon_{k,h,k}\varepsilon_{k+i,h,k}\}\\
     &+\sum_{i=n-h-k+1}^{n-k}E\{\mathbf{x}^\top_k(k)\mathbf{R}^{-1}(k)\mathbf{x}_{k+i}(k)\varepsilon_{k,h,k}\varepsilon_{k+i,h,k}\},
\end{split}
\end{equation}
and
\begin{equation}\label{ptMSPE.3}
\begin{split}
     %& %E[\{\frac{1}{\sqrt{n-h-k+1}}\sum_{i=k}^{n-h}\mathbf{x}^\top_i(k)\varepsilon_{i,h,k}\}\mathbf{R}^{-1}(k)\{\frac{1}{\sqrt{n-h-k+1}}\sum_{i=k}^{n-h}\mathbf{x}_i(k)\varepsilon_{i,h,k}\}]\\
  % =&\frac{1}{n-h-k+1}\sum_{i=k}^{n-h}E\{\mathbf{x}_i^\top(k)\mathbf{R}^{-1}(k)\mathbf{x}_i(k)\varepsilon_{i,h,k}^2\}\\
   %&+\frac{2}{n-h-k+1}\sum_{i=k}^{n-h-1}\sum_{j=k+1}^{n-h}E\{\mathbf{x}_i^\top(k)\mathbf{R}^{-1}(k)\mathbf{x}_j(k)\varepsilon_{i,h,k}\varepsilon_{j,h,k}\}\\
   %=&\mathrm{tr}\{\mathbf{R}^{-1}(k)\mathbf{L}_{0,h}(k)\}+\frac{2}{n-h-k+1}\sum_{i=1}^{n-h-k}(n-h-k+1-i)E\{\mathbf{x}_k^\top(k)\mathbf{R}^{-1}(k)\mathbf{x}_{k+i}(k)\varepsilon_{k,h,k}\varepsilon_{k+i,h,k}\}\\
   &E(Y_n)\\
   =&\mathrm{tr}\{\mathbf{R}^{-1}(k)\mathbf{L}_{0,h}(k)\}+2\sum_{s=1}^{h-1}\mathrm{tr}\{\mathbf{R}^{-1}(k)\mathbf{L}_{s,h}(k)\}\\
   &+2\sum_{i=h}^{n-h-k}E\{\mathbf{x}_k^\top(k)\mathbf{R}^{-1}(k)\mathbf{x}_{k+i}(k)\varepsilon_{k,h,k}\varepsilon_{k+i,h,k}\}\\
   &-\frac{2}{n-h-k+1}\sum_{i=1}^{n-h-k}iE\{\mathbf{x}_k^\top(k)\mathbf{R}^{-1}(k)\mathbf{x}_{k+i}(k)\varepsilon_{k,h,k}\varepsilon_{k+i,h,k}\}.
\end{split}
\end{equation}
In view of \eqref{ptMSPE.2} and \eqref{ptMSPE.3}, it remains to prove that
\begin{equation}\label{ptMSPE.9}
E\{\mathbf{x}_k^\top(k)\mathbf{R}^{-1}(k)\mathbf{x}_{k+n}(k)\varepsilon_{k,h,k}\varepsilon_{k+n,h,k}\}=o(n^{-1}).
\end{equation}
Let $\varepsilon^\ast_{j,h,k}=\varepsilon_{j+h}-\bm{\beta}_h^\top(k)\bm{\varepsilon}_{j}(k)$. By \eqref{AR mean part} and \eqref{AR(k) misspecified multi}, we have for any $0\leq u,w\leq k-1$ and $n\geq4(k+h)$,
\begin{align}\label{ptMSPE.4}
\begin{split}
     & E(x_{k-u}x_{k+n-w}\varepsilon_{k,h,k}\varepsilon_{k+n,h,k}) \\
     =& E[x_{k-u}\varepsilon_{k,h,k}\{x_{k+n-w}\varepsilon_{k+n,h,k}-E(x_{k+n-w}\varepsilon_{k+n,h,k})\}]\\
     =& E[x_{k-u}\varepsilon_{k,h,k}\{\sum_{j=0}^{\lfloor n/2\rfloor}\sum_{l=0}^{\lfloor n/2\rfloor}\alpha_j\alpha_l(\varepsilon_{k+n-w-j}\varepsilon^\ast_{k+n-l,h,k}\\
     &\hspace{15em}-E(\varepsilon_{k+n-w-j}\varepsilon^\ast_{k+n-l,h,k}))\}]\\
     &+E[x_{k-u}\varepsilon_{k,h,k}\{\sum_{j=\lfloor n/2\rfloor+1}^{\infty}\sum_{l=\lfloor n/4\rfloor+1}^{\infty}\alpha_j\alpha_l(\varepsilon_{k+n-w-j}\varepsilon^\ast_{k+n-l,h,k}\\
      &\hspace{15em}-E(\varepsilon_{k+n-w-j}\varepsilon^\ast_{k+n-l,h,k}))\}]\\
     &+E[x_{k-u}\varepsilon_{k,h,k}\{\sum_{j=\lfloor n/4\rfloor+1}^{\lfloor n/2\rfloor}\sum_{l=\lfloor n/2\rfloor+1}^{\infty}\alpha_j\alpha_l(\varepsilon_{k+n-w-j}\varepsilon^\ast_{k+n-l,h,k}\\
      &\hspace{15em}-E(\varepsilon_{k+n-w-j}\varepsilon^\ast_{k+n-l,h,k}))\}]\\
     :=&(E4)+(E5)+(E6),
\end{split}
\end{align}
where the second equality follows from the fact that for $s>t\vee (k+h)$ or $t>s\vee (k+h)$,
\begin{equation}\label{ptMSPE.4.1}
E[x_{k-u}\varepsilon_{k,h,k}\{\varepsilon_{s}\varepsilon_{t}-E(\varepsilon_{s}\varepsilon_{t})\}]=0.
\end{equation}
Note that for all sufficiently large $n$ and $0\leq j\leq \lfloor n/2\rfloor$,
\begin{equation}\label{ptMSPE.5}
\begin{split}
     & E[x_{k-u}\varepsilon_{k,h,k}\{\varepsilon_{k+n-w-j}^2-E(\varepsilon_{k+n-w-j}^2)\}] \\
    = & E[x_{k-u}\varepsilon_{k,h,k}E\{\varepsilon_{k+n-w-j}^2-E(\varepsilon_{k+n-w-j}^2)|\mathcal{F}_{k+n-w-\lfloor\frac{3}{4}n\rfloor}\}].
\end{split}
\end{equation}
This, \eqref{tMSPE.conditional moment}, \eqref{xt moment}, \eqref{beta_h(k) finite}, \eqref{varepsilon t,h,k moment}, \eqref{ptMSPE.4.1}, Minkowski's inequality and H$\ddot{\mathrm{o}}$lder's inequality yield
\begin{equation}\label{ptMSPE.6}
(E4)=o(n^{-1}).
\end{equation}
As for $(E5)$, note that it can be deduced from \eqref{xt moment} and \eqref{tMSPE.moment} that
\begin{equation}\label{E5.1}
\begin{split}
     &\sup\limits_{j,l\geq0}E[x_{k-u}\varepsilon_{k,h,k}\{(\varepsilon_{k+n-w-j}\varepsilon^\ast_{k+n-l,h,k}-E(\varepsilon_{k+n-w-j}\varepsilon^\ast_{k+n-l,h,k})\}] \\
   =  & O(1).
\end{split}
\end{equation}
Thus, we have
\begin{equation}\label{ptMSPE.7}
(E5)\leq C\sum_{j=\lfloor n/2\rfloor+1}^{\infty}\sum_{l=\lfloor n/4\rfloor+1}^{\infty}|\alpha_j||\alpha_l|\leq C(n^{-\iota+1})^2=Cn^{-2\iota+2}=o(n^{-1}),
\end{equation}
where the first inequality follows from \eqref{E5.1} and Minkowski's inequality, and the last equality is ensured by the assumption that $\alpha_i=O(i^{-\iota})$ with $\iota>3/2$. Similarly, one can obtain
\begin{equation}\label{ptMSPE.8}
(E6)=o(n^{-1}).
\end{equation}
Using \eqref{R invese 2norm}, \eqref{ptMSPE.4}, \eqref{ptMSPE.6}--\eqref{ptMSPE.8} and Minkowski's inequality, we obtain \eqref{ptMSPE.9}.
}
\end{proof}

\begin{proof}[Proof of Proposition \ref{prop3.1}]
By \eqref{h.varepsilon square}--\eqref{wt moment}, Minkowski's inequality and Burkholder's inequality, {we have for any $1 \leq \ubar{m}_{n}\leq \bar{m}_{n} \leq n$,
\begin{equation}\label{pp1.1}
\begin{split}
     & \sup\limits_{-\infty<s<\infty}E|\sum_{t=\ubar{m}_{n}}^{\bar{m}_{n}}(\varepsilon_{t+s}^2-E(\varepsilon_{t+s}^2))|^{3+\delta^\star} \\
     =&  \sup\limits_{-\infty<s<\infty}E|\sum_{t=\ubar{m}_{n}}^{\bar{m}_{n}}\sum_{u=0}^{\infty}\mathbf{a}^\top_u\bm{w}_{t+s,u}|^{3+\delta^\star}\\
     \leq& \sup\limits_{-\infty<s<\infty}\{\sum_{u=0}^{\infty}\lVert\mathbf{a}_u\rVert(E\lVert
     \sum_{t=\ubar{m}_{n}}^{\bar{m}_{n}}\bm{w}_{t+s,u}\rVert^{3+\delta^\star})^{\frac{1}{3+\delta^\star}}\}^{3+\delta^\star}\\
     \leq & C(\sum_{u=0}^{\infty}\lVert\mathbf{a}_u\rVert)^{3+\delta^\star}(\bar{m}_{n}-\ubar{m}_{n}+1)^{\frac{3+\delta^\star}{2}}\leq C(\bar{m}_{n}-\ubar{m}_{n}+1)^{\frac{3+\delta^\star}{2}}.
\end{split}
\end{equation}}
Thus, \eqref{tMSPE.moment2} holds. We next prove \eqref{tMSPE.conditional moment}. Observe that by Minkowski's inequality, Burkholder's inequality and \eqref{h.varepsilon square}--\eqref{wt moment},
\begin{equation}\label{pp1.2}
\begin{split}
     & \sup\limits_{-\infty<t<\infty} E|E(\varepsilon_t^2|\mathcal{F}_{t-j})-E(\varepsilon_t^2)|^2 =  \sup\limits_{-\infty<t<\infty} E|E(\sum_{s=0}^{\infty}\mathbf{a}^\top_s\bm{w}_{t,s}|\mathcal{F}_{t-j})|^2\\
    =&\sup\limits_{-\infty<t<\infty} E|\sum_{s=j}^{\infty}\mathbf{a}^\top_s\bm{w}_{t,s}|^2\leq  \sup\limits_{-\infty<t<\infty}\{\sum_{s=j}^{\infty}\lVert\mathbf{a}_s\rVert^2(E\lVert\bm{w}_{t,s}\rVert^2)\}\\
    \leq& C(\sum_{s=j}^{\infty}\lVert\mathbf{a}_s\rVert^2)=o(j^{-2}),
\end{split}
\end{equation}
as $j\to\infty$. Hence, \eqref{tMSPE.conditional moment} follows from \eqref{pp1.2} and H$\ddot{\mathrm{o}}$lder's inequality.
\end{proof}

\section{Proof of Theorem \ref{t MRIC}}\label{appC}
%\setcounter{equation}{0}
%\renewcommand{\theequation}{B.\arabic{equation}}

%To prove \eqref{t MRIC.result},
It suffices to show \eqref{ing129} and \eqref{ing128}, which are
ensured by \eqref{pt.MRIC.1}--\eqref{pt.MRIC.5}:
{for any finite set $J\subset\mathbb{N}$ and $0\leq s\leq h-1$,
\begin{equation}\label{pt.MRIC.1}
\frac{1}{n-h-\tilde{d}+1}\sum_{t=\tilde{d}}^{n-h}\varepsilon_{t,h,J}^2-E(\varepsilon_{t,h,J}^2)=O_p(n^{-1/2}),
\end{equation}
\begin{equation}\label{pt.MRIC.2}
\frac{1}{n-h-\tilde{d}-s+1}\sum_{t=\tilde{d}}^{n-h-s}\mathbf{x}_{t}(J)\mathbf{x}_{t+s}^\top(J)\varepsilon_{t,h,J}\varepsilon_{t+s,h,J}-L_{s,h}(J)=o_p(1),
\end{equation}
\begin{equation}\label{pt.MRIC.3}
\frac{1}{\sqrt{n-h-\tilde{d}+1}}\sum_{t=\tilde{d}}^{n-h}\mathbf{x}_{t}(J)\varepsilon_{t,h,J}=O_p(1),
\end{equation}
\begin{equation}\label{pt.MRIC.4}
\frac{1}{n-h-\tilde{d}+1}\sum_{t=\tilde{d}}^{n-h}\mathbf{x}_{t}(J)\mathbf{x}_{t}^\top(J)-\mathbf{R}(J)=o_p(1),
\end{equation}
and
\begin{equation}\label{pt.MRIC.5}
\sup\limits_{-\infty<t<\infty}E|\varepsilon_{t,h,J}|^4+\sup\limits_{-\infty<t<\infty}E\lVert\mathbf{x}_t(J)\rVert^4<\infty.
\end{equation}}
Whereas \eqref{ing129} follows from \eqref{pt.MRIC.2}--\eqref{pt.MRIC.5},
\eqref{ing128} is guaranteed by \eqref{pt.MRIC.1}, \eqref{pt.MRIC.3}, and \eqref{pt.MRIC.4}.

{We first prove \eqref{pt.MRIC.1}. Observe that from \eqref{xt+h J}, we have
\begin{align}\label{pt.MRIC.1.1}
\begin{split}
     &  \sum_{t=\tilde{d}}^{n-h}\varepsilon_{t,h,J}^2-E(\varepsilon_{t,h,J}^2) \\
   =  &  \sum_{t=\tilde{d}}^{n-h}(x_{t+h}-\bm{\beta}_h^\top(J)\mathbf{x}_{t}(J))^2-E(x_{t+h}-\bm{\beta}_h^\top(J)\mathbf{x}_{t}(J))^2.
     \end{split}
\end{align}
In view of \eqref{pt.MRIC.1.1}, \eqref{pt.MRIC.1} holds if
\begin{equation}\label{pt.MRIC.1.2}
\sup\limits_{-\infty<s_1,s_2<\infty}E|\sum_{t=\tilde{d}}^{n-h}(x_{t+s_1}x_{t+s_2}-E(x_{t+s_1}x_{t+s_2}))|^2\leq Cn,
\end{equation}
which, in turn, is satisfied by the absolute summability of $\{\alpha_i\}$ and \eqref{pt.MRIC.1.3}. Similarly, \eqref{pt.MRIC.3} and \eqref{pt.MRIC.4} can be concluded from \eqref{pt.MRIC.1.3}. In addition, \eqref{tMRIC.moment}, combined with the arguments used in the proofs of \eqref{xt moment} and \eqref{varepsilon t,h,k moment}, leads to \eqref{pt.MRIC.5}.}

{It remains to prove \eqref{pt.MRIC.2}. Note that for any finite set $J\subset\mathbb{N}$ and any $\tilde{d}\leq t \leq n-h-s$,
\begin{align}\label{pt.MRIC.2.1}
\begin{split}
    &\mathbf{x}_{t}(J)\mathbf{x}_{t+s}^\top(J)\varepsilon_{t,h,J}\varepsilon_{t+s,h,J}\\
    = & \big\{\sum_{i=0}^{\infty}\alpha_i\bm{\varepsilon}_{t-i}(J)\big\}
    \big\{\sum_{j=0}^{\infty}\alpha_j\bm{\varepsilon}_{t+s-j}^\top(J)\big\}\\
    &\times
    \left[\sum_{u=0}^{\infty}\alpha_u\{\varepsilon_{t+h-u}-\bm{\beta}_h^\top(J)\bm{\varepsilon}_{t-u}(J)\}\right]\\
    &\times
    \left[\sum_{v=0}^{\infty}\alpha_v\{\varepsilon_{t+s+h-v}-\bm{\beta}_h^\top(J)\bm{\varepsilon}_{t+s-v}(J)\}\right].
    \end{split}
\end{align}
In light of the absolute summability of $\{\alpha_i\}$ and %\eqref{alpha rate}, \eqref{pt.MRIC.1.5},
\eqref{pt.MRIC.2.1}, equation \eqref{pt.MRIC.2} follows if, for any $0\leq s\leq h-1$, the following holds:
\begin{equation*}
\sup\limits_{-\infty< j_1,j_2,j_3,j_4<\infty}E|\sum_{t=\tilde{d}}^{n-h-s}(\varepsilon_{t+j_1}\varepsilon_{t+j_2}\varepsilon_{t+j_3}\varepsilon_{t+j_4}-E(\varepsilon_{t+j_1}\varepsilon_{t+j_2}\varepsilon_{t+j_3}\varepsilon_{t+j_4}))|=o(n),
\end{equation*}
which, in turn, is ensured by \eqref{pt.MRIC.2.6}. The proof of Theorem \ref{t MRIC} is now complete.}

\section{Proofs of Propositions \ref{propA.1} and \ref{propA.2}}\label{supp3.1}
{
\begin{proof}[Proof of Proposition \ref{propA.1}]
Under \eqref{propA.1.1}--\eqref{propA.1.2}, \eqref{tMSPE.moment2} with $\delta^\star=\delta/2$ directly follows from (35) of \cite{Xiao2011}. To prove \eqref{tMSPE.conditional moment}, note that by Minkowski's inequality, we have
\begin{align}\label{pp3.2.1}
\begin{split}
     & \sup\limits_{-\infty<t<\infty} E|E(\varepsilon_t^2|\mathcal{F}_{t-j})-E(\varepsilon_t^2)|^{3/2} \\
   =  & \sup\limits_{-\infty<t<\infty} E|\sum_{i=0}^{\infty} E\{\varepsilon_t^2-E(\varepsilon_t^2)|\mathcal{F}_{t-j-i}\}\\
   &\hspace{10em}-E\{\varepsilon_t^2-E(\varepsilon_t^2)|\mathcal{F}_{t-j-i-1}\}|^{3/2}\\
   \leq & \sup\limits_{-\infty<t<\infty} (\sum_{i=0}^{\infty}[E| E\{\varepsilon_t^2-E(\varepsilon_t^2)|\mathcal{F}_{t-j-i}\}\\
   &\hspace{10em}-E\{\varepsilon_t^2-E(\varepsilon_t^2)|\mathcal{F}_{t-j-i-1}\}|^{3/2}]^{2/3})^{3/2}.
   \end{split}
\end{align}
Moreover, by the strict stationarity of $\varepsilon_t$, \eqref{propA.1.3}, and (i) and (ii) of Theorem 1 of \cite{Wu2005Main}, it follows that
\begin{equation}\label{pp3.2.2}
\begin{split}
     & [E| E\{\varepsilon_t^2-E(\varepsilon_t^2)|\mathcal{F}_{t-j-i}\}-E\{\varepsilon_t^2-E(\varepsilon_t^2)|\mathcal{F}_{t-j-i-1}\}|^{3/2}]^{2/3} \\
    = & [E| E\{\varepsilon_{j+i}^2-E(\varepsilon_{j+i}^2)|\mathcal{F}_{0}\}-E\{\varepsilon_{j+i}^2-E(\varepsilon_{j+i}^2)|\mathcal{F}_{-1}\}|^{3/2}]^{2/3}\\
    \leq & (E|\varepsilon_{j+i}^2-{\varepsilon^\star_{j+i}}^2|^{3/2})^{2/3}\\
    \leq & (E|\varepsilon_{j+i}-\varepsilon^\star_{j+i}|^{3})^{1/3}(E|\varepsilon_{j+i}+\varepsilon^\star_{j+i}|^{3})^{1/3}\\
    \leq & C\delta^\prime_{3}(j+i)\leq C\delta^\prime_{6+\delta}(j+i).
\end{split}
\end{equation}
Combining \eqref{propA.1.2}, \eqref{pp3.2.1} and \eqref{pp3.2.2} gives
\begin{equation}\label{pp3.2.3}
\begin{split}
      & \sup\limits_{-\infty<t<\infty} E|E(\varepsilon_t^2|\mathcal{F}_{t-j})-E(\varepsilon_t^2)|^{3/2}\\
    \leq& \{\sum_{i=0}^{\infty}\delta^\prime_{6+\delta}(j+i)\}^{3/2}=\{\sum_{i=j}^{\infty}\delta^\prime_{6+\delta}(i)\}^{3/2} =o(j^{-3/2}),
\end{split}
\end{equation}
as $j\to\infty$, and the proof is complete.
\end{proof}

\begin{proof}[Proof of Proposition \ref{propA.2}]
Under \eqref{tMRIC.moment} and \eqref{propA.2.3}, \eqref{tMSPE.moment2} with $\delta^\star=\delta/2$ follows from Theorem 1 of \cite{Kim1993}. Furthermore, by combining \eqref{tMRIC.moment}, \eqref{propA.2.3}, and Theorem 14.2 of \cite{Davidson1994}, we obtain
\begin{equation}\label{ppA.2.1}
\begin{split}
    & \sup\limits_{-\infty<t<\infty} E|E(\varepsilon_t^2|\mathcal{F}_{t-j})-E(\varepsilon_t^2)|^{3/2} \\
   \leq  & \{2(2^{2/3}+1)\}^{3/2}\alpha(j)^{1-\frac{3/2}{(6+\delta)/2}}(\sup\limits_{-\infty<t<\infty}E|\varepsilon_t|^{6+\delta})^{\frac{3/2}{6+\delta}} \\
    = & o(j^{-3/2}),
\end{split}
\end{equation}
as $j\to\infty$. Thus, \eqref{tMSPE.conditional moment} holds.
\end{proof}}

\section{{Proofs of Propositions \ref{propB.1} and \ref{propB.2}}}\label{Ex3-S}
{
\begin{proof}[Proof of Proposition \ref{propB.1}]
Let $\tilde{\delta}=\delta/4$. Given any $-\infty<t_1,t_2,t_3,t_4<\infty$, we have
\begin{align}\label{p.pr3.2.1}
\begin{split}
     & (E|\varepsilon_{t_1}\varepsilon_{t_2}\varepsilon_{t_3}\varepsilon_{t_4}-\varepsilon^\star_{t_1}\varepsilon^\star_{t_2}\varepsilon^\star_{t_3}\varepsilon^\star_{t_4}|^{1+\tilde{\delta}})^{\frac{1}{1+\tilde{\delta}}} \\
  \leq & (E|\varepsilon_{t_1}\varepsilon_{t_2}\varepsilon_{t_3}(\varepsilon_{t_4}-\varepsilon^\star_{t_4})|^{1+\tilde{\delta}})^{\frac{1}{1+\tilde{\delta}}}+
  (E|\varepsilon_{t_1}\varepsilon_{t_2}(\varepsilon_{t_3}-\varepsilon^\star_{t_3})\varepsilon^\star_{t_4}|^{1+\tilde{\delta}})^{\frac{1}{1+\tilde{\delta}}}\\
  &+(E|\varepsilon_{t_1}(\varepsilon_{t_2}-\varepsilon^\star_{t_2})\varepsilon^\star_{t_3}\varepsilon^\star_{t_4}|^{1+\tilde{\delta}})^{\frac{1}{1+\tilde{\delta}}}+
  (E|(\varepsilon_{t_1}-\varepsilon^\star_{t_1})\varepsilon^\star_{t_2}\varepsilon^\star_{t_3}\varepsilon^\star_{t_4}|^{1+\tilde{\delta}})^{\frac{1}{1+\tilde{\delta}}}\\
  \leq& (\sup\limits_{-\infty<t<\infty}E|\varepsilon_t|^{4+\delta})^{\frac{3}{4+\delta}}(\delta^\prime_{4+\delta}(t_1)+\delta^\prime_{4+\delta}(t_2)+\delta^\prime_{4+\delta}(t_3)+\delta^\prime_{4+\delta}(t_4)),
\end{split}
\end{align}
where the first and last inequalities follow from Minkowski's inequality and H$\ddot{\mathrm{o}}$lder's inequality, respectively. Utilizing \eqref{propB.1.1}, \eqref{propB.1.2}, and \eqref{p.pr3.2.1}, we can apply (33) of \cite{Xiao2011} to derive \eqref{pt.MRIC.2.6}. In addition, \eqref{pt.MRIC.1.3} can be obtained in the same manner.
\end{proof}

\begin{proof}[Proof of Proposition \ref{propB.2}]
Assumption \eqref{propB.2.1} ensures that $$\{\varepsilon_{t+s_1}\varepsilon_{t+s_2}\varepsilon_{t+s_3}\varepsilon_{t+s_4}-E(\varepsilon_{t+s_1}\varepsilon_{t+s_2}\varepsilon_{t+s_3}\varepsilon_{t+s_4})\}$$ is uniformly integrable for any $-\infty<s_1,s_2,s_3,s_4<\infty$, thereby guaranteeing \eqref{pt.MRIC.2.6} through Theorem 2 of \cite{Andrews1988M}. %see also Theorem 1 of \cite{Hansen2019}. 
Additionally, assumptions \eqref{propB.2.1} and \eqref{App.2} ensure that \eqref{pt.MRIC.1.3} holds  through Theorem 5 of \cite{Kim1994M}.%, where \eqref{propB.2.1} also facilitates the use of arguments based on H$\ddot{\mathrm{o}}$lder's inequality in its proof. \it{(The last sentence is unclear!)}
\end{proof}}

%\begin{remark}\label{rm.s7.2.1}
%The requirement for strong mixing processes to have a $4+\delta$-moment, rather than just a fourth moment as in processes measured by functional dependence in Proposition \ref{prop3.2}, arises because, in the absence of i.i.d. and strict stationarity assumptions, ensuring conditions \eqref{pt.MRIC.1.3}--\eqref{pt.MRIC.2.6} inevitably relies on arguments based on uniform integrability and H$\ddot{\mathrm{o}}$lder's inequality.
%\end{remark}}

\section{Additional Simulation Studies}\label{sec:add sim}
\subsection{Additional Results for Examples 4.1 and 4.2}\label{sec:add sim ex4.1 4.2}
In this section, we investigate the performance of the ratio $R_{n,h}=g_{n,h}(k)/g_{h}(k)$ by considering alternative AR and MA coefficients for the models presented in Examples \ref{ex4.1} and \ref{ex4.2}. As previously defined, $g_{h}(k)$ denotes the theoretical second-order MSPE, and $g_{n,h}(k)$ is the empirical approximation of \eqref{ing2}.

\begin{example}\label{exS.1}
We generate $M=5000$ realizations from the following AR(2) model,
\begin{equation*}%\label{exS.1.1}
x_t=\beta_2 x_{t-2}+\varepsilon_t,
\end{equation*}
where $\beta_2 \in \{0.5, 0.8, -0.8\}$, and 
$\varepsilon_t$ obeys
\eqref{varepsilon basic form} with $\{z_t\}$ being a sequence of
i.i.d. $N(0,1)$ random variables
and $\{\sigma_t\}$ being either a GARCH(1,1) process,
\begin{equation*}
\sigma_t^2=0.4+0.2\varepsilon_{t-1}^2+0.55\sigma_{t-1}^2,
\end{equation*}
or a SV(1) (\eqref{SV model} with $\tilde{p}=1$) process,
\begin{equation*}
(1-0.98B)\log(\sigma_t^2)=0.01+v_t,
\end{equation*}
in which
$\{v_t\}$ is a sequence of i.i.d. $N(0,0.04)$ random variables. This setup mirrors Example \ref{ex4.1}, with the exception that the AR coefficient $\beta_2$ is varied among $\{0.5, 0.8, -0.8\}$. The corresponding values of $g_{h}(k)$ and $R_{n,h}$ are summarized in Table \ref{tabS.1}.
\end{example}

\begin{example}\label{exS.2}
We generate $M=5000$ realizations from
the MA(1) model,
\begin{eqnarray*}
x_t=\varepsilon_t+\alpha_1\varepsilon_{t-1},
\end{eqnarray*}
where $\alpha_1 \in \{0.5, -0.5, 0.8\}$. The remaining settings are identical to those described in Example \ref{exS.1}. This example revisits Example \ref{ex4.2} by varying the MA coefficient $\alpha_1$ across $\{0.5, -0.5, 0.8\}$. The associated results are reported in Table \ref{tabS.2}.
\end{example}

According to Tables \ref{tabS.1} and \ref{tabS.2}, when $n=2000$, the ratio $R_{n,h}$ lies within the interval $(0.9, 1.1)$ in most scenarios. 
Notable exceptions occur in the AR(2) model when $\beta_2=0.8$ at $h=1$ (for both GARCH(1,1) and SV(1) errors) and when $\beta_2=-0.8$ at $h=5$ (only for SV(1) errors). These findings further validate the asymptotic results established in Theorem \ref{tMSPE}. Furthermore, the results appear generally robust to the signs and magnitudes of the AR and MA coefficients ($\beta_2$ and $\alpha_1$).

\begin{table}
\caption{The values of $g_{h}(k)$ and $R_{n,h}$, with $k=1$, $n=500, 2000$, and $h=1,\ldots, 5$, in Example \ref{exS.1}}
\label{tabS.1}
\centering
\begin{tabular}{ccccccccccc} \hline
    GARCH(1,1),\ $\beta_2=0.5$& $h=1$ & 2 & 3 & 4 & 5 \\ \hline
     $g_{h}(1)$ &3.454 & 2.425 & 5.487 & 3.702 & 4.850 \\ \hline
   $R_{500,h}$ &1.821 & 0.902 & 1.072 & 0.936 & 1.065 \\
   $R_{2000,h}$ &0.958 & 0.975 & 1.048 & 0.913 & 0.912 \\ \hline 
      $\beta_2=0.8$& $h=1$ & 2 & 3 & 4 & 5 \\ \hline
    $g_{h}(1)$ &5.327 & 2.132 & 17.786 & 5.694 & 23.364 \\ \hline
   $R_{500,h}$ &1.961 & 0.950 & 1.001 & 0.996 & 1.041 \\
   $R_{2000,h}$ &1.228 & 0.971 & 0.936 & 0.952 & 0.958 \\ \hline 
        $\beta_2=-0.8$& $h=1$ & 2 & 3 & 4 & 5 \\ \hline
    $g_{h}(1)$ &5.327 & 2.132 & 6.876 & 5.694 & 9.399 \\ \hline
   $R_{500,h}$ &1.063 & 0.934 & 0.860 & 0.931 & 0.821 \\
   $R_{2000,h}$ &0.946 & 1.078 & 1.018 & 1.024 & 0.925 \\ \hline \hline
   SV(1),\ $\beta_2=0.5$ & $h=1$ & 2 & 3 & 4 & 5 \\ \hline
    $g_{h}(1)$ & 9.680 & 7.119 & 16.526 & 12.061 & 15.966 \\ \hline
   $R_{500,h}$ & 1.080 & 0.917 & 0.581 & 0.775 & 0.992 \\
   $R_{2000,h}$ & 0.946 & 0.984 & 1.063 & 0.947 & 1.099 \\ \hline
      $\beta_2=0.8$ & $h=1$ & 2 & 3 & 4 & 5 \\ \hline
    $g_{h}(1)$ & 19.188 & 6.780 & 62.828 & 19.251 & 82.631 \\ \hline
   $R_{500,h}$ & 0.815 & 0.887 & 0.588 & 0.843 & 0.298 \\
   $R_{2000,h}$ & 1.132 & 1.003 & 0.900 & 0.928 & 1.060 \\ \hline
      $\beta_2=-0.8$ & $h=1$ & 2 & 3 & 4 & 5 \\ \hline
    $g_{h}(1)$ & 19.188 & 6.780 & 23.531 & 19.251 & 32.331 \\ \hline
   $R_{500,h}$ & 0.918 & 0.724 & 1.004 & 0.785 & 0.755 \\
   $R_{2000,h}$ & 0.978 & 1.011 & 1.026 & 0.985 & 0.888 \\ \hline
  \end{tabular}
\end{table}

\begin{table}
\caption{The values of $g_{h}(k)$ and $R_{n,h}$, with $k=1$, $n=500, 2000$, and $h=1,\ldots, 5$, in Example \ref{exS.2}}
\label{tabS.2}
\centering
\begin{tabular}{ccccccccccc} \hline
    GARCH(1,1),\ $\alpha_1=0.5$& $h=1$ & 2 & 3 & 4 & 5 \\ \hline
     $g_{h}(1)$ &2.698 & 3.955 & 3.626 & 3.380 & 3.195 \\ \hline
   $R_{500,h}$ &1.102 & 0.953 & 0.954 & 0.922 & 0.996 \\
   $R_{2000,h}$ &1.078 & 1.023 & 1.036 & 0.946 & 0.964 \\ \hline 
      $\alpha_1=-0.5$& $h=1$ & 2 & 3 & 4 & 5 \\ \hline
    $g_{h}(1)$ &2.698 & 3.955 & 3.626 & 3.380 & 3.195 \\ \hline
   $R_{500,h}$ &0.772 & 0.990 & 0.984 & 0.965 & 0.948 \\
   $R_{2000,h}$ &0.994 & 0.994 & 0.983 & 0.966 & 0.979 \\ \hline 
        $\alpha_1=0.8$& $h=1$ & 2 & 3 & 4 & 5 \\ \hline
    $g_{h}(1)$ & 2.964 & 5.808 & 5.324 & 4.961 & 4.689 \\ \hline
   $R_{500,h}$ &1.207 & 0.987 & 0.946 & 0.908 & 0.988 \\
   $R_{2000,h}$ &1.033 & 1.034 & 1.034 & 0.938 & 0.936 \\ \hline \hline
   SV(1),\ $\alpha_1=0.5$ & $h=1$ & 2 & 3 & 4 & 5 \\ \hline
    $g_{h}(1)$ & 7.469 & 11.894 & 11.665 & 11.446 & 11.234 \\ \hline
   $R_{500,h}$ & 0.949 & 0.743 & 0.844 & 0.781 & 0.771 \\
   $R_{2000,h}$ & 0.928 & 0.932 & 0.954 & 0.929 & 0.990 \\ \hline
      $\alpha_1=-0.5$ & $h=1$ & 2 & 3 & 4 & 5 \\ \hline
    $g_{h}(1)$ & 7.469 & 11.894 & 11.665 & 11.446 & 11.234 \\ \hline
   $R_{500,h}$ & 0.820 & 0.791 & 0.634 & 0.701 & 0.835 \\
   $R_{2000,h}$ & 0.960 & 0.960 & 0.911 & 0.934 & 1.042 \\ \hline
      $\alpha_1=0.8$ & $h=1$ & 2 & 3 & 4 & 5 \\ \hline
    $g_{h}(1)$ & 7.484 & 17.448 & 17.113 & 16.791 & 16.481 \\ \hline
   $R_{500,h}$ & 1.911 & 0.742 & 0.859 & 0.757 & 0.795 \\
   $R_{2000,h}$ & 0.999 & 0.949 & 0.995 & 0.951 & 0.913 \\ \hline
  \end{tabular}
\end{table}

\subsection{Additional conditional heteroscedastic models: EGARCH and HAR}\label{sec:add EGARCH HAR}
We examine the finite-sample performance of the ratio $R_{n,h}$ under two additional conditional heteroscedastic models for $\{\varepsilon_t\}$, the exponential GARCH (EGARCH) model and the heterogeneous autoregressive (HAR) model (\cite{Corsi2009}). In particular, we consider the following examples.

\begin{example}\label{exS.EGARCH HAR AR2}
We generate $M=5000$ realizations from the following AR(2) model,
\begin{equation*}%\label{exS.1.1}
x_t=-0.5 x_{t-2}+\varepsilon_t,
\end{equation*}
where 
$\varepsilon_t$ obeys
\eqref{varepsilon basic form} with $\{z_t\}$ being a sequence of
i.i.d. $N(0,1)$ random variables
and $\{\sigma_t\}$ being either a EGARCH(1,1) process,
\begin{equation*}
(1-0.95B)\log(\sigma_t^2)=0.05-0.05z_{t-1}+0.15(|z_{t-1}|-E(z_{t-1})),
\end{equation*}
or follows a HAR model,
\begin{equation*}
\begin{split}
   \sigma_t= & 0.08+0.36RV^{(d)}_{t-1}+0.28RV_{t-1}^{(w)}+0.28RV_{t-1}^{(m)}+\tilde{\omega}_{t}, \\
   \sigma_t= & RV^{(d)}_{t}+{\omega}_{t},
\end{split}
\end{equation*}
in which $RV^{(d)}_{t}$, $RV_{t}^{(w)}$, and $RV_{t}^{(m)}$ are respectively the daily, weekly, and monthly observed realized volatilities, and $\tilde{\omega}_{t}$ and ${\omega}_{t}$ are sequences of i.i.d. $N(0,0.04)$ random variables. Alternatively, $\{RV^{(d)}_{t}\}$ can be written as an AR(22) model:
\begin{equation*}
RV^{(d)}_{t}=\theta_0+\sum_{i=1}^{22}\theta_{i}RV^{(d)}_{t-i}+u_t,
\end{equation*}
where $\theta_0=0.08$,
\begin{equation*}
\theta_i=
  \begin{cases}
    0.36+\frac{1}{5}0.28+\frac{1}{22}0.28, & \mbox{if } i=1, \\
    \frac{1}{5}0.28+\frac{1}{22}0.28, & \mbox{if } i=2,\ldots,5, \\
    \frac{1}{22}0.28, & \mbox{if } i=6,\ldots,22,
  \end{cases}
\end{equation*}
and $u_t=\tilde{\omega}_{t}-{\omega}_{t}$. To ensure the positivity of $\sigma_t$ and $RV^{(d)}_{t}$, we truncate these values at a lower bound of 0.0001 whenever they fall below this threshold.
\end{example}

\begin{example}\label{exS.EGARCH HAR MA1}
We generate $M=5000$ realizations from
the MA(1) model,
\begin{eqnarray*}
x_t=\varepsilon_t-0.8\varepsilon_{t-1}.
\end{eqnarray*}
The remaining settings are identical to those described in Example \ref{exS.EGARCH HAR AR2}. 
\end{example}

For the EGARCH(1,1) model, the ratio $R_{n,h}$ consistently remains within the $(0.9, 1.1)$ interval across almost all settings, with the only exceptions occurring at $n=500$ and $h=1$ in Tables \ref{tab:exS.EGARCH HAR AR2} and \ref{tab:exS.EGARCH HAR MA1}. For the HAR model, $R_{n,h}$ falls within the range of $(0.88, 1.1)$ as $n$ reaches $2000$.  These findings suggest the asymptotic results established in Theorem \ref{tMSPE} are applicable to both EGARCH and HAR processes.

\begin{table}
\caption{The values of $g_{h}(k)$ and $R_{n,h}$, with $k=1$, $n=500, 2000$, and $h=1,\ldots, 5$, in Example \ref{exS.EGARCH HAR AR2}}
\label{tab:exS.EGARCH HAR AR2}
\centering
\begin{tabular}{ccccccccccc} \hline
    EGARCH(1,1) & $h=1$ & 2 & 3 & 4 & 5 \\ \hline
    $g_{h}(1)$ & 4.834 & 3.582 & 5.916 & 6.151 & 6.966 \\ \hline
   $R_{500,h}$ & 1.196 & 0.981 & 0.955 & 0.975 & 1.025 \\
   $R_{2000,h}$ & 1.090 & 0.981 & 0.990 & 1.006 & 0.932 \\ \hline \hline
   HAR & $h=1$ & 2 & 3 & 4 & 5 \\ \hline
    $g_{h}(1)$ & 1.939 & 1.422 & 2.357 & 2.458 & 2.794 \\ \hline
   $R_{500,h}$ & 1.123 & 0.995 & 0.974 & 0.965 & 1.047 \\
   $R_{2000,h}$ & 0.924 & 1.007 & 0.883 & 0.968 & 0.949 \\ \hline
  \end{tabular}
\end{table}

\begin{table}
\caption{The values of $g_{h}(k)$ and $R_{n,h}$, with $k=1$, $n=500, 2000$, and $h=1,\ldots, 5$, in Example \ref{exS.EGARCH HAR MA1}}
\label{tab:exS.EGARCH HAR MA1}
\centering
\begin{tabular}{ccccccccccc} \hline
    EGARCH(1,1) & $h=1$ & 2 & 3 & 4 & 5 \\ \hline
    $g_{h}(1)$ & 4.404 & 8.734 & 8.632 & 8.537 & 8.448 \\ \hline
   $R_{500,h}$ & 1.197 & 0.967 & 0.997 & 0.999 & 1.027 \\
   $R_{2000,h}$ & 0.943 & 0.953 & 0.912 & 0.970 & 1.009 \\ \hline \hline
   HAR & $h=1$ & 2 & 3 & 4 & 5 \\ \hline
    $g_{h}(1)$ & 1.863 & 3.470 & 3.429 & 3.408 & 3.385 \\ \hline
   $R_{500,h}$ & 1.215 & 1.050 & 1.082 & 1.041 & 0.988 \\
   $R_{2000,h}$ & 1.030 & 1.078 & 1.017 & 1.009 & 1.012 \\ \hline
  \end{tabular}
\end{table}

\subsection{Fat-tailed innovations}\label{sec:heavy tail errors}
This section investigates the performance of $R_{n,h}$ under fat-tailed innovations.

\begin{example}\label{exS.heavy tail}
We generate $M=5000$ realizations from the following AR(2) model,
\begin{equation}\label{exS.heavy tail.1}
x_t=-0.5 x_{t-2}+\varepsilon_t,
\end{equation}
where 
$\varepsilon_t$ obeys
\eqref{varepsilon basic form} with $\{z_t\}$ being a sequence of
i.i.d. standard Student-$t$ random variables with 10 degrees of freedom, 
and $\{\sigma_t\}$ being either a GARCH(1,1) process,
\begin{equation}\label{exS.heavy tail.2}
\sigma_t^2=0.4+0.2\varepsilon_{t-1}^2+0.55\sigma_{t-1}^2,
\end{equation}
or a SV(1) (\eqref{SV model} with $\tilde{p}=1$) process,
\begin{equation}\label{exS.heavy tail.SV}
(1-0.98B)\log(\sigma_t^2)=0.01+v_t,
\end{equation}
in which
$\{v_t\}$ is a sequence of i.i.d. $N(0,0.04)$ random variables.
\end{example}

\begin{example}\label{exS.heavy tail 2}
We generate $M=5000$ realizations from the AR(2) model \eqref{exS.heavy tail.1}, where 
$\varepsilon_t$ obeys
\eqref{varepsilon basic form} with $\{z_t\}$ being a sequence of
i.i.d. standard Student-$t$ random variables with 7 degrees of freedom, 
and $\{\sigma_t\}$ being either a GARCH(1,1) process,
\begin{equation}\label{exS.heavy tail.3}
\sigma_t^2=0.4+0.1\varepsilon_{t-1}^2+0.8\sigma_{t-1}^2,
\end{equation}
or the SV(1) process \eqref{exS.heavy tail.SV}. Note that the GARCH coefficients are adjusted from \eqref{exS.heavy tail.2} to \eqref{exS.heavy tail.3} to satisfy the moment condition $E|\varepsilon_t|^{6+\delta} < \infty$ for some small $\delta > 0$.
\end{example}

Tables \ref{tab:exS.heavy tail} and \ref{tab:exS.heavy tail 2} report the results for Examples \ref{exS.heavy tail} and \ref{exS.heavy tail 2}. When $z_t$ follows a Student-$t$ distribution with 10 degrees of freedom, $R_{n,h}$ still stays near 1 (between $0.895$ and $1.1$) at $n=2000$. However, when the degrees of freedom are reduced to 7, $R_{n,h}$ exhibits greater volatility, with its fluctuations widening to a range of $(0.837, 1.226)$ at $n=2000$. Nevertheless, as $R_{n,h}$ continues to oscillate around 1, these findings indicate that the asymptotic results established in Theorem \ref{tMSPE} remain valid under fat-tailed innovations.

\begin{table}
\caption{The values of $g_{h}(k)$ and $R_{n,h}$, with $k=1$, $n=500, 2000$, and $h=1,\ldots, 5$, in Example \ref{exS.heavy tail}}
\label{tab:exS.heavy tail}
\centering
\begin{tabular}{ccccccccccc} \hline
    GARCH(1,1) ($\textrm{df}=10$) & $h=1$ & 2 & 3 & 4 & 5 \\ \hline
    $g_{h}(1)$ & 4.364 & 2.994 & 4.513 & 4.324 & 4.528 \\ \hline
   $R_{500,h}$ & 0.950 & 0.735 & 0.843 & 0.817 & 0.904 \\
   $R_{2000,h}$ & 0.973 & 0.943 & 0.961 & 0.936 & 1.014 \\ \hline \hline
   SV(1) ($\textrm{df}=10$) & $h=1$ & 2 & 3 & 4 & 5 \\ \hline
    $g_{h}(1)$ & 9.680 & 7.120 & 11.686 & 12.061 & 13.546 \\ \hline
   $R_{500,h}$ & 0.845 & 0.740 & 1.009 & 0.715 & 0.653 \\
   $R_{2000,h}$ & 0.951 & 0.908 & 0.956 & 0.895 & 0.908 \\ \hline
  \end{tabular}
\end{table}

\begin{table}
\caption{The values of $g_{h}(k)$ and $R_{n,h}$, with $k=1$, $n=500, 2000$, and $h=1,\ldots, 5$, in Example \ref{exS.heavy tail 2}}
\label{tab:exS.heavy tail 2}
\centering
\begin{tabular}{ccccccccccc} \hline
    GARCH(1,1) ($\textrm{df}=7$) & $h=1$ & 2 & 3 & 4 & 5 \\ \hline
    $g_{h}(1)$ & 9.356 & 6.718 & 10.839 & 10.983 & 12.120 \\ \hline
   $R_{500,h}$ & 1.082 & 0.839 & 0.870 & 0.790 & 0.938 \\
   $R_{2000,h}$ & 1.026 & 0.872 & 1.226 & 0.871 & 1.121 \\ \hline \hline
   SV(1) ($\textrm{df}=7$) & $h=1$ & 2 & 3 & 4 & 5 \\ \hline
    $g_{h}(1)$ & 9.680 & 7.120 & 11.686 & 12.061 & 13.546 \\ \hline
   $R_{500,h}$ & 0.473 & 0.720 & 0.715 & 0.761 & 0.729 \\
   $R_{2000,h}$ & 0.869 & 0.894 & 1.198 & 0.837 & 1.111 \\ \hline
  \end{tabular}
\end{table}

\subsection{Additional examples of subset selection}\label{sec:add sub sel}
In this section, we revisit the correctly specified case of DGP (I) from Section \ref{sec:4.2} by replacing the GARCH(1,1) errors with alternative conditional heteroscedastic errors, including EGARCH, HAR, and integrated GARCH (IGARCH) processes. Specifically, we consider the following DGP:
\begin{itemize}
  \item [(III)]
  \begin{equation*}%\label{exS5.2}
\begin{split}
     & x_t=0.8x_{t-1}-0.5x_{t-2}+0.35x_{t-4}+\varepsilon_t, \\
     & \varepsilon_t=\sigma_tz_t,  \\
\end{split}
\end{equation*}
with $\{z_t\}$ being a sequence of
i.i.d. $N(0,1)$ random variables
and $\{\sigma_t\}$ being an EGARCH(1,1) process,
\begin{equation*}
(1-0.95B)\log(\sigma_t^2)=0.05-0.05z_{t-1}+0.15(|z_{t-1}|-E(z_{t-1})),
\end{equation*}
an IGARCH(1,1) process,
\begin{equation*}
\sigma_t^2=0.4+0.45\varepsilon_{t-1}^2+0.55\sigma_{t-1}^2,
\end{equation*}
or a HAR process,
\begin{equation*}
\begin{split}
   \sigma_t= & 0.08+0.36RV^{(d)}_{t-1}+0.28RV_{t-1}^{(w)}+0.28RV_{t-1}^{(m)}+\tilde{\omega}_{t}, \\
   \sigma_t= & RV^{(d)}_{t}+{\omega}_{t},
\end{split}
\end{equation*}
in which $RV^{(d)}_{t}$, $RV_{t}^{(w)}$, and $RV_{t}^{(m)}$ are respectively the daily, weekly, and monthly observed realized volatilities, and $\tilde{\omega}_{t}$ and ${\omega}_{t}$ are sequences of i.i.d. $N(0,0.04)$ random variables.
\end{itemize}

Similar to Section \ref{sec:4.2}, we evaluate and compare the performance of MRIC, AIC, and BIC in identifying the true model. Compared to DGP (I) in Section \ref{sec:4.2}, replacing GARCH(1,1) errors with EGARCH(1,1) or HAR errors yields similar results. Regarding the case of IGARCH(1,1) errors, since the assumptions required for Theorem \ref{t MRIC} fail as no finite variance of $\varepsilon_t$ exists, MRIC appears ineffective in this instance. Likewise, BIC and AIC fail to consistently identify the true model in this case. 

\begin{table}
\caption{Frequency, in 1,000 simulations, of selecting the true model when the data is generated from DGP {\rm (III)}.}
\label{tabS5.2}
\centering
\begin{tabular}{cccc}
EGARCH(1,1)  & AIC & BIC & MRIC \\ \hline
   $n= 200$ & 664 & 936 & 919 \\
   500 & 633 & 950 & 996  \\
   1000 & 629 & 974 & 1000 \\  \hline
  HAR& AIC & BIC & MRIC \\ \hline
   $n=200$ & 658 & 921 & 929 \\
   500 & 662 & 953 & 998  \\
   1000 & 637 & 954 & 1000 \\  \hline
   IGARCH(1,1) & AIC & BIC & MRIC \\ \hline
  $n=200$ & 424 & 720 & 554 \\
   500 & 294 & 635 & 592  \\
   1000 & 244 & 558 & 591 \\  \hline
  \end{tabular}
\end{table}

We next examine the performance of subset selection under fat-tailed innovations. Consider the following DGP:

\begin{itemize}
  \item [(IV)]
  \begin{equation*}
\begin{split}
     & x_t=0.8x_{t-1}-0.5x_{t-2}+0.35x_{t-4}+\varepsilon_t, \\
     & \varepsilon_t=\sigma_tz_t,  \\
\end{split}
\end{equation*}
where $\{z_t\}$ is a sequence of  
i.i.d. standard Student-$t$ random variables with $v=7$ or $10$ degrees of freedom, 
and $\{\sigma_t\}$ follows a GARCH(1,1) process,
\begin{equation*}
\sigma_t^2=0.4+0.2\varepsilon_{t-1}^2+0.55\sigma_{t-1}^2.
\end{equation*}
\end{itemize}

DGP (IV) differs from DGP (I) only in the distributional assumption of $\{z_t\}$, which is changed from i.i.d. Normal to i.i.d. standard Student-$t$. As shown in Table \ref{tabS5.2.2}, using Student-$t$ innovations with 10 degrees of freedom results in only subtle differences compared to the Gaussian case. When the degrees of freedom of $z_t$ decrease to 7, the performance of both BIC and MRIC deteriorates slightly; however, the overall behavior of all three criteria remains consistent with earlier observations.

\begin{table}
\caption{Frequency, in 1,000 simulations, of selecting the true model when the data is generated from DGP {\rm (IV)}.}
\label{tabS5.2.2}
\centering
\begin{tabular}{cccc}
GARCH(1,1) ($\textrm{df}=10$)  & AIC & BIC & MRIC \\ \hline
   $n= 200$ & 599 & 898 & 877 \\
   500 & 600 & 931 & 982  \\
   1000 & 568 & 941 & 994 \\  \hline
GARCH(1,1) ($\textrm{df}=7$)& AIC & BIC & MRIC \\ \hline
   $n=200$ & 584 & 873 & 852 \\
   500 & 564 & 904 & 967  \\
   1000 & 561 & 924 & 990 \\  \hline
  \end{tabular}
\end{table}

\section*{Acknowledgements}
We gratefully acknowledge the valuable feedback received on earlier versions of this work, which contributed to improving the manuscript. We also thank the reviewers for their careful reading and constructive comments.

\begin{funding}
This work was supported by the National Science and Technology Council, Taiwan (R.O.C.), under Grant Nos. 114-2118-M-001-003-MY2 (Huang), 112-2122-M-007-001-MY3 (Ing), and 113-2118-M-390-001 (Yu).
\end{funding}



\bibliographystyle{chicago}

\begin{thebibliography}{26}
% BibTex style file: imsart-number.bst, 2017-11-03
% Default style options (sort=1,type=number).
% Used options (sort=1,type=number).

\bibitem[\protect\citeauthoryear{Akaike}{1969}]{Akaike69}
\textsc{Akaike, H.} (1969). Fitting autoregressive models for prediction. \textit{Annals of the Institute of Statistical Mathematics}
\textbf{21} 243--247.

\bibitem[\protect\citeauthoryear{Andrews}{1988}]{Andrews1988M}
\textsc{Andrews, Donald W. K.} (1988). Laws of large numbers for dependent non-identically distributed random variables. \textit{Econometric Theory} \textbf{4} 458--467.


\bibitem[\protect\citeauthoryear{Bhansali}{1981}]{Bhansali1981}
\textsc{Bhansali, R. J.} (1981). Effects of not knowing the order of an autoregressive process on the
  mean squared error of prediction-I. \textit{Journal of the American Statistical Association}
\textbf{76} 588--597.

\bibitem[\protect\citeauthoryear{Bhansali and Papangelou}{1991}]{Bhansali1991}
\textsc{Bhansali, R. J.} and \textsc{Papangelou, F.} (1991). Convergence of moments of least squares estimators for the
  coefficients of an autoregressive process of unknown order. \textit{Annals of Statistics}
\textbf{19} 1155--1162.

\bibitem[\protect\citeauthoryear{Bollerslev}{1986}]{Bollerslev1986}
\textsc{Bollerslev, T.} (1986). Generalized autoregressive conditional heteroskedasticity. \textit{Journal of Econometrics}
\textbf{31} 307--327.

\bibitem[\protect\citeauthoryear{Bollerslev et al.}{1994}]{Bollerslev1994}
\textsc{Bollerslev, T.}, \textsc{Engle, R. F.} and \textsc{Nelson, D. B.} (1994). ARCH models. In \textit{Handbook of Econometrics} (R. F. E. McFadden and D. L., eds.) 2961--3038.

\bibitem[\protect\citeauthoryear{Bradley}{2005}]{Bradley2005}
\textsc{Bradley, R. C.} (2005). Basic properties of strong mixing conditions. A survey and some open questions. \textit{Probability Surveys} \textbf{2} 107--144.

\bibitem[\protect\citeauthoryear{Brockwell and Davis}{1991}]{Brockwell1991}
\textsc{Brockwell, P.} and \textsc{R. Davis} (1991). \textit{Time Series: Theory and Methods} (2nd Ed.). Springer Series in Statistics. Springer.

\bibitem[\protect\citeauthoryear{Carrasco and Chen}{2002}]{Carrasco2002}
\textsc{Carrasco, M.} and \textsc{Chen, X.} (2002). Mixing and moment properties of various GARCH and stochastic volatility models. \textit{Econometric Theory} \textbf{18} 17--39.

\bibitem[\protect\citeauthoryear{Chan and Ing}{2011}]{Chan2011}
\textsc{Chan, N. H.} and \textsc{Ing, C.-K.} (2011). Uniform moment bounds of Fisher's information with applications to
  time series. \textit{Annals of Statistics}
\textbf{39} 1526--1550.

\bibitem[\protect\citeauthoryear{Chevillon}{2007}]{Chevillon2007}
\textsc{Chevillon, G.} (2007). Direct multi-step estimation and forecasting. \textit{Journal of Economic Surveys}
\textbf{21} 746--785.

\bibitem[\protect\citeauthoryear{Chi et al.}{2021}]{Chi2021}
\textsc{Chi, C.-M.}, \textsc{Ing, C.-K.} and \textsc{Yu, S.-H.} (2021). Negative Moment Bounds for Stochastic Regression Models with
  Deterministic Trends and Their Applications to Prediction Problems. \textit{Statistica Sinica}
\textbf{31} 2215--2237.

\bibitem[\protect\citeauthoryear{Chow and Teicher}{1997}]{Chow1997}
\textsc{Chow, Y.} and \textsc{Teicher, H.} (1997). 
\textit{Probability Theory: Independence, Interchangeability and Martingales}. 
Springer, New York.

\bibitem[\protect\citeauthoryear{Corsi}{2009}]{Corsi2009}
\textsc{Corsi, F.} (2009). A simple approximate long-memory model of realized volatility. \textit{Journal of Financial Econometrics} \textbf{7} 174--196.

\bibitem[\protect\citeauthoryear{Davidson}{1994}]{Davidson1994}
\textsc{Davidson, J.} (1994). 
\textit{Stochastic Limit Theory: An Introduction for Econometricians, Advanced Texts in Econometrics}. 
London: Oxford University Press.

\bibitem[\protect\citeauthoryear{Ding et al.}{1993}]{Ding1993}
\textsc{Ding, Z.}, \textsc{Granger, C. W. J.} and \textsc{Engle, R. F.} (1993). A long memory property of stock market returns and a new model. \textit{Journal of Empirical Finance}
\textbf{1} 83--106.

\bibitem[\protect\citeauthoryear{Eldar and Kutyniok}{2012}]{Eldar2012S}
\textsc{Eldar, Y.} and \textsc{Kutyniok, G.} (2012). 
\textit{Compressed Sensing: Theory and Applications}. 
Cambridge University Press.

\bibitem[\protect\citeauthoryear{Findley}{1991}]{Findley91}
\textsc{Findley, D. F.} (1991).
Counterexamples to parsimony and BIC.
\textit{Annals of the Institute of Statistical Mathematics}
\textbf{43} 505--514.

\bibitem[\protect\citeauthoryear{Findley and Wei}{2002}]{Findley2002}
\textsc{Findley, D. F.} and \textsc{Wei, C. Z.} (2002). AIC, overfitting principles, and the boundedness of moments of inverse
  matrices for vector autoregressions and related models. \textit{Journal of Multivariate Analysis}
\textbf{83} 415--450.

\bibitem[\protect\citeauthoryear{Francq and Zakoian}{2019}]{Francq19}
\textsc{Francq, C.} and \textsc{Zakoian, J.-M.} (2019). \textit{GARCH models: structure, statistical inference and financial applications (2nd ed.)}. John Wiley \& Sons Ltd.

\bibitem[\protect\citeauthoryear{Fuller and Hasza}{1981}]{Fuller1981}
\textsc{Fuller, W. A.} and \textsc{Hasza, D. P.} (1981). Properties of predictors for autoregressive time series. \textit{Journal of the American Statistical Association}
\textbf{76} 155--161.

\bibitem[\protect\citeauthoryear{Glosten et al.}{1993}]{Glosten1993}
\textsc{Glosten, L. R.}, \textsc{Jagannathan, R.} and \textsc{Runkle, D. E.} (1993). On the relation between the expected value and the volatility of the
  nominal excess return on stocks. \textit{The Journal of Finance}
\textbf{48} 1779--1801.

\bibitem[\protect\citeauthoryear{Greenway-Mcgrevy}{2013}]{Ryan2013}
\textsc{Greenway-Mcgrevy, R.} (2013). Multistep prediction of panel vector autoregressive processes. \textit{Econometric Theory}
\textbf{29} 699--734.

%\bibitem[\protect\citeauthoryear{Hansen}{2019}]{Hansen2019}
%\textsc{Hansen, B.E.} (2019). A weak law of large numbers under weak mixing. \textit{unpublished manuscript.}

\bibitem[\protect\citeauthoryear{Hsu et al.}{2019}]{Hsu2019}
\textsc{Hsu, H.-L.}, \textsc{Ing, C.-K.} and \textsc{Tong, H.} (2019). On model selection from a finite family of possibly misspecified time
  series models. \textit{Annals of Statistics}
\textbf{47} 1061--1087.

\bibitem[\protect\citeauthoryear{Ing}{2003}]{Ing2003a}
\textsc{Ing, C.-K.} (2003). Multistep prediction in autoregressive processes. \textit{Econometric Theory}
\textbf{19} 254--279.

\bibitem[\protect\citeauthoryear{Ing and Wei}{2003}]{Ing2003}
\textsc{Ing, C.-K.} and \textsc{Wei, C.-Z.} (2003). On same-realization prediction in an infinite-order autoregressive
  process. \textit{Journal of Multivariate Analysis}
\textbf{85} 130--155.

\bibitem[\protect\citeauthoryear{Inoue and Kilian}{2006}]{Kilian2006}
\textsc{Inoue, A.} and \textsc{Kilian, L.} (2006). 
On the selection of forecasting models.
\textit{Journal of Econometrics}
\textbf{130} 273--306.

\bibitem[\protect\citeauthoryear{Jord\`{a}}{2005}]{OJ05}
\textsc{Jord\`{a}, \`{O}.} (2005). Estimation and inference of impulse responses by local projections. \textit{American Economic Review}
\textbf{95} 161--182.

\bibitem[\protect\citeauthoryear{Kim}{1993}]{Kim1993}
\textsc{Kim, T. Y.} (1993). A note on moment bounds for strong mixing sequences. \textit{Statistics} \& \textit{probability Letters} \textbf{16} 163--168.

\bibitem[\protect\citeauthoryear{Kim}{1994}]{Kim1994M}
\textsc{Kim, T. Y.} (1994). Moment bounds for non-stationary dependent sequences \textit{Journal of Applied Probability} \textbf{31} 731--742.

\bibitem[\protect\citeauthoryear{Kunitomo and Yamamoto}{1985}]{Kunitomo1985a}
\textsc{Kunitomo, N.} and \textsc{Yamamoto, T.} (1985). Properties of predictors in misspecified autoregressive time series
  models. \textit{Journal of the American Statistical Association}
\textbf{80} 941--950.

\bibitem[\protect\citeauthoryear{Lewis and Reinsel}{1988}]{Lewis1988}
\textsc{Lewis, R. A.} and \textsc{Reinsel, G. C.} (1988). Prediction error of multivariate time series with mis-specified
  models. \textit{Journal of Time Series Analysis}
\textbf{9} 43--57.

\bibitem[\protect\citeauthoryear{Ling and McAleer}{2002a}]{Ling2002}
\textsc{Ling, S.} and \textsc{McAleer, M.} (2002a). Necessary and sufficient moment conditions for the GARCH($r,s$) and
  asymmetric power GARCH($r,s$) models. \textit{Econometric Theory}
\textbf{18} 722--729.

\bibitem[\protect\citeauthoryear{Ling and McAleer}{2002b}]{Ling2002joe}
\textsc{Ling, S.} and \textsc{McAleer, M.} (2002b). Stationarity andthe existence of moments of a family of GARCH processes. \textit{Journal of Econometrics}
\textbf{106} 109--117.

\bibitem[\protect\citeauthoryear{Mokkadem}{1988}]{Mokkadem1988}
\textsc{Mokkadem, A.} (1988). Mixing properties of ARMA processes. \textit{Stochastic Processes and their Applications} \textbf{29} 309--315.

\bibitem[\protect\citeauthoryear{Nelson}{1990}]{N1990}
\textsc{Nelson, D. B.} (1990). Stationarity and persistence in the GARCH(1,1) model. \textit{Econometric Theory}
\textbf{6} 318--334.

\bibitem[\protect\citeauthoryear{Papangelou}{1994}]{Papangelou1994}
\textsc{Papangelou, F.} (1994). On a distributional bound arising in autoregressive model fitting. \textit{Journal of Applied Probability}
\textbf{31} 401--408.

\bibitem[\protect\citeauthoryear{Schorfheide}{2005}]{Scho2005}
\textsc{Schorfheide, F.} (2005). VAR forecasting under misspecification. \textit{Journal of Econometrics}
\textbf{128} 99--136.

\bibitem[\protect\citeauthoryear{Shaman and Stine}{1988}]{Shaman1988}
\textsc{Shaman, P.} and \textsc{Stine, R. A.} (1988). The bias of autoregressive coefficient estimators. \textit{Journal of the American Statistical Association}
\textbf{83} 842--848.

\bibitem[\protect\citeauthoryear{Tao}{2021}]{Tao21}
\textsc{Tao, T.} (2021). \textit{
An Introduction to Measure Theory (2nd ed.)}. American Mathematical Society.




\bibitem[\protect\citeauthoryear{Taylor}{1982}]{Taylor1982}
\textsc{Taylor, S. J.} (1982). Financial returns modelled by the product of two stochastic processes
  -- a study of daily sugar prices 1961--79. In \textit{Anderson, O. D. (ed.), Time Series Analysis Theory and
  Practice}, \textbf{1} 203--226. North-Holland, Amsterdam.

\bibitem[\protect\citeauthoryear{Taylor}{1986}]{Taylor1986}
\textsc{Taylor, S. J.} (1986). \textit{Modelling Financial Time Series}. John Wiley, Chichester.


\bibitem[\protect\citeauthoryear{Tsay}{2010}]{Tsay2010}
\textsc{Tsay, R. S.} (2010). \textit{Analysis of Financial Time Series, 3rd Edition}. Wiley Interscience, New York.

\bibitem[\protect\citeauthoryear{
West and Zhao}{2016}]{W2016}
\textsc{West, K. D.} and \textsc{Zhao, Z.} (2016). Regressor and disturbance have moments of all order, least squares estimator has none. \textit{Statistics and Probability Letters}
\textbf{115} 54--59.

\bibitem[\protect\citeauthoryear{Wu}{2005}]{Wu2005Main}
\textsc{Wu, W. B.} (2005). Nonlinear system theory: Another look at dependence. \textit{Proceedings of the National Academy of Sciences} \textbf{102} 14150--14154.

\bibitem[\protect\citeauthoryear{Xiao and Wu}{2011}]{Xiao2011}
\textsc{Xiao, H} and \textsc{Wu, W. B.} (2011). Asymptotic inference of autocovariances of stationary processes. \textit{Available at arXiv:1105.3423.}

\bibitem[\protect\citeauthoryear{Zygmund}{2002}]{Zygmund2002}
\textsc{Zygmund, A.} (2002). \textit{Trigonometric Series}. Cambridge University Press, Cambridge.


\end{thebibliography}






\end{document}
